% Description de l'architecture de l'ordinateur — Informatique 3ème — Cours signé OnBuch+
% Programme officiel MINESEC. Compiler avec : tectonic N02-description-architecture-ordinateur.tex
\newcommand{\DOCMATIERE}{Informatique}
\newcommand{\DOCNIVEAU}{3ème}
\newcommand{\DOCMODULE}{Informatique – classe de 3ème}
\newcommand{\DOCLECON}{N02}
\newcommand{\DOCTITRE}{Description de l'architecture de l'ordinateur}
% OnBuch+ — préambule « cours signés » ENRICHI (compilé avec tectonic / XeLaTeX)
% Système visuel « Pop » d'OnBuch+ : fond crème (jamais blanc), bordures encre
% épaisses, ombres « dures » décalées, titres Archivo Black en capitales.
% Version MATHÉMATIQUES : graphes (pgfplots/tikz), boîtes pédagogiques
% (définition, propriété, méthode, attention, le savais-tu, exercice résolu,
% exercice, TP, bilan), page de garde et signature OnBuch+ ancrée en bas.
%
% Macros attendues dans main.tex AVANT \input{preamble} :
%   \DOCMATIERE  \DOCNIVEAU  \DOCMODULE  \DOCLECON (numéro)  \DOCTITRE
\documentclass[11pt]{article}

\usepackage{fontspec}
\usepackage[french]{babel}
\usepackage[a4paper,margin=2cm,top=3.4cm,bottom=2.8cm,headheight=1.4cm,footskip=1.3cm]{geometry}
\usepackage{amsmath,amssymb,amsfonts}
\usepackage{mathtools} % \xmapsto, \coloneqq... souvent utilisés par les modèles
\usepackage{cancel} % simplifications barrées (bilans de forces)
% \abs{z} (module, valeur absolue) et \norm{u} (norme), utilisés par les modèles
\providecommand{\abs}{}\let\abs\relax\DeclarePairedDelimiter{\abs}{\lvert}{\rvert}
\providecommand{\norm}{}\let\norm\relax\DeclarePairedDelimiter{\norm}{\lVert}{\rVert}
\usepackage[table]{xcolor} % [table] : fournit aussi \rowcolors (alternance de lignes)
\usepackage{pagecolor}
\usepackage{tikz}
\usetikzlibrary{arrows.meta,positioning,calc,shapes.geometric,shapes.misc,shapes.arrows,shapes.symbols,decorations.pathmorphing,decorations.pathreplacing,decorations.markings,patterns,shadows,mindmap,trees,fit,backgrounds,matrix,intersections,angles,quotes}
\usepackage{pgfplots}
\usepackage[version=4]{mhchem} % équations nucléaires \ce{^{235}_{92}U}
\usepackage[european,siunitx]{circuitikz} % schémas électriques
\pgfplotsset{compat=1.18}
\usepgfplotslibrary{fillbetween,groupplots} % aires entre courbes (calcul intégral)
\usepackage{enumitem}
\usepackage{fancyhdr}
% [most] charge toutes les bibliothèques tcolorbox (listings, minted, raster,
% documentation...) et gonfle inutilement la mémoire TeX sur un document long
% (~300 boîtes) : on ne charge que ce qui est réellement utilisé.
\usepackage[skins,breakable]{tcolorbox}
\usepackage{graphicx}
\usepackage{colortbl}
\usepackage{booktabs}
\usepackage{tabularx}
\usepackage{diagbox} % case d'en-tête diagonale des tableaux à double entrée
% Colonnes extensibles centrée / gauche / droite (C, L, R), souvent utilisées par les modèles
\newcolumntype{C}{>{\centering\arraybackslash}X}
\newcolumntype{L}{>{\raggedright\arraybackslash}X}
\newcolumntype{R}{>{\raggedleft\arraybackslash}X}
\usepackage{array}
\usepackage{multirow}
\usepackage{multicol}
\usepackage{siunitx}
% parse-numbers=false : accepte aussi \SI{2,5 \times 10^{-3}}{...}
\sisetup{locale=FR,output-decimal-marker={,},group-minimum-digits=4,parse-numbers=false}
\DeclareSIUnit\ohm{\ensuremath{\symup{\Omega}}} % U+2126 absent de la police
\DeclareSIUnit\division{div} % réglages d'oscilloscope (ms/div, V/div)
\DeclareSIUnit\tr{tr} % tours (vitesse de rotation, ex. \SI{1500}{\tr\per\minute})
\DeclareSIUnit\bit{bit}
\DeclareSIUnit\byte{o} % octet
\DeclareSIUnit\inch{po} % pouce
\DeclareSIUnit\ppp{ppp} % points par pouce (résolution d'image)
\usepackage{qrcode}
% Blocs de code source (PHP, C, HTML, JavaScript, SQL) — spécifique à la
% matière Informatique : listings + habillage aux couleurs OnBuch+ (voir le
% style \lstset ci-dessous, à côté des autres boîtes pédagogiques).
\usepackage{listings}
% \degree : commande fréquemment utilisée par les modèles pour un angle ou une
% température en dehors de \SI{}{} (non fournie par les paquets chargés).
\providecommand{\degree}{^{\circ}}
% \qed (fin de démonstration), utilisé par les modèles sans amsthm
\providecommand{\qed}{\hfill$\square$}
% \barre : double barre des valeurs interdites dans un tableau de variations (tkz-tab)
\providecommand{\barre}{\ensuremath{\|}}
% environnement proof (amsthm non chargé) utilisé par les modèles
\ifdefined\proof\else\newenvironment{proof}[1][Démonstration]{\par\noindent\textit{#1.}\ }{\qed\par}\fi
\usepackage{lastpage}
\usepackage{hyperref}
\hypersetup{hidelinks}

% ============================================================
% Palette « Pop » OnBuch+ — mêmes hex que la classe P (plus_ui.dart)
% ============================================================
\definecolor{poporange}{HTML}{F59321}
\definecolor{popdark}{HTML}{E07A0C}
\definecolor{popcream}{HTML}{FFF6E8}
\definecolor{popink}{HTML}{1C1714}
\definecolor{poppurple}{HTML}{7A5AE0}
\definecolor{popgreen}{HTML}{2E9E6A}
\definecolor{poppink}{HTML}{E0457B}
\definecolor{popblue}{HTML}{2F7DE1}
\definecolor{popgold}{HTML}{C98A12}
\definecolor{popmuted}{HTML}{8A7F76}
\definecolor{popline}{HTML}{EADFD2}
\definecolor{popgray}{HTML}{8A7F76}
\colorlet{popgrey}{popgray}
% Alias tolérants (le modèle nomme parfois les couleurs différemment)
\colorlet{popwhite}{white}
\colorlet{popblack}{popink}
\colorlet{popred}{poppink}
\colorlet{popyellow}{popgold}
% Teintes claires pour fonds de tableaux / zones
\colorlet{popblueL}{popblue!10!white}
\colorlet{poporangeL}{poporange!14!white}
\colorlet{popgreenL}{popgreen!12!white}
\colorlet{poppurpleL}{poppurple!12!white}
\colorlet{poppinkL}{poppink!12!white}
\colorlet{popgoldL}{popgold!14!white}

\pagecolor{popcream}

% ============================================================
% Typographie
% ============================================================
\defaultfontfeatures{Ligatures=TeX}
\setmainfont{PlusJakartaSans-Medium.ttf}[Path=../../../fonts/,BoldFont=PlusJakartaSans-Bold.ttf]
% Maths en sans-serif (Fira Math) avec chiffres et lettres droites de Plus
% Jakarta, pour que les formules se fondent dans le texte.
\usepackage[math-style=ISO,bold-style=ISO]{unicode-math}
\setmathfont{FiraMath-Regular.otf}[Path=../../../fonts/]
\setmathfont{PlusJakartaSans-Medium.ttf}[Path=../../../fonts/,range={up/num,up/{Latin,latin},\mathup}]
\usepackage{icomma} % virgule décimale sans espace en mode math
% Caractères mathématiques Unicode absents des polices de texte (Plus Jakarta,
% Archivo Black) : rendus par leur équivalent mathématique, en texte comme en
% formule — sinon ils s'affichent en carré vide (ex. « Arithmétique dans ℤ »).
\usepackage{newunicodechar}
\newunicodechar{ℝ}{\ensuremath{\mathbb{R}}}
\newunicodechar{ℤ}{\ensuremath{\mathbb{Z}}}
\newunicodechar{ℂ}{\ensuremath{\mathbb{C}}}
\newunicodechar{ℕ}{\ensuremath{\mathbb{N}}}
\newunicodechar{ℚ}{\ensuremath{\mathbb{Q}}}
\newunicodechar{𝔻}{\ensuremath{\mathbb{D}}}
\newunicodechar{α}{\ensuremath{\alpha}}
\newunicodechar{β}{\ensuremath{\beta}}
\newunicodechar{γ}{\ensuremath{\gamma}}
\newunicodechar{δ}{\ensuremath{\delta}}
\newunicodechar{ε}{\ensuremath{\varepsilon}}
\newunicodechar{θ}{\ensuremath{\theta}}
\newunicodechar{λ}{\ensuremath{\lambda}}
\newunicodechar{μ}{\ensuremath{\mu}}
\newunicodechar{σ}{\ensuremath{\sigma}}
\newunicodechar{τ}{\ensuremath{\tau}}
\newunicodechar{φ}{\ensuremath{\varphi}}
\newunicodechar{ψ}{\ensuremath{\psi}}
\newunicodechar{ω}{\ensuremath{\omega}}
\newunicodechar{Σ}{\ensuremath{\Sigma}}
% majuscules produites par \MakeUppercase dans les titres (α-aminés -> Α)
\newunicodechar{Α}{\ensuremath{\alpha}}
\newunicodechar{Β}{\ensuremath{\beta}}
\newunicodechar{Γ}{\ensuremath{\Gamma}}
\newunicodechar{Δ}{\ensuremath{\Delta}}
\newunicodechar{Ε}{\ensuremath{\varepsilon}}
\newunicodechar{Θ}{\ensuremath{\Theta}}
\newunicodechar{Λ}{\ensuremath{\Lambda}}
\newunicodechar{Μ}{\ensuremath{\mu}}
\newunicodechar{Τ}{\ensuremath{\tau}}
\newunicodechar{Φ}{\ensuremath{\Phi}}
\newunicodechar{Ψ}{\ensuremath{\Psi}}
\newunicodechar{Ω}{\ensuremath{\Omega}}
\newunicodechar{↦}{\ensuremath{\mapsto}}
\newunicodechar{∈}{\ensuremath{\in}}
\newunicodechar{∉}{\ensuremath{\notin}}
\newunicodechar{∧}{\ensuremath{\wedge}}
\newunicodechar{⇔}{\ensuremath{\Leftrightarrow}}
\newunicodechar{⟺}{\ensuremath{\Longleftrightarrow}}
\newunicodechar{⇒}{\ensuremath{\Rightarrow}}
\newunicodechar{⟹}{\ensuremath{\Longrightarrow}}
\newunicodechar{∘}{\ensuremath{\circ}}
\newunicodechar{∪}{\ensuremath{\cup}}
\newunicodechar{∩}{\ensuremath{\cap}}
\newunicodechar{≡}{\ensuremath{\equiv}}
\newunicodechar{⊂}{\ensuremath{\subset}}
\newunicodechar{⊥}{\ensuremath{\perp}}
\newunicodechar{∃}{\ensuremath{\exists}}
\newunicodechar{∀}{\ensuremath{\forall}}
\newunicodechar{∛}{\ensuremath{\sqrt[3]{\,}}}
\newunicodechar{∜}{\ensuremath{\sqrt[4]{\,}}}
\newunicodechar{⃗}{\ensuremath{{}^{\to}}}
\newunicodechar{ⁿ}{\ifmmode^{n}\else\textsuperscript{\ensuremath{n}}\fi}
\newunicodechar{ˣ}{\ifmmode^{x}\else\textsuperscript{\ensuremath{x}}\fi}
\newunicodechar{ᵇ}{\ifmmode^{b}\else\textsuperscript{\ensuremath{b}}\fi}
\newunicodechar{ʳ}{\ifmmode^{r}\else\textsuperscript{\ensuremath{r}}\fi}
\newunicodechar{ᵘ}{\ifmmode^{u}\else\textsuperscript{\ensuremath{u}}\fi}
\newunicodechar{ʸ}{\ifmmode^{y}\else\textsuperscript{\ensuremath{y}}\fi}
\newunicodechar{ᵃ}{\ifmmode^{a}\else\textsuperscript{\ensuremath{a}}\fi}
\newunicodechar{ᵉ}{\ifmmode^{e}\else\textsuperscript{\ensuremath{e}}\fi}
\newunicodechar{ᵏ}{\ifmmode^{k}\else\textsuperscript{\ensuremath{k}}\fi}
\newunicodechar{ᵖ}{\ifmmode^{p}\else\textsuperscript{\ensuremath{p}}\fi}
\newunicodechar{ᴺ}{\ifmmode^{N}\else\textsuperscript{\ensuremath{N}}\fi}
\newunicodechar{ᶜ}{\ifmmode^{c}\else\textsuperscript{\ensuremath{c}}\fi}
\newunicodechar{ʰ}{\ifmmode^{h}\else\textsuperscript{\ensuremath{h}}\fi}
\newunicodechar{ᵠ}{\ifmmode^{q}\else\textsuperscript{\ensuremath{q}}\fi}
\newunicodechar{ₙ}{\ifmmode_{n}\else\textsubscript{\ensuremath{n}}\fi}
\newunicodechar{ₐ}{\ifmmode_{a}\else\textsubscript{\ensuremath{a}}\fi}
\newunicodechar{ᵢ}{\ifmmode_{i}\else\textsubscript{\ensuremath{i}}\fi}
\newunicodechar{ᵣ}{\ifmmode_{r}\else\textsubscript{\ensuremath{r}}\fi}
\newunicodechar{ₖ}{\ifmmode_{k}\else\textsubscript{\ensuremath{k}}\fi}
\newunicodechar{ᵦ}{\ifmmode_{\beta}\else\textsubscript{\ensuremath{\beta}}\fi}
\newunicodechar{ₓ}{\ifmmode_{x}\else\textsubscript{\ensuremath{x}}\fi}
\newfontfamily\popdisplay{ArchivoBlack-Regular.ttf}[Path=../../../fonts/]
\newfontfamily\poplogo{SpaceGrotesk-Bold.ttf}[Path=../../../fonts/]
\newfontfamily\popsemi{PlusJakartaSans-SemiBold.ttf}[Path=../../../fonts/]
\newfontfamily\popxbold{PlusJakartaSans-ExtraBold.ttf}[Path=../../../fonts/]
\newfontfamily\popxboldls{PlusJakartaSans-ExtraBold.ttf}[Path=../../../fonts/,LetterSpace=3.0]

\setlist[enumerate]{leftmargin=1.7em,itemsep=5pt,topsep=4pt}
\setlist[itemize]{leftmargin=1.7em,itemsep=4pt,topsep=4pt,label={\color{poporange}\textbullet}}
\setlength{\parindent}{0pt}
\setlength{\parskip}{5pt}
\renewcommand{\arraystretch}{1.25}

% pgfplots : style Pop par défaut
\pgfplotsset{
  every axis/.append style={axis line style={popink,line width=1pt},tick style={popink},
    grid style={popline,line width=0.5pt},label style={font=\small},tick label style={font=\footnotesize}},
  popcurve/.style={poporange,line width=1.8pt,no markers,smooth},
}
% Même style disponible dans un \draw TikZ simple (hors pgfplots)
\tikzset{popcurve/.style={poporange,line width=1.8pt}}
\tikzset{popgrid/.style={popline,very thin}}
\colorlet{popcurve}{poporange} % le nom du style est parfois utilisé comme couleur

% ============================================================
% Pill / squiggle
% ============================================================
\newcommand{\poppill}[3]{%
  \tikz[baseline=-0.55ex]\node[draw=popink,line width=1.2pt,rounded corners=2.6pt,%
    fill=#2,inner xsep=6.5pt,inner ysep=3pt]{%
    \popxboldls\fontsize{7}{7}\selectfont\color{#3}\MakeUppercase{#1}};%
}
\newcommand{\popsquiggle}[1]{%
  \tikz[baseline=-0.4ex]{
    \draw[#1,line width=2.6pt,line cap=round]
      plot [smooth] coordinates {(0,0.09) (0.34,0.01) (0.68,0.21) (1.02,0.01) (1.36,0.10)};
  }%
}
% Mot-clé surligné (marqueur orange sous le texte)
\newcommand{\cle}[1]{{\popxbold\color{popdark}#1}}

% ============================================================
% En-tête / pied
% ============================================================
\pagestyle{fancy}
\fancyhf{}
\renewcommand{\headrulewidth}{0pt}
\renewcommand{\footrulewidth}{0pt}
\lhead{\raisebox{-0.15ex}{\poplogo\fontsize{13}{13}\selectfont\color{popink}OnBuch\color{poporange}+}}
\rhead{\poppill{\DOCMATIERE\ \textbullet\ \DOCNIVEAU\ \textbullet\ Leçon \DOCLECON}{white}{popink}}
\lfoot{\popxbold\fontsize{7.6}{7.6}\selectfont\color{popmuted}\MakeUppercase{Cours signé OnBuch+}\ \textbullet\ {\popsemi\DOCTITRE}}
\rfoot{\tikz[baseline=-0.6ex]\node[rounded corners=8.5pt,fill=popink,minimum height=17pt,minimum width=17pt,inner xsep=5pt,inner sep=0]{\popxbold\fontsize{8}{8}\selectfont\color{popcream}\thepage\,/\,\pageref*{LastPage}};}

% ============================================================
% Titres
% ============================================================
\newcommand{\chaptitre}[2]{%
  \begin{tcolorbox}[enhanced,colback=poporange,colframe=popink,boxrule=2.6pt,arc=11pt,
      drop shadow=popink,left=15pt,right=15pt,top=12pt,bottom=11pt,before skip=3pt,after skip=18pt]
    \poppill{#2}{popink}{popcream}\par\vspace{9pt}
    {\raggedright\hyphenpenalty=10000\exhyphenpenalty=10000\popdisplay\fontsize{22}{24}\selectfont\color{popcream}\MakeUppercase{#1}\par}
  \end{tcolorbox}
}
% Section numérotée : pastille orange avec le numéro + titre Archivo
\newcounter{popsec}
\newcommand{\coursec}[1]{%
  \stepcounter{popsec}\setcounter{popsub}{0}%
  \par\vspace{16pt}\noindent
  \begin{minipage}{\linewidth}
  \tikz[baseline=-0.7ex]\node[draw=popink,line width=1.6pt,fill=poporange,rounded corners=4pt,minimum size=22pt,inner sep=0]{\popdisplay\fontsize{12}{12}\selectfont\color{popcream}\thepopsec};%
  \hspace{8pt}{\popdisplay\fontsize{15}{17}\selectfont\color{popink}\MakeUppercase{#1}}\par
  \vspace{4pt}\noindent{\color{poporange}\rule{36pt}{3.6pt}}
  \end{minipage}\par\nopagebreak\vspace{8pt}
}
\newcounter{popsub}
\newcommand{\courssub}[1]{%
  \stepcounter{popsub}%
  \par\vspace{9pt}\noindent{\popxbold\fontsize{11.5}{13}\selectfont\color{popink}{\color{poporange}\thepopsec.\thepopsub}\hspace{6pt}#1}\par\nopagebreak\vspace{4pt}
}

% ============================================================
% Boîtes pédagogiques — même recette : fond blanc, bordure épaisse colorée,
% ombre dure encre, badge plein. Titre optionnel [#1].
% ============================================================
\tcbset{popbase/.style={breakable,enhanced,colback=white,boxrule=2.3pt,arc=9pt,
  shadow={3pt}{-3pt}{0pt}{popink},left=14pt,right=14pt,top=11pt,bottom=11pt,before skip=11pt,after skip=15pt,
  fontupper=\popsemi\fontsize{10.6}{14.5}\selectfont\color{popink},
  fontlower=\popsemi\fontsize{10.6}{14.5}\selectfont\color{popink}}}
% Coche dessinée (le glyphe ✓ n'existe pas dans les polices du document)
\newcommand{\popcheck}[1]{\tikz[baseline=-0.5ex]\draw[#1,line width=1.6pt,line cap=round,line join=round](-0.13,0.0)--(-0.04,-0.09)--(0.13,0.1);}
\providecommand{\coche}{\popcheck{popgreen}}
\providecommand{\resultat}[1]{\par\smallskip\noindent\fcolorbox{popgreen}{popgreenL}{\parbox{\dimexpr\linewidth-2\fboxsep-2\fboxrule\relax}{\textbf{Résultat.} #1}}\par\smallskip}
\newcommand{\popboxhead}[3]{\poppill{#1}{#2}{white}\ifx\relax#3\relax\else\hspace{6pt}{\popxbold\fontsize{10.6}{12}\selectfont\color{popink}#3}\fi\par\vspace{7pt}}

\newtcolorbox{aretenir}[1][]{popbase,colframe=popgreen,before upper={\popboxhead{À retenir}{popgreen}{#1}}}
\newtcolorbox{exemplebox}[1][]{popbase,colframe=poporange,before upper={\popboxhead{Exemple}{poporange}{#1}}}
\newtcolorbox{definition}[1][]{popbase,colframe=popblue,before upper={\popboxhead{Définition}{popblue}{#1}}}
\newtcolorbox{propriete}[1][]{popbase,colframe=poppurple,before upper={\popboxhead{Propriété}{poppurple}{#1}}}
\newtcolorbox{methode}[1][]{popbase,colframe=popgold,before upper={\popboxhead{Méthode}{popgold}{#1}}}
\newtcolorbox{attention}[1][]{popbase,colframe=poppink,before upper={\popboxhead{Attention}{poppink}{#1}}}
\newtcolorbox{experience}[1][]{popbase,colframe=popblue,colback=popblueL,before upper={\popboxhead{Expérience}{popblue}{#1}}}
% « Le savais-tu ? » : carte violette pleine (contexte, culture, Cameroun)
\newtcolorbox{savaistu}[1][]{popbase,colback=poppurple,colframe=popink,
  fontupper=\popsemi\fontsize{10.4}{14.2}\selectfont\color{white},
  before upper={\poppill{Le savais-tu\,?}{white}{poppurple}\ifx\relax#1\relax\else\hspace{6pt}{\popxbold\color{white}#1}\fi\par\vspace{7pt}}}
% Exercice résolu : énoncé en haut, \tcblower puis la solution sur fond crème
\newcounter{popexo}\newcounter{popexr}
\newtcolorbox{exoresolu}[1][]{popbase,colframe=poporange,colbacklower=popcream,
  segmentation style={popink,line width=1.2pt,dashed},
  before upper={\stepcounter{popexr}\popboxhead{Exercice résolu \thepopexr}{poporange}{#1}},
  before lower={\poppill{Solution}{popink}{popcream}\par\vspace{6pt}}}
% Exercice (non résolu en place) — la correction est dans la partie Corrigés
\newtcolorbox{exercice}[1][]{popbase,colframe=popink,
  before upper={\stepcounter{popexo}\popboxhead{Exercice \thepopexo}{popink}{#1}}}
% Corrigé
\newtcolorbox{corrige}[1][]{popbase,colframe=popgreen,colback=popgreenL,
  before upper={\popboxhead{Corrigé}{popgreen}{#1}}}
% Objectifs / prérequis / bilan
\newtcolorbox{objectifs}{popbase,colframe=popink,colback=poporangeL,
  before upper={\popboxhead{Objectifs de la leçon}{popink}{}}}
\newtcolorbox{prerequis}{popbase,colframe=popink,colback=popblueL,
  before upper={\popboxhead{Prérequis}{popblue}{}}}
\newtcolorbox{bilan}{popbase,colframe=popink,colback=white,boxrule=2.8pt,
  before upper={\popboxhead{Fiche bilan}{poporange}{L'essentiel à connaître pour l'examen}}}
% Figure encadrée avec légende
\newtcolorbox{popfigure}[1][]{enhanced,breakable=false,colback=white,colframe=popline,boxrule=1.4pt,arc=8pt,
  left=10pt,right=10pt,top=10pt,bottom=8pt,before skip=10pt,after skip=12pt,halign=center,
  lower separated=false,halign lower=center,
  fontlower=\popsemi\fontsize{9}{11.5}\selectfont\color{popmuted},
  #1}
\newcommand{\legende}[1]{\par\vspace{6pt}{\popsemi\fontsize{9}{11.5}\selectfont\color{popmuted}#1}}

% ============================================================
% Bloc de code source — \begin{lstlisting}[language=Php]...\end{lstlisting}
% (ou language=C, HTML, JavaScript, SQL ; option omise = pseudo-code brut).
% Même recette visuelle que les figures (\popfigure) : cadre encre épais,
% fond clair (popcream), légende possible juste après via \legende{...}
% comme pour une figure (listings ne sait pas arrondir ses coins : on garde
% un cadre rectangulaire, cohérent avec les autres tableaux du document).
% CONTRAT : les modèles ne doivent utiliser QUE \begin{lstlisting}[...] tel
% quel (pas de \begin{tcblisting}, pas de \verb pour un bloc multi-lignes).
% ============================================================
\lstdefinestyle{poplst}{
  backgroundcolor=\color{popcream},
  basicstyle=\popsemi\fontsize{8.6}{11}\selectfont\color{popink},
  keywordstyle=\bfseries\color{popblue},
  commentstyle=\itshape\color{popmuted},
  stringstyle=\color{popgreen},
  numberstyle=\tiny\color{popmuted},
  identifierstyle=\color{popink},
  frame=l,
  rulecolor=\color{popink},
  framerule=1.6pt,
  framesep=7pt,
  framexleftmargin=7pt,
  framexrightmargin=7pt,
  xleftmargin=2pt,
  xrightmargin=2pt,
  breaklines=true,
  breakatwhitespace=false,
  showstringspaces=false,
  showspaces=false,
  showtabs=false,
  tabsize=2,
  columns=flexible,
  keepspaces=true,
  numbers=none,
  aboveskip=10pt,
  belowskip=10pt,
  captionpos=b,
}
\lstset{style=poplst, language=PHP}
% La distribution TeX embarquée ne fournit pas le fichier de langue
% JavaScript de listings (lstlang*.dat) : on la redéfinit nous-mêmes
% pour que \begin{lstlisting}[language=JavaScript] fonctionne.
\lstdefinelanguage{JavaScript}{
  keywords={break,case,catch,class,const,continue,debugger,default,delete,do,
    else,export,extends,finally,for,function,if,import,in,instanceof,let,new,
    return,static,super,switch,this,throw,try,typeof,var,void,while,with,yield,
    async,await,of,get,set},
  keywordstyle=\bfseries\color{popblue},
  ndkeywords={true,false,null,undefined,NaN,Infinity},
  ndkeywordstyle=\bfseries\color{poporange},
  identifierstyle=\color{popink},
  sensitive=true,
  comment=[l]{//},
  morecomment=[s]{/*}{*/},
  string=[b]",
  morestring=[b]',
  morestring=[b]`,
}
% Même limitation pour CSS et les fichiers de configuration (apacheconf,
% nginx, ini...) : ils ne sont pas fournis. CSS et un style générique de
% type "config" (commentaires # ou ;, pas de mots-clés) couvrent les besoins.
\lstdefinelanguage{CSS}{
  keywords={color,background,background-color,margin,padding,border,
    border-radius,font,font-size,font-weight,font-family,width,height,
    display,position,top,left,right,bottom,float,clear,text-align,
    line-height,list-style,cursor,overflow,box-shadow,transition,
    flex,grid,align-items,justify-content,z-index},
  keywordstyle=\bfseries\color{popblue},
  identifierstyle=\color{popink},
  sensitive=false,
  comment=[s]{/*}{*/},
  string=[b]",
  morestring=[b]',
}
\lstdefinelanguage{apacheconf}{
  sensitive=false,
  morecomment=[l]{\#},
}
\lstdefinelanguage{Apache}{
  sensitive=false,
  morecomment=[l]{\#},
}
\lstdefinelanguage{text}{sensitive=false}
\lstdefinelanguage{powershell}{
  keywords={New-Object,New-SmbShare,Get-Acl,Set-Acl,Get-ChildItem,Set-Content,
    Get-Content,Write-Host,ForEach-Object,Where-Object,Select-Object,
    Test-Path,Remove-Item,Copy-Item,Move-Item,Start-Process,Stop-Process,
    If,Else,ElseIf,ForEach,While,Function,Return,Try,Catch},
  keywordstyle=\bfseries\color{popblue},
  identifierstyle=\color{popink},
  sensitive=false,
  comment=[l]{\#},
  string=[b]",
  morestring=[b]',
}
\lstdefinelanguage{ini}{
  sensitive=false,
  morecomment=[l]{;},
  morecomment=[l]{\#},
}

% ============================================================
% Signature OnBuch+
% ============================================================
\newcommand{\poplogoinline}[1]{{\poplogo\fontsize{#1}{#1}\selectfont\color{popink}OnBuch\color{poporange}+}}
\newcommand{\signatureonbuch}{%
  \begin{tcolorbox}[enhanced,colback=popink,colframe=popink,boxrule=2.2pt,arc=10pt,
      drop shadow=poporange,left=14pt,right=14pt,top=10pt,bottom=10pt,before skip=0pt,after skip=0pt]
    \tikz[baseline=-0.7ex]\node[circle,fill=popgreen,draw=popcream,line width=1.3pt,minimum size=22pt,inner sep=0]{\popcheck{white}};%
    \hspace{9pt}\begin{minipage}[c]{0.62\linewidth}
      {\popxbold\fontsize{12}{13}\selectfont\color{popcream}Signé\ \textbullet\ {\poplogo OnBuch\color{poporange}+}}\par\vspace{2pt}
      {\raggedright\popsemi\fontsize{8}{10.5}\selectfont\color{popline}Ludovic A.\\\MakeUppercase{Conforme au programme officiel MINESEC}\par}
    \end{minipage}\hfill
    {\poplogo\fontsize{22}{22}\selectfont\color{popcream}OnBuch\color{poporange}+}
  \end{tcolorbox}%
}

% ============================================================
% Page de garde
% #1 titre · #2 matière · #3 niveau · #4 module · #5 n° leçon · #6 durée ·
% #7 description / objectifs (texte ou itemize)
% ============================================================
\newcommand{\pagegarde}[7]{%
  \thispagestyle{empty}
  \newgeometry{a4paper,margin=1.7cm,top=1.6cm,bottom=1.4cm}%
  \begin{tikzpicture}[remember picture,overlay]
    \fill[poporange!18] ([xshift=-3.2cm,yshift=-3.6cm]current page.north east) circle (4.6cm);
    \fill[poppurple!14] ([xshift=2.4cm,yshift=7.6cm]current page.south west) circle (3.8cm);
  \end{tikzpicture}%
  {\poplogo\fontsize{30}{30}\selectfont\color{popink}OnBuch\color{poporange}+}\hfill
  \raisebox{0.6ex}{\poppill{Cours signé \textbullet\ Premium}{popink}{popcream}}\par
  \vspace{1pt}\popsquiggle{poppurple}\par
  \vspace{8pt}
  \begin{tcolorbox}[enhanced,colback=poporange,colframe=popink,boxrule=2.8pt,arc=13pt,
      drop shadow=popink,left=18pt,right=18pt,top=12pt,bottom=13pt,after skip=10pt]
    \poppill{#2 \textbullet\ #3}{popink}{popcream}\hspace{5pt}\poppill{#4}{white}{popink}\par\vspace{8pt}
    {\popdisplay\fontsize{14}{14}\selectfont\color{popink}LEÇON #5}\par\vspace{4pt}
    {\raggedright\hyphenpenalty=10000\exhyphenpenalty=10000\popdisplay\fontsize{25}{27}\selectfont\color{popcream}\MakeUppercase{#1}\par}
    \vspace{8pt}
    {\popxbold\fontsize{9.5}{9.5}\selectfont\color{popink}\MakeUppercase{Durée conseillée : #6}}
  \end{tcolorbox}
  \begin{tcolorbox}[enhanced,colback=white,colframe=popink,boxrule=2.2pt,arc=10pt,
      drop shadow=popink,left=16pt,right=16pt,top=10pt,bottom=10pt,after skip=8pt]
    \poppill{Dans cette leçon}{poppurple}{white}\par\vspace{5pt}
    \popsemi\fontsize{9.3}{12.6}\selectfont\color{popink}#7
  \end{tcolorbox}
  \noindent\begin{minipage}[t]{0.487\linewidth}
    \begin{tcolorbox}[enhanced,colback=popgreen,colframe=popink,boxrule=2.2pt,arc=10pt,
        drop shadow=popink,left=11pt,right=11pt,top=10pt,bottom=10pt,height=3.1cm,valign=center]
      \tcbox[colback=white,colframe=popink,boxrule=1.3pt,arc=5pt,boxsep=0pt,nobeforeafter,
        left=3pt,right=3pt,top=3pt,bottom=3pt,valign=center]{\qrcode[height=1.9cm]{https://play.google.com/store/apps/details?id=com.luvvix.onbuch&hl=fr}}%
      \hspace{8pt}\begin{minipage}[c]{0.5\linewidth}
        {\popdisplay\fontsize{11}{12.5}\selectfont\color{white}\MakeUppercase{Télécharge OnBuch}}\par\vspace{4pt}
        {\raggedright\popsemi\fontsize{8.4}{11}\selectfont\color{white}Gratuit \textbullet\ Google Play\\Tuteur IA, annales, concours, résultats\par}
      \end{minipage}
    \end{tcolorbox}
  \end{minipage}\hfill
  \begin{minipage}[t]{0.487\linewidth}
    \begin{tcolorbox}[enhanced,colback=poppurple,colframe=popink,boxrule=2.2pt,arc=10pt,
        drop shadow=popink,left=11pt,right=11pt,top=10pt,bottom=10pt,height=3.1cm,valign=center]
      \tikz[baseline=-0.7ex]\node[circle,fill=white,draw=popink,line width=1.3pt,minimum size=1.6cm,inner sep=0]{\popdisplay\fontsize{24}{24}\selectfont\color{poppurple}+};%
      \hspace{8pt}\begin{minipage}[c]{0.6\linewidth}
        {\popdisplay\fontsize{11}{12.5}\selectfont\color{white}\MakeUppercase{Passe à OnBuch+}}\par\vspace{4pt}
        {\raggedright\popsemi\fontsize{8.4}{11}\selectfont\color{white}Tous les cours signés, Tuteur IA illimité, fiches PDF — onglet OnBuch+ de l'app\par}
      \end{minipage}
    \end{tcolorbox}
  \end{minipage}
  \vspace{8pt}
  \popcontenu
  \vfill
  \signatureonbuch
  \restoregeometry
  \newpage
}

% Rangée « Ce cours contient » (page de garde)
\newcommand{\poptile}[3]{% #1 couleur #2 gros texte #3 légende
  \begin{tcolorbox}[enhanced,colback=#1,colframe=popink,boxrule=2pt,arc=9pt,shadow={2.5pt}{-2.5pt}{0pt}{popink},
      left=6pt,right=6pt,top=9pt,bottom=9pt,halign=center,valign=center,height=1.75cm,width=\linewidth,nobeforeafter]
    {\popdisplay\fontsize{12.5}{13}\selectfont\color{white}\MakeUppercase{#2}}\par\vspace{3pt}
    {\centering\popsemi\fontsize{7.8}{9.5}\selectfont\color{white}#3\par}
  \end{tcolorbox}}
\newcommand{\popcontenu}{%
  \par\noindent{\popxboldls\fontsize{7.5}{7.5}\selectfont\color{popmuted}CE COURS SIGNÉ CONTIENT}\par\vspace{7pt}
  \noindent\begin{minipage}[t]{0.235\linewidth}\poptile{popblue}{Cours}{complet, illustré, pas à pas}\end{minipage}\hfill
  \begin{minipage}[t]{0.235\linewidth}\poptile{poporange}{Méthodes}{et exercices résolus}\end{minipage}\hfill
  \begin{minipage}[t]{0.235\linewidth}\poptile{poppink}{Exercices}{type examen + corrigés}\end{minipage}\hfill
  \begin{minipage}[t]{0.235\linewidth}\poptile{popgreen}{Fiche bilan}{l'essentiel à retenir}\end{minipage}\par}

% Sommaire
\newtcolorbox{sommaire}{enhanced,colback=white,colframe=popink,boxrule=2.2pt,arc=10pt,
  drop shadow=popink,left=18pt,right=18pt,top=13pt,bottom=13pt,before skip=6pt,after skip=20pt,
  before upper={\poppill{Sommaire}{poppurple}{white}\par\vspace{9pt}\popsemi\fontsize{10.6}{14.5}\selectfont\color{popink}}}

% Signature de fin de cours — poussée tout en bas de la dernière page
\newcommand{\findecours}{%
  \par\vfill\nopagebreak
  \begin{center}\popsquiggle{poporange}\end{center}\vspace{4pt}
  \signatureonbuch
}
\AtBeginDocument{\def\llbracket{\mathopen{[\mkern-3.5mu[}}\def\rrbracket{\mathclose{]\mkern-3.5mu]}}} % intervalles d'entiers (glyphe ⟦ absent de la police)
% Glyphes absents de Fira Math : redéfinis à partir de glyphes disponibles
\AtBeginDocument{%
  \def\setminus{\mathbin{\backslash}}\let\smallsetminus\setminus
  \def\vdots{\mathord{\vbox{\baselineskip3.2pt\lineskiplimit0pt\kern4pt\hbox{.}\hbox{.}\hbox{.}}}}%
  \def\ddots{\mathinner{\mkern1mu\raise6.5pt\hbox{.}\mkern2mu\raise3.6pt\hbox{.}\mkern2mu\raise0.7pt\hbox{.}\mkern1mu}}%
  \def\triangle{\mathord{\scalebox{0.9}{$\symup{\Delta}$}}}%
  \def\nabla{\mathord{\raisebox{\depth}{\scalebox{1}[-1]{$\symup{\Delta}$}}}}%
  \def\checkmark{\popcheck{popdark}}%
  \def\star{\mathbin{\ast}}\def\bigstar{\mathord{\ast}}%
  \def\diamond{\mathbin{\scalebox{0.6}{$\Diamond$}}}%
  \def\bigcup{\mathop{\mathchoice{\raisebox{-0.3em}{\scalebox{1.7}{$\cup$}}}{\scalebox{1.25}{$\cup$}}{\cup}{\cup}}\displaylimits}%
  \def\bigcap{\mathop{\mathchoice{\raisebox{-0.3em}{\scalebox{1.7}{$\cap$}}}{\scalebox{1.25}{$\cap$}}{\cap}{\cap}}\displaylimits}%
  \def\lesseqqgtr{\mathrel{\lessgtr}}\def\bigoplus{\mathop{\oplus}\displaylimits}%
}
\providecommand{\ssi}{si et seulement si} % abréviation fréquente
\AtBeginDocument{\providecommand{\wideparen}{\overparen}} % arc de cercle

%% FIGFIX-BEGIN v4 — correctifs globaux des figures (docs/onbuch-generation/tools/figfix)
\usepackage{adjustbox}\usepackage{etoolbox}
% toute figure TikZ/pgfplots plus large que la ligne est réduite à la largeur disponible
\BeforeBeginEnvironment{tikzpicture}{\begin{adjustbox}{max width=\linewidth}}
\AfterEndEnvironment{tikzpicture}{\end{adjustbox}}
% légendes pgfplots lisibles par-dessus les courbes
\pgfplotsset{every axis/.append style={legend style={fill=white,fill opacity=0.92,text opacity=1,draw=black!15,inner sep=3pt}}}
% étiquettes d'arêtes (syntaxe edge["..."]) avec fond blanc
\tikzset{every edge quotes/.append style={fill=white,inner sep=1.2pt}}
%% FIGFIX-END
\begin{document}

\pagegarde{Description de l'architecture de l'ordinateur}{Informatique}{3ème}{Informatique – classe de 3ème}{N02}{2 séances (fiche de progression harmonisée)}{Cette leçon présente l'architecture générale d'un ordinateur (unité centrale, mémoires, périphériques, bus) et introduit les deux grandes familles d'architectures de processeurs : RISC et CISC. L'élève apprend à décrire, comparer et justifier le rôle de chaque composant dans des situations concrètes de la vie courante.
\vspace{4pt}
\begin{multicols}{2}\small
\begin{itemize}
\item À la fin de cette leçon, je sais définir l'architecture d'un ordinateur et citer ses grandes unités fonctionnelles
\item À la fin de cette leçon, je sais décrire le rôle de l'unité centrale de traitement (CPU), de l'unité arithmétique et logique et de l'unité de commande
\item À la fin de cette leçon, je sais distinguer les différents types de mémoires (RAM, ROM, mémoire de masse) et leur usage
\item À la fin de cette leçon, je sais expliquer le rôle des bus et des périphériques d'entrée/sortie
\item À la fin de cette leçon, je sais décrire le cycle d'exécution d'une instruction (fetch-decode-execute)
\item À la fin de cette leçon, je sais définir les architectures RISC et CISC et citer leurs caractéristiques
\end{itemize}\end{multicols}}

\begin{sommaire}
\begin{multicols}{2}
\begin{enumerate}
\item Qu'est-ce que l'architecture d'un ordinateur ?
\item L'unité centrale de traitement (CPU)
\item Les mémoires et les bus
\item Le cycle d'exécution d'une instruction
\item Les architectures RISC et CISC
\item Appliquer et justifier : démarche argumentée
\item Activité d'intégration
\item Méthodes et exercices résolus
\item Exercices
\item Corrigés des exercices
\item Fiche bilan
\end{enumerate}
\end{multicols}
\end{sommaire}

\begin{objectifs}
\begin{itemize}
\item À la fin de cette leçon, je sais définir l'architecture d'un ordinateur et citer ses grandes unités fonctionnelles
\item À la fin de cette leçon, je sais décrire le rôle de l'unité centrale de traitement (CPU), de l'unité arithmétique et logique et de l'unité de commande
\item À la fin de cette leçon, je sais distinguer les différents types de mémoires (RAM, ROM, mémoire de masse) et leur usage
\item À la fin de cette leçon, je sais expliquer le rôle des bus et des périphériques d'entrée/sortie
\item À la fin de cette leçon, je sais décrire le cycle d'exécution d'une instruction (fetch-decode-execute)
\item À la fin de cette leçon, je sais définir les architectures RISC et CISC et citer leurs caractéristiques
\item À la fin de cette leçon, je sais comparer RISC et CISC et identifier dans quels appareils chacune est utilisée
\item À la fin de cette leçon, je sais rédiger et justifier une réponse argumentée sur le choix d'un équipement informatique adapté à un besoin concret
\end{itemize}
\end{objectifs}

\begin{prerequis}
\begin{itemize}
\item Notion de système informatique : matériel (hardware) et logiciel (software) vus en 4ème
\item Identification des périphériques courants (clavier, souris, écran, imprimante)
\item Unités de mesure de l'information : bit, octet, Ko, Mo, Go
\item Notion de programme et d'instruction (début d'année de 3ème)
\item Notion de fréquence d'horloge (Hertz, MHz, GHz) — rappel de physique
\end{itemize}
\end{prerequis}

\newpage

\coursec{Qu'est-ce que l'architecture d'un ordinateur ?}

\courssub{Une situation pour entrer dans la leçon}

Imagine la scène : nous sommes dans un cybercafé animé du quartier Mvog-Ada à Yaoundé. Deux ordinateurs de bureau sont côte à côte, même marque apparente, même prix affiché sur l'étiquette. Tu t'assois devant le premier pour taper ton mémoire de fin d'année sous \texttt{LibreOffice Writer} : le curseur clignote instantanément, l'enregistrement est fluide, tu peux même écouter de la musique en fond. Tu changes de poste pour le second : dès que tu ouvres le navigateur avec trois onglets (le portail \texttt{www.minesec.gov.cm}, une vidéo YouTube et ton webmail), la machine \emph{rame}, le ventilateur hurle, la souris saccade. Pourtant, de l'extérieur, les deux boîtiers noirs se ressemblent comme deux gouttes d'eau !

Autre mystère : le smartphone Android de ton grand frère (un Tecno ou un Infinix tout neuf) tourne avec une puce \texttt{ARM} (architecture RISC), tandis que l'ordinateur du secrétariat du collège utilise un processeur \texttt{Intel} ou \texttt{AMD} (architecture CISC/x86). Pourquoi cette différence ? Qu'est-ce qui se cache \emph{sous le capot} et qui explique la vitesse, la consommation d'énergie, le prix, les possibilités d'évolution ?

\begin{aretenir}[Problématique de la leçon]
L'\textbf{architecture d'un ordinateur} définit l'organisation interne, l'interconnexion et le fonctionnement coordonné de ses composants matériels. C'est le « plan de la maison » qui détermine si l'ordinateur sera un sprinter (téléphone), un marathonien (serveur) ou un véhicule tout-terrain (PC bureautique). Comprendre cette architecture, c'est pouvoir \textbf{choisir}, \textbf{dépanner} et \textbf{programmer} efficacement.
\end{aretenir}

\courssub{De l'analogie de la maison à la définition rigoureuse}

\begin{definition}[Architecture d'un ordinateur]
L'\cle{architecture d'un ordinateur} est l'ensemble des règles, des structures et des méthodes qui décrivent l'\textbf{organisation}, l'\textbf{interconnexion} et le \textbf{fonctionnement} des composants matériels (hardware) permettant le traitement automatique de l'information. Elle précise \emph{quoi} fait chaque bloc, \emph{comment} ils communiquent et \emph{dans quel ordre} les opérations s'enchaînent.
\end{definition}

Pour ancrer l'intuition, faisons l'analogie avec une \textbf{maison} :
\begin{itemize}
    \item Le \textbf{plan d'architecte} (murs, pièces, couloirs, gaines électriques) $\rightarrow$ l'\textbf{architecture matérielle} (schéma des bus, emplacement CPU, slots RAM, ports USB).
    \item Les \textbf{pièces spécialisées} (cuisine, chambre, bureau, salon) $\rightarrow$ les \textbf{unités fonctionnelles} (CPU = bureau de travail, RAM = plan de travail immédiat, Disque dur = armoire d'archives, Clavier/Souris = porte d'entrée/sortie).
    \item Les \textbf{couloirs et portes} qui permettent de circuler entre les pièces $\rightarrow$ les \textbf{bus} (bus de données, d'adresses, de contrôle) qui transportent l'information.
    \item Les \textbf{occupants} qui suivent des recettes de cuisine ou des dossiers de travail $\rightarrow$ le \textbf{logiciel} (système d'exploitation, applications) qui \emph{pilote} le matériel.
\end{itemize}

\begin{attention}[Piège classique : Apparence vs Architecture]
Ne confonds jamais le \textbf{boîtier} (la tour noire, le châssis du portable, la coque du téléphone) avec l'\textbf{architecture}. Deux ordinateurs aux boîtiers identiques peuvent avoir des architectures radicalement différentes (processeur 32 bits vs 64 bits, bus PCIe 3.0 vs 4.0, RAM DDR4 vs DDR5). L'architecture est \emph{invisible} de l'extérieur ; elle se lit sur la fiche technique (datasheet) et se comprend par le schéma logique.
\end{attention}

\courssub{Le modèle de Von Neumann : la référence fondatrice}

Avant 1945, les machines comme l'ENIAC étaient programmées en rebranchant physiquement des câbles et en actionnant des commutateurs : changer de programme prenait des jours. En juin 1945, dans un rapport resté célèbre (\emph{First Draft of a Report on the EDVAC}), le mathématicien \textbf{John von Neumann} (né János Neumann à Budapest, naturalisé américain) propose une organisation révolutionnaire qui équipe encore 99\% de nos machines actuelles (PC, serveurs, smartphones, objets connectés).

\begin{definition}[Modèle de Von Neumann (Machine à programme enregistré)]
Le \cle{modèle de Von Neumann} définit une architecture d'ordinateur fondée sur trois principes majeurs :
\begin{enumerate}
    \item \textbf{Mémoire unifiée unique} : les \emph{instructions} (le programme) et les \emph{données} (nombres, textes, images) sont stockées dans la \textbf{même mémoire centrale}, adressée linéairement (chaque case a une adresse numérique).
    \item \textbf{Programme enregistré} : le programme est chargé en mémoire avant exécution, comme n'importe quelle donnée. Pour changer de tâche, il suffit de charger un autre programme en mémoire (gain de temps immense).
    \item \textbf{Exécution séquentielle} : une \textbf{unité de contrôle} lit les instructions une par une dans l'ordre des adresses (sauf saut explicite), les décode et commande leur exécution par l'unité arithmétique et logique.
\end{enumerate}
\end{definition}

\begin{savaistu}[John von Neumann (1903--1957) : l'architecte du monde numérique]
Mathématicien d'origine hongroise, figure centrale de l'Institut d'études avancées de Princeton, von Neumann a posé les bases de l'informatique moderne \emph{avant même que les transistors n'existent} ! Son modèle sépare la \textbf{structure} (hardware) du \textbf{contenu} (software), rendant l'ordinateur \textbf{universel} : la même machine peut traiter la paie des agents de la CAMTEL, simuler la météo sur Douala, ou faire tourner un jeu vidéo, pourvu qu'on lui fournisse le bon programme en mémoire. C'est cette flexibilité qui a permis l'explosion du numérique au Cameroun comme partout ailleurs.
\end{savaistu}

\begin{propriete}[Conséquences directes du modèle Von Neumann]
\begin{itemize}
    \item \textbf{Flexibilité totale} : une machine unique pour tous les usages (bureautique, CAO, serveur web, IA).
    \item \textbf{Goulot d'étranglement de Von Neumann} : le bus unique partagé entre instructions et données limite le débit (le CPU attend souvent la mémoire). Les architectures modernes (caches, bus séparés Harvard modifié, multicœurs) cherchent à l'atténuer.
    \item \textbf{Sécurité} : comme code et données cohabitent, un programme malveillant peut modifier son propre code ou celui d'un autre (débordement de tampon, injection de code) $\rightarrow$ nécessité de protections matérielles (bit NX/DEP, ASLR) et logicielles.
\end{itemize}
\end{propriete}

\courssub{Les quatre grandes unités fonctionnelles}

Le modèle de Von Neumann se concrétise en \textbf{quatre blocs matériels} interconnectés par des bus. C'est la « vue en blocs » que tout technicien doit savoir dessiner et expliquer.

\begin{center}
\begin{tabularx}{\linewidth}{|l|X|X|}
\hline
\rowcolor{popblueL}
\textbf{Unité fonctionnelle} & \textbf{Rôle principal (Analogie)} & \textbf{Composants clés (Exemples actuels)} \\
\hline
\textbf{Unité Centrale de Traitement (CPU / Processeur)} & \emph{Le « cerveau » / Le chef de chantier} : lit, décode, exécute les instructions, pilote les autres unités. & Cœur(s) d'exécution (ALU, FPU), Unité de contrôle, Registres, Caches L1/L2/L3. Ex : Intel Core i5-12400, AMD Ryzen 5 5600G, ARM Cortex-A78 (mobile). \\
\hline
\textbf{Mémoire Centrale (RAM / Main Memory)} & \emph{Le « plan de travail » / L'établi} : stockage \textbf{volatile} (perdu à l'arrêt) et \textbf{accès rapide} des programmes et données en cours d'exécution. & Barrettes DDR4/DDR5 (DIMM pour tour, SO-DIMM pour portable). Capacité typique 2025 : \SI{8}{Go} (entrée de gamme) à \SI{32}{Go}+ (stations de travail). \\
\hline
\textbf{Unités d'Entrée / Sortie (E/S / Périphériques)} & \emph{Les « portes et fenêtres »} : communication avec le monde extérieur (humain, réseau, autres machines). & \textbf{Entrée} : clavier, souris, scanner, webcam, microphone, capteurs. \textbf{Sortie} : écran, imprimante, haut-parleur. \textbf{Entrée/Sortie} : disque SSD/HDD, clé USB, carte réseau (RJ45, Wi-Fi), modem 4G/5G. \\
\hline
\textbf{Bus (Système d'interconnexion)} & \emph{Les « couloirs et autoroutes »} : voies de communication partagées (fils conducteurs) transportant adresses, données, signaux de contrôle. & \textbf{Bus système} (CPU $\leftrightarrow$ RAM/Chipset), \textbf{Bus d'extension} (PCIe pour GPU/SSD NVMe), \textbf{Bus périphériques} (USB 3.2/4, SATA, I2C, SPI). \\
\hline
\end{tabularx}
\end{center}

\courssub{Schéma global de l'architecture de Von Neumann}

Voici la représentation canonique que tu dois savoir reproduire à main levée (schéma fonctionnel, pas photo réaliste) :

\begin{popfigure}
\begin{tikzpicture}[
    node distance=1.8cm and 2.5cm,
    block/.style={draw=popdark, thick, rounded corners, fill=popblueL, minimum height=1.8cm, minimum width=3.2cm, align=center, font=\small},
    bus/.style={draw=poporange, very thick, -{Stealth[length=3mm]}, shorten >=2pt, shorten <=2pt},
    busbi/.style={draw=poporange, very thick, -{Stealth[length=3mm]}, double, double distance=1.5pt, shorten >=2pt, shorten <=2pt},
    labelstyle/.style={font=\footnotesize, text=popdark, midway, sloped, above},
]
% Nœuds principaux
\node[block, align=center] (cpu){\textbf{UNITÉ CENTRALE (CPU)}\\\texttt{ALU + Unité de contrôle + Registres + Cache}};
\node[block, right=of cpu, align=center] (mem){\textbf{MÉMOIRE CENTRALE (RAM)}\\\texttt{Programmes + Données}\\\texttt{Adresses 0 $\to$ N}};
\node[block, below=of cpu, align=center] (io){\textbf{UNITÉS D'ENTRÉE / SORTIE}\\\texttt{Clavier, Souris, Écran, Disque SSD, Réseau, USB}};

% Bus système (triangle CPU-MEM-IO)
\draw[busbi] (cpu.north east) -- node[labelstyle] {\texttt{Bus de données (bidirectionnel)}} (mem.north west);
\draw[bus] (cpu.east |- cpu.north) -- ++(0.3,0) |- node[labelstyle, pos=0.7, above] {\texttt{Bus d'adresses (unidirectionnel CPU vers MEM/IO)}} (mem.west |- mem.north);
\draw[bus] (cpu.east |- cpu.south) -- ++(0.3,0) |- node[labelstyle, pos=0.3, below] {\texttt{Bus de contrôle (unidirectionnel CPU vers MEM/IO)}} (io.west |- io.north);

% Flèches E/S vers l'extérieur
\draw[bus] (io.west) -- ++(-1.2,0) node[left, font=\footnotesize, text=popmuted, align=center]{Périphériques\\externes};
\draw[bus] (io.east) -- ++(1.2,0) node[right, font=\footnotesize, text=popmuted] {Réseau / Internet};

% Légendes internes CPU (optionnel, pour richesse)
\node[font=\tiny, text=popmuted, anchor=north] at (cpu.south) {Horloge : \SIrange{2.5}{4.5}{\giga\hertz} typique};
\node[font=\tiny, text=popmuted, anchor=north] at (mem.south) {Volatile -- Accès octet/mot};
\node[font=\tiny, text=popmuted, anchor=north] at (io.south) {Contrôleurs dédiés (GPU, NIC, USB Ctrl)};
\end{tikzpicture}
\legende{Architecture de Von Neumann : les 4 unités fonctionnelles reliées par le bus système (données, adresses, contrôle).}
\end{popfigure}

\courssub{Architecture matérielle vs Pilotage logiciel : la distinction cruciale}

\begin{propriete}[Séparation Architecture / Microarchitecture / Logiciel]
\begin{itemize}
    \item \textbf{Architecture (ISA -- Instruction Set Architecture)} : le \emph{contrat} visible du programmeur (jeu d'instructions, registres, modes d'adressage, modèle mémoire). Ex : \texttt{x86-64}, \texttt{ARMv8-A}, \texttt{RISC-V}. C'est \textbf{quoi} la machine sait faire.
    \item \textbf{Microarchitecture (Organisation interne)} : l'\emph{implémentation} concrète d'une ISA (pipeline, prédiction de branchement, nombre de cœurs, taille des caches, fréquence). Ex : Intel \emph{Golden Cove} vs AMD \emph{Zen 4} (tous deux x86-64 mais microarchitectures différentes). C'est \textbf{comment} elle le fait (performance, consommation).
    \item \textbf{Logiciel (Système + Applications)} : la suite d'instructions binaires qui \emph{pilote} le matériel via l'ISA. Le même programme compilé pour x86 ne tourne \textbf{pas} sur ARM sans recompilation (ou émulation lente).
\end{itemize}
\end{propriete}

\begin{exemplebox}[Identifier les 4 unités sur le poste de la salle multimédia]
\emph{Consigne :} Devant l'ordinateur du collège (tour HP ProDesk 400 G9, écran 24''), pointe du doigt et nomme :
\begin{enumerate}
    \item \textbf{CPU} : sous le gros ventirad (radiateur + ventilateur) sur la carte mère, socket LGA1700, puce Intel Core i3-12100 (4 cœurs, ISA x86-64).
    \item \textbf{RAM} : 1 barrette DDR4 \SI{3200}{\mega\hertz} de \SI{8}{Go} insérée dans le slot DIMM noir (volatile, perte données si coupure courant).
    \item \textbf{E/S} : \emph{Entrée} : clavier/souris USB (contrôleur USB dans chipset), webcam USB. \emph{Sortie} : écran DisplayPort (GPU intégré UHD 730), haut-parleurs jack. \emph{E/S mixte} : SSD M.2 NVMe \SI{256}{Go} (disque système), port RJ45 (carte réseau Realtek).
    \item \textbf{Bus} : invisibles sur la carte mère (pistes cuivrées sous le vernis vert) : bus DDR4 vers RAM, bus PCIe 4.0 x4 vers SSD, bus DMI 4.0 vers chipset (qui gère USB, SATA, Audio, LAN).
\end{enumerate}
\emph{Observation :} Si tu ouvres le \texttt{Gestionnaire de tâches} (Ctrl+Maj+\'Echap) $\rightarrow$ onglet \texttt{Performance}, tu vois en temps réel l'activité du \textbf{CPU} (graphique \%), de la \textbf{Mémoire} (Go utilisés / Go totaux), du \textbf{Disque} (SSD, débit Mo/s) et du \textbf{Réseau} (Wi-Fi/Ethernet, Mbit/s). C'est la visualisation logicielle de l'architecture matérielle !
\end{exemplebox}

\courssub{Exercice d'application résolu}

\begin{exoresolu}[Analyse d'une fiche technique réelle]
\textbf{Énoncé :} Tu veux acheter un ordinateur portable pour tes études (bureautique, navigation, programmation légère Python/C). Voici deux fiches techniques simplifiées. Lequel choisir et \textbf{justifie} ton choix en reliant les composants à l'architecture Von Neumann.

\begin{center}
\begin{tabularx}{\linewidth}{|l|X|X|}
\hline
\rowcolor{popblueL}
\textbf{Composant} & \textbf{Modèle A (Entrée de gamme)} & \textbf{Modèle B (Milieu de gamme)} \\
\hline
CPU & Intel Celeron N4500 (2 cœurs, 2 threads, \SI{1.1}{\giga\hertz}--\SI{2.8}{\giga\hertz}, \SI{4}{Mo} cache) & AMD Ryzen 5 7530U (6 cœurs, 12 threads, \SI{2.0}{\giga\hertz}--\SI{4.5}{\giga\hertz}, \SI{16}{Mo} cache) \\
RAM & \SI{4}{Go} DDR4 \SI{2400}{\mega\hertz} \textbf{soudée} (non extensible) & \SI{16}{Go} DDR4 \SI{3200}{\mega\hertz} (1 slot libre, extensible à \SI{32}{Go}) \\
Stockage & \SI{128}{Go} eMMC (soudé, lent \SI{~250}{Mo/s}) & \SI{512}{Go} SSD NVMe PCIe 3.0 (amovible, rapide \SI{~3000}{Mo/s}) \\
Écran & 14'' HD ($1366\times768$) TN & 14'' Full HD ($1920\times1080$) IPS \\
\hline
\end{tabularx}
\end{center}

\tcblower
\textbf{Solution argumentée :}

\textbf{Choix : Modèle B sans hésitation.} Voici l'analyse architecturale :

\begin{enumerate}
    \item \textbf{CPU (Unité de traitement)} : Le Ryzen 5 (6 cœurs / 12 threads, cache \SI{16}{Mo}, fréquence max \SI{4.5}{\giga\hertz}) offre une \textbf{puissance de calcul 5 à 8 fois supérieure} au Celeron (2 cœurs / 2 threads, cache \SI{4}{Mo}, \SI{2.8}{\giga\hertz}). Pour la compilation (gcc, javac) et le multitâche (navigateur + IDE + terminal), le parallélisme (threads) et le cache L3 évitent les attentes mémoire (goulot Von Neumann).
    \item \textbf{RAM (Mémoire centrale)} : \SI{4}{Go} est \textbf{critiquement insuffisant} en 2025 : Windows 10/11 + navigateur (3 onglets) + VS Code $\approx$ \SI{3.5}{Go}--\SI{4}{Go} utilisés $\rightarrow$ \textbf{swapping} constant sur disque (gel de la machine). \SI{16}{Go} assure la fluidité. L'extensibilité (slot libre) prolonge la durée de vie (investissement durable).
    \item \textbf{Stockage (E/S mixte)} : eMMC \SI{128}{Go} = technologie carte SD, contrôleur simple, \textbf{débit séquentiel \SI{~250}{Mo/s}}, usure rapide, non remplaçable. SSD NVMe \SI{512}{Go} = contrôleur avancé, \textbf{débit \SI{~3000}{Mo/s}} (12$\times$ plus vite), remplaçable. Temps de démarrage : ~45 s (eMMC) vs ~10 s (NVMe). Ouverture d'applications : instantanée sur NVMe.
    \item \textbf{Bus \& Équilibre global} : Le modèle A crée un \textbf{déséquilibre architectural} : CPU faible + RAM étriquée + stockage lent $\rightarrow$ le maillon le plus faible (stockage/RAM) bloque le CPU. Le modèle B est \textbf{équilibré} : CPU performant alimenté par RAM suffisante et stockage rapide.
\end{enumerate}

\textbf{Conclusion :} Le prix plus élevé du Modèle B s'explique par une \textbf{architecture cohérente} (pas de goulot d'étranglement), garante de productivité sur 4--5 ans. Le Modèle A est une « machine à frustration » pour un étudiant.
\end{exoresolu}

\courssub{Synthèse de la section}

\begin{aretenir}[Ce qu'il faut retenir de la section 1]
\begin{itemize}
    \item L'\textbf{architecture} est le plan logique d'organisation du matériel (invisible de l'extérieur).
    \item Le \textbf{modèle de Von Neumann} (1945) impose : \textbf{mémoire unifiée} (programme + données), \textbf{programme enregistré}, \textbf{exécution séquentielle pilotée par une unité de contrôle}.
    \item \textbf{4 unités fonctionnelles} : \textbf{CPU} (calcul/contrôle), \textbf{RAM} (stockage rapide volatile), \textbf{E/S} (interface monde extérieur), \textbf{Bus} (interconnexion adresses/données/contrôle).
    \item Distinction \textbf{ISA (architecture logique)} vs \textbf{Microarchitecture (implémentation physique)} vs \textbf{Logiciel (pilote)}.
    \item Un bon choix d'ordinateur repose sur l'\textbf{équilibre architectural} (CPU, RAM, Stockage, Bus cohérents) et non sur un seul composant « marketing ».
\end{itemize}
\end{aretenir}

\begin{exercice}[Vérification des acquis -- Section 1]
\begin{enumerate}
    \item Explique par une analogie (autre que la maison) la différence entre \textbf{architecture} et \textbf{boîtier}.
    \item Cite les \textbf{trois principes} du modèle de Von Neumann. Pourquoi dit-on « machine à programme enregistré » ?
    \item Sur le schéma de Von Neumann, trace le chemin suivi par une instruction \texttt{ADD A, B} depuis la mémoire jusqu'à son exécution dans l'ALU, en nommant les bus empruntés.
    \item Un PC a un CPU très puissant (i9 dernier cri) mais seulement \SI{4}{Go} de RAM DDR3 et un disque dur mécanique 5400 tr/min. Quel \textbf{principe architectural} est violé ? Quel symptôme l'utilisateur observera-t-il ?
    \item \textbf{Recherche personnelle} : Quelle est l'architecture (ISA) du processeur de ton smartphone ? (Indice : \texttt{Paramètres $\to$ À propos du téléphone $\to$ Informations matériel / Processeur}). Est-ce du RISC ou du CISC ? (On verra la différence en section 5).
\end{enumerate}
\end{exercice}

\coursec{L'unité centrale de traitement (CPU)}

Dans la section précédente, nous avons découvert que l'architecture d'un ordinateur décrit l'organisation de ses composants : l'unité centrale, les mémoires, les périphériques et les bus qui les relient. Nous avons vu que tous ces éléments coopèrent pour traiter l'information. Mais un composant joue un rôle tellement central qu'il mérite une étude à part entière : le \cle{processeur}. C'est lui qui exécute réellement les programmes. Ouvrons donc le capot et observons ce qui se passe à l'intérieur de cette petite puce de quelques centimètres carrés.

\courssub{Définition et rôle du CPU}

\textbf{Une situation de la vie courante.}\par

Imagine le proviseur d'un collège de Yaoundé le jour de la rentrée. Les élèves arrivent avec leurs dossiers (les données), les enseignants appliquent les consignes (les programmes), les surveillants rangent les salles (les périphériques). Mais qui décide, coordonne et vérifie que tout se déroule correctement ? Le proviseur. Dans un ordinateur, ce « proviseur » s'appelle le processeur : sans lui, la mémoire, le disque dur et l'écran ne sont que des objets inertes.

\begin{definition}[Le CPU ou processeur]
L'\cle{unité centrale de traitement}, appelée \cle{CPU} (de l'anglais \emph{Central Processing Unit}) ou \cle{processeur}, est le composant de l'ordinateur qui \textbf{exécute les instructions des programmes}. Il lit les instructions et les données depuis la mémoire, les traite, puis produit les résultats. On le surnomme souvent le \emph{cerveau de l'ordinateur}.
\end{definition}

Concrètement, le processeur est une puce électronique (un circuit intégré) enfichée sur la carte mère. Quand tu double-cliques sur l'icône d'un logiciel de traitement de texte, c'est le processeur qui exécute une à une les millions d'instructions de ce logiciel : afficher une fenêtre, attendre une frappe au clavier, afficher la lettre tapée, etc.

\begin{aretenir}[L'essentiel sur le CPU]
Le CPU est le composant qui \textbf{exécute les instructions}. Tout ce qu'un ordinateur fait (calculer, afficher, enregistrer, jouer un son) passe obligatoirement par lui. Il est constitué de trois grandes parties : l'unité de commande, l'unité arithmétique et logique, et les registres.
\end{aretenir}

\courssub{L'unité de commande (UC)}

\textbf{Intuition : le chef d'orchestre.}\par

Dans une fanfare, chaque musicien sait jouer de son instrument, mais c'est le chef d'orchestre qui indique \emph{quand} jouer et \emph{quelle} note. De même, dans le processeur, les calculs ne se font pas n'importe comment : un composant supervise tout le travail.

\begin{definition}[L'unité de commande]
L'\cle{unité de commande} (UC), appelée aussi unité de contrôle, est la partie du processeur qui \textbf{lit les instructions en mémoire, les décode} (c'est-à-dire les comprend) et \textbf{coordonne le travail de tous les autres composants} : elle ordonne à l'UAL de calculer, demande à la mémoire de fournir une donnée, ou commande l'envoi d'un résultat vers un périphérique.
\end{definition}

L'unité de commande ne calcule rien elle-même : elle \emph{dirige}. Son travail se déroule en trois temps :
\begin{enumerate}
\item \textbf{Lire} : elle va chercher la prochaine instruction dans la mémoire vive.
\item \textbf{Décoder} : elle analyse l'instruction pour comprendre ce qu'elle demande (une addition ? une comparaison ? un déplacement de donnée ?).
\item \textbf{Commander} : elle envoie les ordres aux composants concernés pour que l'instruction soit exécutée.
\end{enumerate}

\courssub{L'unité arithmétique et logique (UAL)}

\textbf{Intuition : la calculatrice du processeur.}\par

Quand un élève résout un problème de mathématiques, il réfléchit (il décide quoi faire), puis il calcule (additions, comparaisons). Si l'unité de commande est la partie qui « réfléchit et organise », il faut bien une partie qui « calcule » : c'est l'UAL.

\begin{definition}[L'unité arithmétique et logique]
L'\cle{unité arithmétique et logique} (\cle{UAL}, en anglais \emph{ALU}) est la partie du processeur qui \textbf{effectue les opérations} :
\begin{itemize}
\item les \textbf{opérations arithmétiques} : addition, soustraction, multiplication, division ;
\item les \textbf{opérations logiques} : comparaisons (égal, plus grand, plus petit) et opérations logiques (ET, OU, NON).
\end{itemize}
\end{definition}

Chaque fois qu'un programme calcule la moyenne d'un élève, vérifie si un mot de passe est correct ou teste si un stock est épuisé, c'est l'UAL qui effectue l'opération, sur ordre de l'unité de commande.

\begin{exemplebox}[Qui fait quoi dans le processeur ?]
Un programme de bulletin de notes calcule : $moyenne = (note1 + note2) / 2$.
\begin{itemize}
\item L'\textbf{unité de commande} lit l'instruction, la décode et comprend qu'il faut additionner puis diviser ;
\item l'\textbf{UAL} effectue l'addition $note1 + note2$, puis la division par $2$ ;
\item l'unité de commande ordonne ensuite l'enregistrement du résultat en mémoire.
\end{itemize}
\end{exemplebox}

\courssub{Les registres : les mémoires internes du processeur}

\textbf{Intuition : l'ardoise de l'élève.}\par

Pour calculer $12 \times 8 + 4$, un élève retient le résultat intermédiaire $96$ dans sa tête ou le note sur son ardoise avant d'ajouter $4$. Le processeur a le même besoin : il doit conserver temporairement, \emph{à l'intérieur de lui-même}, les données sur lesquelles il travaille, car aller les chercher en mémoire vive à chaque instant serait beaucoup trop lent. Ces petites mémoires internes s'appellent les registres.

\begin{definition}[Registre]
Un \cle{registre} est une \textbf{petite mémoire très rapide située à l'intérieur du processeur}, qui stocke temporairement une donnée, une instruction ou une adresse pendant le traitement. Les registres sont les mémoires les plus rapides de l'ordinateur, mais aussi les plus petites.
\end{definition}

Trois registres sont à connaître en particulier :

\begin{itemize}
\item le \cle{compteur ordinal} (CO, en anglais \emph{Program Counter}) : il contient l'\textbf{adresse de la prochaine instruction à exécuter}. Grâce à lui, le processeur sait toujours où il en est dans le programme ;
\item le \cle{registre d'instruction} (RI) : il contient l'\textbf{instruction en cours d'exécution}, le temps que l'unité de commande la décode ;
\item l'\cle{accumulateur} : il contient les \textbf{données et les résultats} des opérations effectuées par l'UAL.
\end{itemize}

\begin{definition}[Le compteur ordinal]
Le \cle{compteur ordinal} est le registre qui mémorise l'\textbf{adresse en mémoire de la prochaine instruction à exécuter}. Après chaque instruction, il est automatiquement mis à jour pour pointer vers l'instruction suivante : c'est lui qui garantit que le programme avance dans le bon ordre.
\end{definition}

\begin{popfigure}
\begin{center}
\begin{tikzpicture}[
  box/.style={draw=popblue, thick, fill=popblueL, rounded corners=2pt, minimum width=3.2cm, minimum height=1cm, align=center, font=\small},
  reg/.style={draw=poppurple, thick, fill=poppurpleL, rounded corners=2pt, minimum width=2.6cm, minimum height=0.7cm, align=center, font=\footnotesize},
  mem/.style={draw=popgreen, thick, fill=popgreenL, rounded corners=2pt, minimum width=3.6cm, minimum height=1.2cm, align=center, font=\small},
  lien/.style={-{Stealth[length=2.5mm]}, thick, popdark},
  cpu/.style={draw=poporange, very thick, dashed, rounded corners=6pt, inner sep=8pt}
]
% Mémoire centrale (hors CPU)
\node[mem, align=center] (ram) at (0,4.6){\textbf{Mémoire centrale (RAM)}\\ programme + données};
% Composants du CPU
\node[box, align=center] (uc) at (-3.2,1.6){\textbf{Unité de}\\ \textbf{commande (UC)}\\ \footnotesize lit, décode, coordonne};
\node[box, align=center] (ual) at (3.2,1.6){\textbf{UAL}\\ \footnotesize calculs et comparaisons};
\node[reg] (co) at (-3.2,-0.2) {Compteur ordinal (CO)};
\node[reg] (ri) at (0,-0.2) {Registre d'instruction (RI)};
\node[reg] (acc) at (3.2,-0.2) {Accumulateur};
% Cadre CPU
\begin{scope}[on background layer]
\node[cpu, fit=(uc)(ual)(co)(ri)(acc), label={[poporange, font=\small\bfseries]below:Processeur (CPU)}] {};
\end{scope}
% Liaisons
\draw[lien] (ram.south -| uc) -- (uc.north) node[midway, left, font=\footnotesize, align=center]{instruction};
\draw[lien] (uc.south) -- (co.north);
\draw[lien] (uc.east) -- (ual.west) node[midway, above, font=\footnotesize]{ordres};
\draw[lien] (ri.north) -- (uc.south east);
\draw[lien] (acc.north) -- (ual.south);
\draw[lien, popgreen] (ram.south east) to[out=-60,in=60] node[midway, right, font=\footnotesize, align=center]{données} (acc.east);
\end{tikzpicture}
\end{center}
\legende{Schéma interne du CPU : l'unité de commande pilote l'UAL et les registres, et dialogue avec la mémoire centrale.}
\end{popfigure}

\courssub{L'horloge et la fréquence}

\textbf{Intuition : le métronome.}\par

Un orchestre joue en suivant un tempo. Le processeur aussi : ses opérations ne se déroulent pas n'importe quand, mais au rythme d'un signal régulier appelé \cle{horloge}. Chaque « tic » de l'horloge autorise une étape de travail : c'est un \cle{cycle}.

\begin{definition}[Horloge et fréquence]
L'\cle{horloge} du processeur est un signal électrique régulier qui \textbf{rythme l'exécution des instructions}. La \cle{fréquence} est le nombre de cycles (de « tics ») par seconde. Elle s'exprime en \textbf{hertz} (Hz) et ses multiples :
\begin{itemize}
\item $1$ MHz $= 1$ million de cycles par seconde ;
\item $1$ GHz $= 1$ milliard de cycles par seconde.
\end{itemize}
\end{definition}

Plus la fréquence est élevée, plus le processeur enchaîne vite les cycles, et donc plus il exécute d'instructions en une seconde.

\begin{exemplebox}[Un processeur à 2,4 GHz]
Un processeur annoncé à \SI{2,4}{GHz} effectue environ $2,4$ \textbf{milliards de cycles par seconde} :
\[ 2,4 \text{ GHz} = 2,4 \times 10^{9} \text{ cycles/seconde} = 2\,400\,000\,000 \text{ cycles/seconde.} \]
Autrement dit, un cycle dure environ :
\[ \frac{1}{2,4 \times 10^{9}} \approx 0,42 \text{ nanoseconde,} \]
soit moins d'un demi-milliardième de seconde ! C'est dans ce laps de temps minuscule que le processeur lit, décode ou exécute une étape d'instruction.
\end{exemplebox}

\begin{attention}[La fréquence ne dit pas tout !]
Erreur fréquente : croire qu'un processeur à \SI{3}{GHz} est \emph{forcément} plus rapide qu'un processeur à \SI{2,4}{GHz}. C'est faux en général, car la rapidité dépend aussi :
\begin{itemize}
\item du \textbf{nombre de cœurs} : plusieurs cœurs permettent de traiter plusieurs tâches en même temps ;
\item de l'\textbf{architecture interne} : à fréquence égale, un processeur récent accomplit souvent plus de travail par cycle qu'un ancien.
\end{itemize}
La fréquence n'est donc qu'\emph{un} indice de performance parmi d'autres.
\end{attention}

\begin{savaistu}[Les processeurs à plusieurs cœurs]
Les processeurs modernes possèdent \textbf{plusieurs cœurs} (\emph{dual-core}, \emph{quad-core}, etc.) : ce sont en réalité \textbf{plusieurs unités de traitement réunies dans une seule puce}. Un processeur à 4 cœurs peut exécuter 4 programmes (ou 4 parties d'un même programme) \emph{en même temps}, comme 4 élèves qui résoudraient chacun un exercice différent sur la même table. C'est pourquoi un téléphone ou un ordinateur récent reste fluide même quand plusieurs applications sont ouvertes.
\end{savaistu}

\courssub{Lire l'étiquette d'un processeur}

Quand on achète un ordinateur dans une boutique de Douala ou qu'on consulte les propriétés d'un poste du CDI, on trouve une description commerciale du processeur. Savoir la lire est un savoir-faire concret.

\begin{methode}[Lire les caractéristiques d'un processeur]
\begin{enumerate}
\item Repérer la \textbf{marque et le modèle} : par exemple Intel (Core i3, i5, i7) ou AMD (Ryzen 3, 5, 7).
\item Repérer la \textbf{fréquence} : exprimée en GHz, elle indique le rythme de l'horloge.
\item Repérer le \textbf{nombre de cœurs} : il indique combien d'unités de traitement la puce contient.
\item Conclure en comparant \emph{ensemble} ces informations, jamais la fréquence seule.
\end{enumerate}
\end{methode}

\begin{exoresolu}[Lecture d'une étiquette de machine]
Énoncé. Sur l'affiche d'un magasin, on lit : « Ordinateur HP — Processeur Intel Core i5, 4 cœurs, 2,4 GHz ». Décris le processeur de cette machine et calcule le nombre de cycles qu'il effectue en une seconde.
\tcblower
Solution détaillée.
\begin{itemize}
\item \textbf{Marque et modèle} : Intel Core i5, un processeur de milieu de gamme.
\item \textbf{Nombre de cœurs} : 4 cœurs, donc 4 unités de traitement dans une seule puce, capables de travailler simultanément.
\item \textbf{Fréquence} : \SI{2,4}{GHz}, c'est-à-dire que l'horloge bat $2,4$ milliards de fois par seconde.
\end{itemize}
Nombre de cycles par seconde :
\[ 2,4 \text{ GHz} = 2,4 \times 10^{9} = 2\,400\,000\,000 \text{ cycles par seconde.} \]
Conclusion : ce processeur effectue environ $2,4$ milliards de cycles par seconde, répartis sur 4 cœurs. Pour juger de sa rapidité réelle, il faut tenir compte des 4 cœurs et de l'architecture, pas seulement des \SI{2,4}{GHz}.
\end{exoresolu}

\begin{exercice}[À toi de jouer]
Dans le CDI de ton collège, un poste affiche : « AMD Ryzen 3, 2 cœurs, 3,1 GHz ».
\begin{enumerate}
\item Donne la marque, le nombre de cœurs et la fréquence de ce processeur.
\item Combien de cycles ce processeur effectue-t-il en une seconde ?
\item Un camarade affirme : « Ce processeur à \SI{3,1}{GHz} est forcément plus rapide qu'un Core i5 à 4 cœurs et \SI{2,4}{GHz} ». Rédige une réponse argumentée pour nuancer cette affirmation.
\end{enumerate}
\end{exercice}

\begin{corrige}[Exercice : À toi de jouer]
\begin{enumerate}
\item Marque : AMD (modèle Ryzen 3) ; nombre de cœurs : 2 ; fréquence : \SI{3,1}{GHz}.
\item $3,1 \text{ GHz} = 3,1 \times 10^{9} = 3\,100\,000\,000$ cycles par seconde.
\item Cette affirmation est trop rapide. Certes, \SI{3,1}{GHz} est supérieur à \SI{2,4}{GHz}, mais le Ryzen 3 n'a que 2 cœurs contre 4 pour le Core i5 : ce dernier peut traiter davantage de tâches simultanément. De plus, l'architecture interne (le travail accompli à chaque cycle) diffère d'un modèle à l'autre. La fréquence seule ne détermine donc pas la rapidité : il faut comparer aussi le nombre de cœurs et l'architecture.
\end{enumerate}
\end{corrige}

\begin{aretenir}[Bilan de la section]
\begin{itemize}
\item Le \textbf{CPU} (processeur) exécute les instructions des programmes : c'est le cerveau de l'ordinateur.
\item L'\textbf{unité de commande} lit, décode les instructions et coordonne les autres composants.
\item L'\textbf{UAL} effectue les calculs arithmétiques et les comparaisons logiques.
\item Les \textbf{registres} (compteur ordinal, registre d'instruction, accumulateur) sont de petites mémoires internes très rapides.
\item L'\textbf{horloge} rythme le travail du processeur ; la \textbf{fréquence}, en GHz, compte les cycles par seconde.
\item Pour évaluer un processeur, on regarde la marque, le nombre de cœurs \emph{et} la fréquence.
\end{itemize}
\end{aretenir}

\coursec{Les mémoires et les bus}

Dans la section précédente, nous avons découvert l'\cle{unité centrale de traitement} (CPU), véritable « cerveau » de l'ordinateur, capable d'effectuer des calculs et de commander les autres composants. Mais un cerveau sans mémoire ne sert à rien : pour calculer, le processeur a besoin d'avoir sous la main les \cle{données} et les \cle{programmes}. Où sont-ils rangés ? Comment voyagent-ils jusqu'au processeur ? C'est toute la question de cette section : nous allons étudier les \cle{mémoires} de l'ordinateur et les \cle{bus}, ces « routes » qui relient tous les composants entre eux.

\courssub{Le rôle général des mémoires}

\begin{definition}[Mémoire]
Une \cle{mémoire} est un composant électronique capable de \textbf{stocker} (conserver) des informations : des \textbf{données} (nombres, textes, images, sons) et des \textbf{programmes} (suites d'instructions). On peut y \textbf{écrire} des informations et les y \textbf{lire} ensuite.
\end{definition}

Pourquoi un ordinateur possède-t-il \emph{plusieurs} types de mémoires, et pas une seule ? Parce qu'aucune mémoire n'est parfaite sur tous les plans : les mémoires très rapides coûtent cher et sont petites, les mémoires très grandes sont lentes. L'ordinateur combine donc plusieurs mémoires complémentaires, chacune spécialisée dans un rôle. On distingue trois grandes familles : la mémoire vive, la mémoire morte et les mémoires de masse.

\courssub{La mémoire vive (RAM)}

Imagine que tu fais tes devoirs sur une table : les cahiers ouverts sur la table sont ceux sur lesquels tu travailles \emph{en ce moment}. La mémoire vive joue le rôle de cette table de travail.

\begin{definition}[Mémoire vive (RAM)]
La \cle{mémoire vive}, appelée \cle{RAM} (\emph{Random Access Memory}, mémoire à accès aléatoire), est la mémoire de \textbf{travail} de l'ordinateur. Elle stocke \textbf{temporairement} les programmes en cours d'exécution et les données qu'ils manipulent. Sa capacité se mesure en gigaoctets (Go) : 2 Go, 4 Go, 8 Go, etc.
\end{definition}

Quand tu ouvres un logiciel (par exemple un traitement de texte), le programme est \textbf{copié du disque dur vers la RAM}, puis le processeur travaille avec cette copie, beaucoup plus rapide d'accès. C'est pourquoi un ordinateur avec plus de RAM peut faire tourner confortablement plusieurs programmes en même temps : la « table de travail » est plus grande.

\begin{propriete}[Volatilité de la RAM (admise)]
La RAM est une mémoire \cle{volatile} : son contenu \textbf{disparaît dès que l'ordinateur est éteint} (ou débranché). En revanche, la ROM et les mémoires de masse (disque dur, SSD, clé USB) sont \cle{non volatiles} : elles \textbf{conservent} leurs informations même sans courant électrique.
\end{propriete}

\begin{attention}[Travail non enregistré = travail perdu]
Conséquence pratique de la volatilité de la RAM : si tu rédiges un devoir au traitement de texte et que le courant se coupe (délestage !) avant que tu aies \textbf{enregistré} ton fichier, tout ce qui était dans la RAM est perdu. Enregistre donc régulièrement ton travail sur le disque dur (touche \texttt{Ctrl+S}) !
\end{attention}

\courssub{La mémoire morte (ROM)}

Quand tu appuies sur le bouton d'allumage, la RAM est vide (elle est volatile). Qui donne alors au processeur ses toutes premières instructions ? C'est le rôle de la mémoire morte.

\begin{definition}[Mémoire morte (ROM)]
La \cle{mémoire morte}, appelée \cle{ROM} (\emph{Read Only Memory}, mémoire en lecture seule), est une mémoire \textbf{permanente} et \textbf{non volatile}. Son contenu est écrit une fois pour toutes par le fabricant et ne s'efface pas, même ordinateur éteint. Elle contient le \cle{programme de démarrage} de l'ordinateur, appelé \cle{BIOS} (\emph{Basic Input Output System}), qui teste les composants à l'allumage puis lance le chargement du système d'exploitation.
\end{definition}

On dit « lecture seule » car l'utilisateur ne peut pas modifier son contenu en utilisation normale : on peut seulement le \textbf{lire}. C'est une sécurité : le programme de démarrage est ainsi protégé contre tout effacement accidentel.

\courssub{Les mémoires de masse}

Revenons à notre comparaison : si la RAM est la table de travail, les mémoires de masse sont l'\textbf{armoire} où l'on range tous ses cahiers, livres et affaires de façon durable.

\begin{definition}[Mémoire de masse]
Une \cle{mémoire de masse} (ou mémoire de stockage) est une mémoire \textbf{non volatile} de grande capacité, destinée à conserver \textbf{durablement} les fichiers de l'utilisateur (documents, photos, vidéos) et les programmes installés. Exemples : le \cle{disque dur}, le \cle{SSD}, la \cle{clé USB}, la \cle{carte mémoire} (carte SD d'un téléphone).
\end{definition}

\begin{center}
\begin{tabularx}{\linewidth}{|l|X|X|}\hline
\rowcolor{popblueL} \textbf{Mémoire de masse} & \textbf{Principe} & \textbf{Caractéristiques} \\ \hline
Disque dur & Plateaux magnétiques en rotation lus par une tête de lecture & Grande capacité (500 Go, 1 To), bon marché, mais fragile et assez lent \\ \hline
SSD & Mémoire flash (électronique, sans pièce mobile) & Rapide, silencieux, résistant aux chocs, mais plus cher \\ \hline
Clé USB & Mémoire flash & Petite, transportable, capacité moyenne (8 à 128 Go) \\ \hline
Carte mémoire & Mémoire flash & Très petite, utilisée dans téléphones et appareils photo \\ \hline
\end{tabularx}
\end{center}

\courssub{La hiérarchie des mémoires}

Récapitulons : le processeur possède ses propres petites mémoires internes, les \cle{registres}, ultra-rapides mais minuscules ; puis vient la RAM, rapide mais volatile ; enfin les mémoires de masse, très grandes mais lentes. On obtient une \cle{hiérarchie des mémoires} que l'on représente par une pyramide.

\begin{propriete}[Loi de la hiérarchie des mémoires (admise)]
Plus une mémoire est \textbf{rapide}, plus elle est \textbf{chère} au gigaoctet et plus sa \textbf{capacité} est \textbf{petite}. Inversement, plus une mémoire est grande et bon marché, plus elle est lente. Ainsi : registres (très rapides, très petits) $\rightarrow$ RAM (rapide, quelques Go) $\rightarrow$ mémoires de masse (lentes, centaines de Go).
\end{propriete}

\begin{popfigure}
\begin{center}
\begin{tikzpicture}
  % Pyramide en 3 couches
  \fill[poporange!80]  (0,3.4) -- (-1.1,2.2) -- (1.1,2.2) -- cycle;
  \fill[popblue!70]    (-1.1,2.2) -- (-2.2,1.0) -- (2.2,1.0) -- (1.1,2.2) -- cycle;
  \fill[popgreen!70]   (-2.2,1.0) -- (-3.3,-0.2) -- (3.3,-0.2) -- (2.2,1.0) -- cycle;
  \draw[popdark, thick] (0,3.4) -- (-3.3,-0.2) -- (3.3,-0.2) -- cycle;
  \draw[popdark] (-1.1,2.2) -- (1.1,2.2);
  \draw[popdark] (-2.2,1.0) -- (2.2,1.0);
  % Étiquettes des couches
  \node[white, font=\small\bfseries] at (0,2.65) {Registres};
  \node[white, font=\small\bfseries] at (0,1.6) {RAM};
  \node[white, font=\small\bfseries] at (0,0.4) {Mémoires de masse (disque dur, SSD)};
  % Flèche vitesse (vers le haut)
  \draw[-{Stealth[length=3mm]}, line width=1.2pt, poppurple] (-4.6,-0.2) -- (-4.6,3.4)
      node[midway, left, align=center, font=\small, text=poppurple] {Vitesse\\croissante\\(et prix)};
  % Flèche capacité (vers le bas)
  \draw[-{Stealth[length=3mm]}, line width=1.2pt, poppink] (4.6,3.4) -- (4.6,-0.2)
      node[midway, right, align=center, font=\small, text=poppink] {Capacité\\croissante\\(Go, To)};
\end{tikzpicture}
\end{center}
\legende{La pyramide de la hiérarchie des mémoires : en haut, les registres du processeur (très rapides, très petits) ; en bas, les mémoires de masse (lentes, très grandes capacités).}
\end{popfigure}

\courssub{Les bus : les routes de l'information}

Nos composants sont maintenant connus : le CPU, la mémoire centrale (RAM et ROM), les mémoires de masse et les périphériques. Mais comment communiquent-ils ? Par des « routes » électriques appelées \cle{bus}.

\begin{definition}[Bus]
Un \cle{bus} est un ensemble de fils conducteurs (ou de pistes gravées sur la carte mère) qui \textbf{transportent les informations} entre les composants de l'ordinateur : processeur, mémoires et périphériques. On distingue trois types de bus selon la nature des informations transportées.
\end{definition}

\begin{propriete}[Les trois types de bus]
\begin{itemize}
  \item Le \cle{bus de données} transporte les \textbf{données} elles-mêmes (un nombre, un caractère, une instruction). Il est \textbf{bidirectionnel} : les données circulent dans les deux sens (lecture et écriture).
  \item Le \cle{bus d'adresses} transporte l'\textbf{adresse} de la case mémoire (ou du périphérique) concernée par l'échange. Il est \textbf{unidirectionnel} : c'est toujours le processeur qui indique l'adresse.
  \item Le \cle{bus de commande} (ou bus de contrôle) transporte les \textbf{ordres} : « lire », « écrire », ainsi que les signaux de synchronisation.
\end{itemize}
\end{propriete}

Une image pour bien comprendre : le bus d'adresses indique \emph{où} (le numéro de la maison), le bus de commande indique \emph{quoi faire} (déposer ou retirer un colis), et le bus de données transporte \emph{le colis} lui-même.

\begin{popfigure}
\begin{center}
\begin{tikzpicture}[
    bloc/.style={rectangle, rounded corners=3pt, draw=popdark, thick, align=center, font=\small\bfseries, minimum height=1.1cm},
    bus/.style={line width=2.2pt},
    lbl/.style={font=\footnotesize, align=center}]
  \node[bloc, fill=poporange!60, minimum width=3.2cm, align=center] (cpu) at (0,0){CPU\\(processeur)};
  \node[bloc, fill=popblue!50, minimum width=3.2cm, align=center] (mem) at (7.5,2.2){Mémoire centrale\\RAM + ROM};
  \node[bloc, fill=popgreen!55, minimum width=3.2cm, align=center] (per) at (7.5,-2.2){Périphériques\\clavier, écran, disque};
  % Bus de données (bidirectionnel)
  \draw[bus, popblue, {Stealth[length=3mm]}-{Stealth[length=3mm]}] (cpu.east) -- (mem.west);
  \draw[bus, popblue, {Stealth[length=3mm]}-{Stealth[length=3mm]}] (cpu.east) -- (per.west);
  \node[lbl, text=popblue] at (3.75,2.75) {Bus de données (bidirectionnel)};
  \node[lbl, text=popblue] at (3.75,-1.6) {Bus de données};
  % Bus d'adresses (unidirectionnel)
  \draw[bus, poppurple, -{Stealth[length=3mm]}] ([yshift=-0.35cm]cpu.east) -- ([yshift=-0.35cm]mem.west);
  \draw[bus, poppurple, -{Stealth[length=3mm]}] ([yshift=-0.35cm]cpu.east) -- ([yshift=-0.35cm]per.west);
  \node[lbl, text=poppurple] at (3.75,0.9) {Bus d'adresses (CPU $\rightarrow$ mémoire)};
  % Bus de commande
  \draw[bus, poppink, -{Stealth[length=3mm]}] ([yshift=-0.7cm]cpu.east) -- ([yshift=-0.7cm]mem.west);
  \draw[bus, poppink, -{Stealth[length=3mm]}] ([yshift=-0.7cm]cpu.east) -- ([yshift=-0.7cm]per.west);
  \node[lbl, text=poppink] at (3.75,-2.9) {Bus de commande (ordres lire / écrire)};
\end{tikzpicture}
\end{center}
\legende{Les trois bus relient le CPU à la mémoire centrale et aux périphériques : données (dans les deux sens), adresses et commandes (du CPU vers les autres composants).}
\end{popfigure}

\courssub{Les périphériques reliés par les bus}

Grâce aux bus, le processeur dialogue aussi avec les \cle{périphériques}. On les classe en trois catégories :

\begin{itemize}
  \item les périphériques d'\cle{entrée}, qui font \textbf{entrer} les informations dans l'ordinateur : clavier, souris, scanner, microphone ;
  \item les périphériques de \cle{sortie}, qui \textbf{restituent} les résultats : écran, imprimante, haut-parleurs ;
  \item les périphériques \cle{mixtes} (entrée/sortie), qui font les deux : l'\textbf{écran tactile} (il affiche et détecte les doigts), le disque dur, la clé USB, la carte réseau.
\end{itemize}

\begin{exoresolu}[Identifier mémoires, bus et périphériques]
Énoncé. Dans un cybercafé de Yaoundé, le gérant affirme : « Mon ordinateur a 4 Go de RAM et 500 Go de disque dur ; l'écran tactile de la caisse est un périphérique de sortie. » 1) Quel est le rôle de chacune des deux mémoires citées ? 2) Corrige l'erreur concernant l'écran tactile. 3) Par quels bus le processeur récupère-t-il un fichier client enregistré sur le disque dur ?
\tcblower
Solution détaillée.
\begin{enumerate}
  \item Les \textbf{4 Go de RAM} constituent la mémoire de travail : ils accueillent temporairement les programmes en cours d'exécution (navigateur, logiciel de facturation) et leurs données ; ce contenu disparaît à l'extinction. Les \textbf{500 Go de disque dur} forment la mémoire de masse : ils conservent durablement tous les fichiers, même ordinateur éteint.
  \item L'écran tactile n'est pas un simple périphérique de sortie : il \textbf{affiche} (sortie) mais aussi \textbf{détecte les touchers} du caissier (entrée). C'est donc un périphérique \textbf{mixte} (entrée/sortie).
  \item Le processeur place l'adresse du fichier sur le \textbf{bus d'adresses}, envoie l'ordre « lire » sur le \textbf{bus de commande}, et le contenu du fichier revient du disque dur vers la RAM par le \textbf{bus de données}.
\end{enumerate}
\end{exoresolu}

\courssub{Choisir les bonnes capacités selon l'usage}

\begin{methode}[Choisir la RAM et le disque selon l'usage]
Pour conseiller l'achat d'un ordinateur, procède en trois étapes :
\begin{enumerate}
  \item \textbf{Identifier l'usage principal} : bureautique (traitement de texte, tableur, Internet), jeux, ou montage vidéo.
  \item \textbf{Choisir la RAM en conséquence} : environ 4 Go suffisent pour la bureautique ; 8 Go sont conseillés pour les jeux ; 8 à 16 Go pour le montage vidéo, très gourmand.
  \item \textbf{Choisir le stockage selon la quantité de fichiers} : 250 à 500 Go pour des documents ; 1 To ou plus si l'on stocke beaucoup de vidéos et de jeux.
\end{enumerate}
Justifie toujours ta réponse en reliant la capacité choisie au besoin réel de l'utilisateur.
\end{methode}

\begin{exemplebox}[Comparer deux annonces du marché de Douala]
Sur une annonce d'ordinateur à Douala, on lit : « 4 Go de RAM, disque dur 500 Go ». Cela signifie que la \textbf{mémoire de travail} fait 4 Go : la machine convient à la bureautique et à Internet, mais ralentira si l'on ouvre trop de programmes à la fois. Les \textbf{500 Go de disque dur} permettent de ranger durablement des centaines de milliers de photos ou plus d'une centaine de films. Une seconde annonce « 8 Go de RAM, SSD 256 Go » offre une mémoire de travail double et un stockage plus rapide mais deux fois plus petit : elle convient mieux à un usage exigeant avec moins de fichiers volumineux.
\end{exemplebox}

\begin{attention}[Ne pas confondre RAM et stockage]
Erreur très fréquente : quand un téléphone annonce « \textbf{64 Go} », il s'agit de l'\textbf{espace de stockage} (mémoire de masse, pour les photos, applications et fichiers), et \textbf{pas} de la RAM ! La RAM du téléphone est beaucoup plus petite (souvent 2 à 8 Go) et est indiquée séparément. De même, « manque de mémoire » sur un téléphone signifie presque toujours « stockage plein », alors que la RAM n'y est pour rien. Retiens : \textbf{RAM = mémoire de travail (quelques Go, volatile)} ; \textbf{stockage = mémoire de masse (dizaines ou centaines de Go, permanente)}.
\end{attention}

\begin{aretenir}[L'essentiel de la section]
\begin{itemize}
  \item Les mémoires stockent les \textbf{données} et les \textbf{programmes}.
  \item La \textbf{RAM} (mémoire vive) est la mémoire de travail : rapide mais \textbf{volatile}.
  \item La \textbf{ROM} (mémoire morte) est permanente et contient le \textbf{BIOS}, le programme de démarrage.
  \item Les \textbf{mémoires de masse} (disque dur, SSD, clé USB, carte mémoire) conservent durablement les fichiers.
  \item Hiérarchie : plus une mémoire est rapide, plus elle est chère et petite (registres $\rightarrow$ RAM $\rightarrow$ disque).
  \item Les \textbf{bus} relient les composants : bus de \textbf{données} (bidirectionnel), bus d'\textbf{adresses} et bus de \textbf{commande}.
  \item Les périphériques sont d'\textbf{entrée}, de \textbf{sortie} ou \textbf{mixtes} (écran tactile).
\end{itemize}
\end{aretenir}

\begin{exercice}[À toi de jouer]
1) Un élève veut acheter un ordinateur pour le montage vidéo. On lui propose : (a) 2 Go de RAM, 1 To de disque ; (b) 8 Go de RAM, 500 Go de disque. Lequel conseilles-tu ? Justifie ta réponse en deux phrases. 2) Pour chaque périphérique, indique s'il est d'entrée, de sortie ou mixte : souris, imprimante, écran tactile, clé USB, haut-parleur. 3) Explique pourquoi, à l'allumage de l'ordinateur, le premier programme exécuté se trouve forcément dans la ROM et non dans la RAM.
\end{exercice}

\coursec{Le cycle d'exécution d'une instruction}

\courssub{Transition : du matériel à l'action}

Dans la section précédente, nous avons découvert les \textbf{mémoires} (RAM, ROM, cache) qui stockent les données et les programmes, et les \textbf{bus} (adresses, données, contrôle) qui servent d'autoroutes pour transporter l'information entre le processeur, la mémoire et les périphériques. Nous savons maintenant \emph{où} se trouvent les instructions et \emph{comment} elles circulent.

Mais un ordinateur ne se contente pas de stocker des instructions : il doit les \textbf{exécuter} les unes après les autres, sans arrêt, à une vitesse vertigineuse. Comment le processeur transforme-t-il une suite de bits stockée en mémoire en une action concrète (une addition, un affichage, un test) ? C'est précisément le rôle du \textbf{cycle d'exécution d'une instruction}, aussi appelé \textbf{cycle machine} ou \textbf{cycle \textit{Fetch-Decode-Execute}}. C'est le « battement de cœur » du processeur, la boucle infinie qui fait vivre n'importe quel programme, du plus simple script au jeu vidéo le plus complexe.

\courssub{Qu'est-ce qu'une instruction machine ?}

Avant de décrire le cycle, posons une définition précise. Un programme, aux yeux du processeur, n'est pas du texte lisible (comme du Python ou du C), mais une suite de \textbf{mots binaires} de longueur fixe (16, 32 ou 64 bits selon l'architecture).

\begin{definition}[Instruction machine]
Une \cle{instruction machine} est une opération élémentaire, codée en binaire, que le processeur (CPU) est capable de comprendre et d'exécuter directement sans traduction. Elle spécifie :
\begin{itemize}
    \item L'\textbf{opération} à réaliser (code opération ou \cle{opcode} : addition, chargement, saut, comparaison…).
    \item Les \textbf{opérandes} (données) sur lesquelles porte l'opération : soit des valeurs immédiates, soit des adresses de registres, soit des adresses mémoire.
\end{itemize}
L'ensemble des instructions qu'un processeur connaît constitue son \cle{jeu d'instructions} (ISA — \textit{Instruction Set Architecture}).
\end{definition}

\begin{attention}[Ne pas confondre instruction et ligne de code]
Une ligne de code en langage haut niveau (ex. \texttt{a = b + c} en C) se traduit généralement en \textbf{plusieurs} instructions machine (charger \texttt{b}, charger \texttt{c}, additionner, stocker dans \texttt{a}). L'instruction machine est l'unité atomique du travail du processeur.
\end{attention}

\courssub{Le cycle \textit{Fetch – Decode – Execute} : les 4 étapes fondamentales}

Le processeur répète inlassablement la même séquence de quatre étapes pour chaque instruction du programme. Ce cycle est piloté par l'\cle{unité de contrôle} (UC) à l'intérieur du CPU.

\begin{propriete}[Les 4 étapes du cycle d'exécution]
Pour chaque instruction, le processeur effectue dans l'ordre :
\begin{enumerate}
    \item \textbf{Fetch (Recherche / Chargement)} : L'UC va chercher l'instruction en mémoire à l'adresse indiquée par le \cle{compteur ordinal} (\cle{Program Counter} ou \cle{PC}), puis incrémente le PC pour qu'il pointe vers l'instruction suivante.
    \item \textbf{Decode (Décodage)} : L'UC analyse le code opération (opcode) de l'instruction pour déterminer quelle opération effectuer et identifie les opérandes (registres, adresses mémoire, valeurs immédiates). Elle prépare les chemins de données internes (bus internes, multiplexeurs) vers l'UAL.
    \item \textbf{Execute (Exécution)} : L'\cle{UAL} (\textit{Unité Arithmétique et Logique}) ou un autre organe spécialisé (FPU pour les flottants, unité de branchement) réalise l'opération demandée (calcul, test, déplacement de données, saut).
    \item \textbf{Writeback / Update PC (Écriture résultat \& Mise à jour PC)} : Le résultat est écrit dans le registre de destination ou en mémoire. Le compteur ordinal (PC) contient déjà l'adresse de l'instruction suivante (mis à jour à l'étape 1 ou modifié par un saut à l'étape 3). Le cycle recommence.
\end{enumerate}
\end{propriete}

\begin{popfigure}
\begin{tikzpicture}[
    >=Stealth,
    node distance=2.5cm and 3cm,
    box/.style={rectangle, rounded corners, draw=popdark, thick, fill=popcream, minimum width=3.2cm, minimum height=1.6cm, align=center, font=\small},
    arrow/.style={->, thick, draw=popblue},
    arrowback/.style={->, thick, draw=poporange, dashed}
]
% Nœuds
\node[box, align=center] (fetch){\textbf{1. FETCH}\\(Recherche)\\\textit{PC $\to$ MAR}\\\textit{Bus Adresses}\\\textit{Bus Données $\to$ MDR/RI}\\\textit{PC $\leftarrow$ PC + 1}};
\node[box, right=of fetch, align=center] (decode){\textbf{2. DECODE}\\(Décodage)\\\textit{UC analyse Opcode}\\\textit{Identifie opérandes}\\\textit{Prépare UAL / Registres}};
\node[box, below=of decode, align=center] (execute){\textbf{3. EXECUTE}\\(Exécution)\\\textit{UAL : Calcul / Logique}\\\textit{ou Accès Mémoire}\\\textit{ou Branchement}};
\node[box, left=of execute, align=center] (writeback){\textbf{4. WRITEBACK}\\(Écriture / MAJ PC)\\\textit{Résultat $\to$ Registre / Mém}\\\textit{Drapeaux (Flags) MAJ}\\\textit{PC prêt pour suivant}};

% Flèches principales
\draw[arrow] (fetch.east) -- node[above, font=\tiny] {Instruction brute} (decode.west);
\draw[arrow] (decode.south) -- node[right, font=\tiny] {Signaux de contrôle} (execute.north);
\draw[arrow] (execute.west) -- node[below, font=\tiny] {Résultat / Nouvel PC} (writeback.east);
\draw[arrow] (writeback.north) -- node[left, font=\tiny] {Prochaine instruction} (fetch.south);

% Flèche retour (boucle)
\draw[arrowback, bend left=15] (writeback.north west) to node[left=2mm, font=\tiny, text=poporange] {Boucle infinie} (fetch.north west);

% Légendes registres internes (optionnel, petit)
\node[below=0.2cm of fetch, font=\tiny, text=popmuted] {Registres : PC, MAR, MDR, RI};
\node[below=0.2cm of decode, font=\tiny, text=popmuted] {Unité de Contrôle (UC)};
\node[below=0.2cm of execute, font=\tiny, text=popmuted] {UAL, Registres Généraux};
\node[below=0.2cm of writeback, font=\tiny, text=popmuted] {Registres, Mémoire, Flags};
\end{tikzpicture}
\legende{Le cycle Fetch–Decode–Execute–Writeback : le cœur battant du processeur. Les flèches pleines montrent le flux principal ; la flèche pointillée orange rappelle la répétition continue.}
\end{popfigure}

\courssub{Détail pas à pas de chaque étape}

Regardons de plus près ce qui se passe dans les registres internes du processeur (architecture de type \textit{Von Neumann} à bus unique simplifiée).

\begin{methode}[Décrire le déroulement d'une instruction (Rédaction guidée)]
Pour expliquer l'exécution d'une instruction lors d'un contrôle ou à l'oral, suis ce canevas en 4 paragraphes :
\begin{enumerate}
    \item \textbf{Fetch} : Cite le registre \texttt{PC} (compteur ordinal), le bus d'adresses, la mémoire, le bus de données, le registre \texttt{RI} (Registre d'Instruction) ou \texttt{MDR}. Note l'incrémentation de \texttt{PC}.
    \item \textbf{Decode} : Cite l'\textbf{Unité de Contrôle (UC)}, l'analyse de l'\textbf{opcode}, l'identification des opérandes (registres sources, adresse mémoire, valeur immédiate), la génération des signaux de contrôle vers l'UAL et les multiplexeurs.
    \item \textbf{Execute} : Cite l'\textbf{UAL}, l'opération précise (addition, ET logique, décalage, comparaison…), les entrées (registres sources), la sortie (résultat), la mise à jour des \textbf{drapeaux} (Zero, Carry, Sign, Overflow).
    \item \textbf{Writeback / MAJ PC} : Indique où va le résultat (registre destination \texttt{Rd} ou mémoire via \texttt{MAR}/\texttt{MDR}). Précise que \texttt{PC} pointe déjà sur l'instruction suivante (sauf saut/branchement où \texttt{PC} a été chargé avec une nouvelle adresse à l'étape Execute).
\end{enumerate}
Utilise le vocabulaire précis : \emph{bus d'adresses}, \emph{bus de données}, \emph{bus de contrôle}, \emph{registre}, \emph{UAL}, \emph{drapeaux}.
\end{methode}

\courssub{Exemple concret : additionner 5 et 3}

Imaginons un petit programme en langage machine simplifié (architecture à accumulateur unique pour la clarté, comme un vieux 8080 ou un microcontrôleur simple). L'instruction est : \texttt{ADD \#3} (Ajoute la valeur immédiate 3 à l'accumulateur).
Supposons que l'accumulateur (\texttt{ACC}) contient déjà \texttt{5}. L'instruction se trouve à l'adresse \texttt{@100}.

\begin{exemplebox}[Déroulement pas à pas : \texttt{ADD \#3} (ACC = 5)]
\textbf{État initial :} \texttt{PC = @100}, \texttt{ACC = 5}, Mémoire[@100] = \texttt{ADD \#3} (code binaire : ex. \texttt{1100 0011} où \texttt{1100}=opcode ADD immédiat, \texttt{0011}=3).

\begin{enumerate}
    \item \textbf{Fetch} :
    \begin{itemize}
        \item L'UC place \texttt{PC (@100)} sur le \textbf{bus d'adresses} (via registre \texttt{MAR}).
        \item Signal \texttt{READ} sur bus de contrôle.
        \item La mémoire place l'instruction \texttt{ADD \#3} sur le \textbf{bus de données}.
        \item L'UC la charge dans le \textbf{Registre d'Instruction (RI)}.
        \item \textbf{PC $\leftarrow$ PC + 1} (devient \texttt{@101}).
    \end{itemize}
    \item \textbf{Decode} :
    \begin{itemize}
        \item L'UC extrait l'opcode \texttt{1100} $\to$ « Addition immédiate ».
        \item L'UC extrait l'opérande \texttt{0011} (valeur 3).
        \item L'UC configure l'UAL : entrée A = \texttt{ACC}, entrée B = valeur immédiate 3, opération = ADD. Signale que le résultat ira dans \texttt{ACC}.
    \end{itemize}
    \item \textbf{Execute} :
    \begin{itemize}
        \item L'UAL calcule : \texttt{5 + 3 = 8}.
        \item L'UAL met à jour les \textbf{drapeaux} : \texttt{Z=0} (résultat non nul), \texttt{C=0} (pas de retenue sur 8 bits), \texttt{N=0} (positif), \texttt{V=0} (pas de débordement signé).
    \end{itemize}
    \item \textbf{Writeback} :
    \begin{itemize}
        \item Le résultat \texttt{8} est écrit dans le registre \texttt{ACC}.
        \item \texttt{PC} vaut déjà \texttt{@101} (prochaine instruction).
    \end{itemize}
\end{enumerate}
\textbf{État final :} \texttt{ACC = 8}, \texttt{PC = @101}, Drapeaux mis à jour. Le cycle repart pour l'instruction à \texttt{@101}.
\end{exemplebox}

\begin{exoresolu}[Analyse d'une instruction de chargement : \texttt{LOAD @200}]
\textbf{Énoncé :} L'instruction \texttt{LOAD @200} (charger le contenu de l'adresse mémoire @200 dans l'accumulateur) se trouve à l'adresse \texttt{@150}. \texttt{PC=@150}. La mémoire @200 contient la valeur \texttt{42}. Décrire les 4 étapes du cycle.

\textbf{Solution détaillée :}
\begin{enumerate}
    \item \textbf{Fetch} : \texttt{PC (@150)} $\to$ Bus Adresses $\to$ \texttt{MAR}. Lecture mémoire. Bus Données $\to$ \texttt{RI} (contient \texttt{LOAD @200}). \texttt{PC $\leftarrow$ @151}.
    \item \textbf{Decode} : UC décode opcode \texttt{LOAD} (adressage direct). Identifie l'adresse effective \texttt{@200} (partie adresse de l'instruction). Prépare un cycle mémoire \texttt{READ} pour l'étape Execute. Sélectionne \texttt{ACC} comme destination.
    \item \textbf{Execute} (souvent en 2 sous-cycles mémoire pour chargement) :
        \begin{itemize}
            \item \texttt{@200} $\to$ Bus Adresses (via \texttt{MAR}). Signal \texttt{READ}.
            \item Mémoire place \texttt{42} sur Bus Données $\to$ \texttt{MDR} (Registre Données Mémoire).
            \item Transfert \texttt{MDR} $\to$ Bus Interne $\to$ Entrée UAL (passage direct) $\to$ Sortie UAL.
        \end{itemize}
    \item \textbf{Writeback} : Sortie UAL (\texttt{42}) $\to$ \texttt{ACC}. Drapeaux mis à jour (\texttt{Z=0}, \texttt{N=0}...). \texttt{PC=@151}.
\end{enumerate}
\textbf{Remarque :} Les instructions d'accès mémoire (LOAD, STORE) prennent souvent plus de temps (plusieurs cycles d'horloge) que les instructions registre-à-registre (ADD, AND) car elles nécessitent des allers-retours supplémentaires sur le bus système.
\end{exoresolu}

\courssub{La cadence : des milliards de cycles par seconde}

Tu te demandes peut-être : « Tout ça pour une simple addition ? » Oui, mais le processeur ne fait \textbf{qu'une seule chose à la fois} (dans un cœur simple) : il enchaîne ces 4 étapes à une vitesse folle.

\begin{savaistu}[La vitesse de la lumière... électrique]
Un processeur moderne (ton smartphone, ton PC) tourne à une fréquence de \textbf{2 à 4 GHz} (GigaHertz).
\cle{1 GHz = 1 milliard de cycles par seconde}.
À 3 GHz, le processeur effectue \textbf{3 milliards de cycles d'horloge par seconde}.
Comme une instruction simple prend environ 1 à 4 cycles d'horloge (grâce au \textit{pipelining} que tu approfondiras plus tard), ton téléphone peut exécuter \textbf{presque 1 milliard d'instructions par seconde \textit{par cœur}}.
Quand tu joues à un jeu 3D, le GPU (processeur graphique) en fait encore bien plus en parallèle (des milliers de cœurs) pour calculer la position, la couleur et l'ombre de chaque pixel, 60 fois par seconde. Tout repose sur ce petit cycle \textit{Fetch-Decode-Execute} répété à l'infini !
\end{savaistu}

\begin{aretenir}[Ce qu'il faut retenir du cycle d'exécution]
\begin{itemize}
    \item Le programme est une suite d'\textbf{instructions machine} en mémoire.
    \item Le \textbf{Compteur Ordinal (PC)} pointe toujours sur l'instruction \textbf{suivante} à exécuter.
    \item Le cycle \textbf{Fetch $\to$ Decode $\to$ Execute $\to$ Writeback} est la boucle fondamentale du CPU.
    \item L'\textbf{Unité de Contrôle (UC)} dirige le ballet : bus, registres, UAL.
    \item L'\textbf{UAL} fait le « sale boulot » (calculs, logique).
    \item La fréquence d'horloge (Hz) donne le rythme ; les architectures modernes (pipeline, superscalaire, multicœur) permettent de faire \emph{chevaucher} plusieurs cycles pour aller encore plus vite.
\end{itemize}
\end{aretenir}

\courssub{Exercices d'application}

\begin{exercice}[Cycle d'exécution et registres]
On considère un processeur simplifié avec les registres : \texttt{PC} (Compteur Ordinal), \texttt{MAR} (Registre d'Adresse Mémoire), \texttt{MDR} (Registre de Données Mémoire), \texttt{RI} (Registre d'Instruction), \texttt{ACC} (Accumulateur), \texttt{UAL}.
L'instruction \texttt{STORE @300} (sauvegarder \texttt{ACC} en @300) est à l'adresse \texttt{@500}. \texttt{PC=@500}, \texttt{ACC=12}.
\begin{enumerate}
    \item Complète le tableau de suivi des registres pour les 4 étapes (Fetch, Decode, Execute, Writeback). Indique le contenu de \texttt{PC}, \texttt{MAR}, \texttt{MDR}, \texttt{RI}, \texttt{ACC} après chaque étape.
    \item Sur quel bus circule l'adresse \texttt{@300} ? À quelle étape ?
    \item Quelle est la valeur de \texttt{PC} à la fin du cycle ?
\end{enumerate}
\end{exercice}

\begin{corrige}[Exercice n°1]
\begin{enumerate}
    \item \textbf{Tableau de suivi :}
    \begin{center}
    \begin{tabularx}{\linewidth}{|l|>{\centering\arraybackslash}X|>{\centering\arraybackslash}X|>{\centering\arraybackslash}X|>{\centering\arraybackslash}X|>{\centering\arraybackslash}X|}
    \hline
    \rowcolor{popblueL}
    \textbf{Étape} & \textbf{PC} & \textbf{MAR} & \textbf{MDR / RI} & \textbf{ACC} & \textbf{Action clé} \\
    \hline
    Initial & @500 & -- & -- & 12 & -- \\
    \hline
    Fetch & @501 & @500 & \texttt{STORE @300} (RI) & 12 & Lecture instr. @500, PC+1 \\
    \hline
    Decode & @501 & @500 & \texttt{STORE @300} & 12 & UC décode opcode STORE, adresse @300 \\
    \hline
    Execute & @501 & @300 & 12 (MDR $\leftarrow$ ACC) & 12 & Adresse @300 dans MAR, Donnée (12) dans MDR, Signal WRITE \\
    \hline
    Writeback & @501 & @300 & 12 & 12 & Mémoire[@300] $\leftarrow$ 12. Fin. \\
    \hline
    \end{tabularx}
    \end{center}
    \item L'adresse \texttt{@300} circule sur le \textbf{bus d'adresses} lors de l'étape \textbf{Execute} (chargement dans \texttt{MAR} pour l'écriture mémoire).
    \item \texttt{PC = @501} (inchangé depuis la fin du Fetch, car STORE n'est pas une instruction de saut).
\end{enumerate}
\end{corrige}

\begin{exercice}[Ordre des étapes et rôle de l'UAL]
Parmi les affirmations suivantes, coche les vraies et corrige les fausses en justifiant.
\begin{enumerate}
    \item L'étape Decode précède toujours l'étape Fetch.
    \item L'UAL intervient uniquement lors de l'étape Execute.
    \item Le Compteur Ordinal (PC) est incrémenté lors de l'étape Writeback.
    \item Une instruction de saut (JMP) modifie le PC lors de l'étape Execute.
    \item L'instruction \texttt{ADD R1, R2} (R1 $\leftarrow$ R1 + R2) ne nécessite pas d'accès à la mémoire principale lors de l'étape Execute.
\end{enumerate}
\end{exercice}

\begin{corrige}[Exercice n°2]
\begin{enumerate}
    \item \textbf{Faux.} L'ordre est immuable : \textbf{Fetch $\to$ Decode $\to$ Execute $\to$ Writeback}. On ne peut pas décoder avant d'avoir cherché l'instruction.
    \item \textbf{Vrai.} L'UAL est l'organe d'exécution. Lors du Fetch, Decode et Writeback, ce sont l'UC, les bus et les registres qui travaillent. L'opération logique/arithmétique de l'instruction se fait à l'étape Execute.
    \item \textbf{Faux.} Le PC est incrémenté (ou mis à jour) dès la fin de l'étape \textbf{Fetch} (pour pointer sur l'instruction suivante) ou modifié à l'étape Execute (saut). L'étape Writeback écrit le résultat de l'opération (dans un registre ou mémoire), pas le PC.
    \item \textbf{Vrai.} L'instruction \texttt{JMP @addr} charge la nouvelle adresse \texttt{@addr} dans le \texttt{PC} pendant l'étape Execute (via l'UAL ou un chemin direct), changeant ainsi le flux du programme.
    \item \textbf{Vrai.} C'est une instruction \textbf{registre-à-registre}. Les opérandes sont dans les registres internes (R1, R2), le résultat va dans R1. Aucun bus système (mémoire RAM) n'est sollicité pendant Execute. C'est pour cela qu'elles sont les plus rapides.
\end{enumerate}
\end{corrige}

\coursec{Les architectures RISC et CISC}

Dans la section précédente, nous avons suivi le cycle d'exécution d'une instruction : le processeur va chercher l'instruction en mémoire, la décode, l'exécute, puis recommence. Mais une question reste en suspens : \emph{quelles instructions un processeur sait-il exécuter ?} Tous les processeurs ne se ressemblent pas : celui d'un ordinateur de bureau et celui d'un smartphone ne sont pas construits de la même façon. Pour comprendre cette différence, nous devons découvrir la notion de \cle{jeu d'instructions}, puis les deux grandes familles d'architectures : \cle{CISC} et \cle{RISC}.

\courssub{Le jeu d'instructions : le vocabulaire du processeur}

\begin{definition}[Jeu d'instructions]
Le \cle{jeu d'instructions} (en anglais \emph{instruction set}) est l'ensemble de toutes les instructions élémentaires qu'un processeur sait reconnaître et exécuter : addition, soustraction, comparaison, lecture ou écriture en mémoire, saut vers une autre instruction, etc.
\end{definition}

On peut comparer le jeu d'instructions au \textbf{vocabulaire} d'une langue. Deux personnes qui parlent la même langue se comprennent ; de même, un programme écrit pour un certain jeu d'instructions ne fonctionne que sur les processeurs qui connaissent ce jeu. C'est pourquoi une application conçue pour un ordinateur ne s'installe pas directement sur un téléphone : les deux appareils ne parlent pas le même « langage machine ».

Or, les concepteurs de processeurs ont fait deux choix opposés pour construire ce vocabulaire :

\begin{itemize}
  \item soit un vocabulaire \textbf{riche et complexe}, avec beaucoup de mots très expressifs : c'est l'approche \cle{CISC} ;
  \item soit un vocabulaire \textbf{réduit et simple}, avec peu de mots mais des mots très rapides à prononcer : c'est l'approche \cle{RISC}.
\end{itemize}

\courssub{L'architecture CISC : des instructions complexes}

\begin{definition}[Architecture CISC]
Une architecture \cle{CISC} (\emph{Complex Instruction Set Computer}, « ordinateur à jeu d'instructions complexe ») est une architecture dans laquelle le processeur possède un \textbf{grand nombre d'instructions}, dont certaines sont \textbf{complexes} : une seule instruction peut réaliser plusieurs opérations à la fois (par exemple lire une donnée en mémoire, l'additionner à une autre et ranger le résultat).
\end{definition}

L'idée historique du CISC (années 1970-1980) était de faciliter le travail des programmeurs : plus le processeur en fait en une seule instruction, plus les programmes sont \textbf{courts}. Par exemple, une instruction CISC peut signifier : « va chercher le nombre à l'adresse 1000, ajoute-le au nombre à l'adresse 2000, et range le résultat à l'adresse 3000 » — tout cela en \emph{une seule} instruction.

Mais cette richesse a un prix : pour comprendre et exécuter des instructions aussi complexes, le processeur doit être lui-même très \textbf{complexe}, avec de nombreux circuits électroniques. Il devient plus difficile à fabriquer, plus gourmand en énergie, et il chauffe davantage. De plus, une instruction complexe met \textbf{plusieurs cycles d'horloge} à s'exécuter.

L'exemple le plus célèbre d'architecture CISC est la famille \cle{x86} d'\textbf{Intel} (et AMD), qui équipe la grande majorité des ordinateurs de bureau et des ordinateurs portables, y compris ceux des salles multimédias des collèges et lycées du Cameroun.

\courssub{L'architecture RISC : des instructions simples et rapides}

\begin{definition}[Architecture RISC]
Une architecture \cle{RISC} (\emph{Reduced Instruction Set Computer}, « ordinateur à jeu d'instructions réduit ») est une architecture dans laquelle le processeur ne possède qu'un \textbf{petit nombre d'instructions simples}, chacune réalisant une seule opération élémentaire, et conçues pour s'exécuter \textbf{très rapidement}, généralement en un seul cycle d'horloge.
\end{definition}

L'approche RISC (années 1980) part d'un constat : les programmeurs (et surtout les logiciels qui traduisent les programmes, les \emph{compilateurs}) n'utilisent presque jamais les instructions très complexes. Autant donc ne garder que des instructions \textbf{simples}, mais les exécuter à toute vitesse. Une tâche complexe est alors découpée en plusieurs instructions simples.

\begin{propriete}[Rapidité des instructions RISC (admise)]
Dans une architecture RISC, les instructions sont simples et régulières : elles s'exécutent généralement en \textbf{un seul cycle d'horloge}. Le processeur est donc plus \textbf{simple à construire}, plus \textbf{rapide} et plus \textbf{économe en énergie} qu'un processeur CISC de génération comparable.
\end{propriete}

La figure suivante montre comment une même tâche (additionner deux nombres rangés en mémoire) est réalisée par une seule instruction CISC complexe, ou par une suite de quatre instructions RISC simples.

\begin{popfigure}
\begin{center}
\begin{tikzpicture}[font=\small, >=Stealth]
  % Bloc CISC
  \node[draw=poporange, fill=poporangeL, rounded corners, minimum width=7.2cm, minimum height=1.3cm, align=center] (cisc) at (0,0)
    {\textbf{1 instruction CISC}\\ \texttt{ADD [3000], [1000], [2000]}\\ (lit, additionne et range en une fois)};
  % Blocs RISC
  \node[draw=popblue, fill=popblueL, rounded corners, minimum width=7.2cm, minimum height=0.8cm, align=center] (r1) at (0,-2.2) {\texttt{LOAD R1, [1000]} \quad charger le 1\ier{} nombre};
  \node[draw=popblue, fill=popblueL, rounded corners, minimum width=7.2cm, minimum height=0.8cm, align=center] (r2) at (0,-3.2) {\texttt{LOAD R2, [2000]} \quad charger le 2\ieme{} nombre};
  \node[draw=popblue, fill=popblueL, rounded corners, minimum width=7.2cm, minimum height=0.8cm, align=center] (r3) at (0,-4.2) {\texttt{ADD R3, R1, R2} \quad additionner};
  \node[draw=popblue, fill=popblueL, rounded corners, minimum width=7.2cm, minimum height=0.8cm, align=center] (r4) at (0,-5.2) {\texttt{STORE [3000], R3} \quad ranger le résultat};
  \node[above=0.1cm of cisc, text=poporange, font=\bfseries] {Approche CISC};
  \node[text=popblue, font=\bfseries] at (0,-1.5) {Approche RISC : 4 instructions simples, chacune en 1 cycle};
  \draw[->, thick, popmuted] (cisc.south) -- (r1.north);
\end{tikzpicture}
\end{center}
\legende{Une instruction CISC complexe équivaut à plusieurs instructions RISC simples. Chaque instruction RISC est exécutée en un cycle d'horloge : le processeur reste simple et rapide.}
\end{popfigure}

L'exemple le plus connu d'architecture RISC est \cle{ARM}, qui équipe la quasi-totalité des \textbf{smartphones} et \textbf{tablettes} du monde, ainsi que de nombreux objets connectés (montres, téléviseurs, routeurs).

\courssub{Comparaison synthétique RISC / CISC}

\begin{popfigure}
\begin{center}
\begin{tikzpicture}[font=\small]
  % En-têtes
  \fill[popblue] (-4.2,0) rectangle (4.2,0.7);
  \node[white, font=\bfseries] at (-2.8,0.35) {Critère};
  \node[white, font=\bfseries] at (0,0.35) {CISC};
  \node[white, font=\bfseries] at (2.8,0.35) {RISC};
  % Ligne 1
  \fill[popcream] (-4.2,-0.9) rectangle (4.2,0);
  \node[anchor=west] at (-4.1,-0.45) {\textbf{Instructions}};
  \node[align=center] at (0,-0.45) {nombreuses et\\ complexes};
  \node[align=center] at (2.8,-0.45) {peu nombreuses\\ et simples};
  % Ligne 2
  \fill[popblueL!40] (-4.2,-1.8) rectangle (4.2,-0.9);
  \node[anchor=west] at (-4.1,-1.35) {\textbf{Vitesse}};
  \node[align=center] at (0,-1.35) {instruction complexe :\\ plusieurs cycles};
  \node[align=center] at (2.8,-1.35) {instruction simple :\\ 1 cycle en général};
  % Ligne 3
  \fill[popcream] (-4.2,-2.7) rectangle (4.2,-1.8);
  \node[anchor=west] at (-4.1,-2.25) {\textbf{Énergie}};
  \node[align=center] at (0,-2.25) {forte consommation,\\ chauffe beaucoup};
  \node[align=center] at (2.8,-2.25) {faible consommation,\\ chauffe peu};
  % Ligne 4
  \fill[popblueL!40] (-4.2,-3.6) rectangle (4.2,-2.7);
  \node[anchor=west] at (-4.1,-3.15) {\textbf{Usages}};
  \node[align=center] at (0,-3.15) {ordinateurs de bureau,\\ portables (Intel x86)};
  \node[align=center] at (2.8,-3.15) {smartphones, tablettes,\\ objets connectés (ARM)};
  % Contours
  \draw[popline] (-4.2,0.7) rectangle (4.2,-3.6);
  \draw[popline] (-1.4,0.7) -- (-1.4,-3.6);
  \draw[popline] (1.4,0.7) -- (1.4,-3.6);
  \draw[popline] (-4.2,0) -- (4.2,0);
  \draw[popline] (-4.2,-0.9) -- (4.2,-0.9);
  \draw[popline] (-4.2,-1.8) -- (4.2,-1.8);
  \draw[popline] (-4.2,-2.7) -- (4.2,-2.7);
\end{tikzpicture}
\end{center}
\legende{Tableau comparatif des architectures CISC et RISC selon quatre critères : instructions, vitesse, consommation d'énergie et domaines d'application.}
\end{popfigure}

\begin{aretenir}[L'essentiel]
\begin{itemize}
  \item \cle{CISC} : beaucoup d'instructions complexes $\Rightarrow$ \textbf{programmes courts} mais \textbf{processeur complexe}, qui consomme et chauffe beaucoup (ex. Intel x86, ordinateurs de bureau).
  \item \cle{RISC} : peu d'instructions simples exécutées en un cycle $\Rightarrow$ \textbf{processeur simple, rapide et économe en énergie} (ex. ARM, smartphones et tablettes).
\end{itemize}
\end{aretenir}

\begin{exemplebox}[Pourquoi la batterie d'un smartphone tient-elle si longtemps ?]
Un smartphone possède une batterie beaucoup plus petite qu'un ordinateur portable, pourtant il tient souvent une journée entière. L'explication est en grande partie architecturale : son processeur \textbf{ARM} est de type \cle{RISC}. Ses instructions simples nécessitent moins de circuits électroniques, donc moins d'énergie électrique, et le processeur chauffe peu — c'est pour cela qu'un smartphone n'a pas besoin de ventilateur. À l'inverse, le processeur \cle{CISC} d'un ordinateur portable consomme davantage et nécessite un ventilateur pour évacuer la chaleur, ce qui vide la batterie plus vite.
\end{exemplebox}

\begin{attention}[« Reduced » ne veut pas dire « moins puissant » !]
Erreur fréquente : croire qu'un processeur \cle{RISC} (« réduit ») est moins puissant qu'un processeur CISC. C'est faux ! « Réduit » concerne le \textbf{nombre d'instructions du jeu}, pas la puissance de la machine. En exécutant chaque instruction simple en un seul cycle, un processeur RISC est souvent \textbf{plus efficace} qu'un CISC. Les processeurs ARM des smartphones actuels réalisent des tâches considérables (vidéo, jeux, photos) avec une remarquable efficacité.
\end{attention}

\begin{savaistu}[ARM partout, même dans ta poche]
Plus de \textbf{90 \%} des smartphones du monde, y compris ceux utilisés chaque jour au Cameroun pour communiquer, payer par mobile money ou suivre des cours en ligne, fonctionnent avec des processeurs \cle{ARM} d'architecture \cle{RISC}. L'entreprise ARM ne fabrique d'ailleurs aucune puce : elle conçoit des architectures que les fabricants (Samsung, MediaTek, Qualcomm...) adaptent à leurs appareils.
\end{savaistu}

\begin{exoresolu}[Identifier une architecture]
Énoncé. Le collège bilingue de Nkol-Eton veut acheter du matériel. On hésite entre : (a) des ordinateurs de bureau à processeur Intel x86 pour la salle multimédia ; (b) des tablettes à processeur ARM pour le club informatique. Pour chaque choix, indique l'architecture du processeur (RISC ou CISC) et justifie un avantage de ce choix.
\tcblower
Solution détaillée.
\begin{enumerate}
  \item \textbf{Ordinateurs de bureau Intel x86} : architecture \cle{CISC}. Avantage : le jeu d'instructions riche permet de faire fonctionner les logiciels classiques de bureautique (traitement de texte, tableur) ; les programmes sont courts et l'offre logicielle est immense. La consommation électrique élevée n'est pas gênante car l'appareil est branché sur secteur.
  \item \textbf{Tablettes ARM} : architecture \cle{RISC}. Avantage : les instructions simples s'exécutent en un cycle et consomment peu d'énergie, ce qui donne une \textbf{grande autonomie de batterie} et évite la chauffe — idéal pour des élèves qui se déplacent avec l'appareil toute la journée.
\end{enumerate}
Conclusion : le choix de l'architecture dépend de l'usage : CISC pour la puissance logicielle sur secteur, RISC pour la mobilité et l'économie d'énergie.
\end{exoresolu}

\begin{exercice}[À toi de jouer]
\begin{enumerate}
  \item Définis le jeu d'instructions d'un processeur.
  \item Donne deux différences entre les architectures RISC et CISC.
  \item Ton petit frère affirme : « mon téléphone est moins puissant qu'un ordinateur parce que son processeur est réduit ». Corrige son erreur en deux phrases.
\end{enumerate}
\end{exercice}

\begin{corrige}[Exercice]
\begin{enumerate}
  \item Le jeu d'instructions est l'ensemble de toutes les instructions élémentaires qu'un processeur sait reconnaître et exécuter (addition, lecture en mémoire, comparaison, etc.).
  \item Par exemple : le CISC possède de nombreuses instructions complexes alors que le RISC n'en possède que peu, et simples ; le CISC consomme beaucoup d'énergie alors que le RISC est économe (on pouvait aussi citer : plusieurs cycles par instruction contre un seul, ordinateurs de bureau contre smartphones).
  \item « Réduit » signifie que le \emph{jeu d'instructions} est réduit, pas la puissance : le processeur RISC exécute chaque instruction simple en un seul cycle, il est donc très efficace. Les smartphones réalisent des tâches exigeantes (vidéo, jeux) grâce à cette efficacité.
\end{enumerate}
\end{corrige}

\begin{savaistu}[Hors-programme : et demain ?]
Pour ta culture générale (ceci dépasse le programme de 3ème) : la frontière entre RISC et CISC s'estompe. Les processeurs Intel modernes traduisent en interne leurs instructions complexes en micro-opérations simples de type RISC, et les puces \textbf{Apple M} des ordinateurs Mac récents sont d'architecture ARM (RISC) tout en rivalisant avec les meilleurs processeurs de bureau. La tendance actuelle \textbf{mélange les deux approches} pour combiner leurs avantages.
\end{savaistu}

\coursec{Appliquer et justifier : démarche argumentée}

Après avoir exploré les architectures \cle{RISC} et \cle{CISC}, nous possédons maintenant toutes les clés techniques pour analyser un besoin réel et proposer un choix d'équipement informatique \emph{argumenté}. Au BEPC, comme dans la vie professionnelle, il ne suffit pas de donner une réponse : il faut la \textbf{justifier} avec le vocabulaire précis du cours (fréquence, nombre de cœurs, taille de la RAM, type de bus, architecture du processeur). Cette section t'apprend la démarche rigoureuse attendue à l'examen.

\courssub{La méthode en trois étapes : Besoin → Choix → Argumentation}

\begin{methode}[Rédiger une réponse justifiée en informatique]
\textbf{Étape 1 : Identifier le besoin réel (Analyse de la situation)} \\
Lis attentivement l'énoncé. Repère les \textbf{usages} (bureautique, montage vidéo, jeu, navigation web, programmation), les \textbf{contraintes} (budget, mobilité/autonomie, environnement réseau) et le \textbf{profil utilisateur} (secrétaire, élève, développeur, graphiste). Note les mots-clés : « fluidité », « multitâche », « stockage massif », « autonomie 8h ».

\textbf{Étape 2 : Sélectionner les caractéristiques techniques pertinentes (Transposition)} \\
Fais le lien entre l'usage et l'architecture vue en cours :
\begin{itemize}
    \item \textbf{Bureautique / Navigation / Cours} $\rightarrow$ Processeur 2 à 4 cœurs (architecture CISC x86-64 ou ARMv8), RAM 4 à 8 Go, SSD 256 Go. Priorité : coût maîtrisé, silence, fiabilité.
    \item \textbf{Multitâche lourd / VM / Compilation} $\rightarrow$ Processeur 6 à 8 cœurs (CISC haute performance), RAM 16 Go minimum, SSD NVMe 512 Go+. Priorité : fréquence, cache L3, bande passante bus mémoire.
    \item \textbf{Mobilité / Autonomie / Smartphone} $\rightarrow$ Architecture \cle{RISC} (ARM), gravure fine (4 nm, 3 nm), gestion d'énergie \texttt{big.LITTLE}, batterie $\ge$ 5000 mAh. Priorité : performance par Watt.
\end{itemize}

\textbf{Étape 3 : Rédiger l'argumentation structurée (Justification)} \\
Structure ta réponse par \textbf{phrases complètes} reliant \emph{besoin} $\rightarrow$ \emph{composant} $\rightarrow$ \emph{rôle technique}.
\begin{enumerate}
    \item Cite le composant choisi (ex : « Un processeur \textbf{quad-core} à \textbf{2,5 GHz} »).
    \item Explique \textbf{pourquoi} au regard de l'architecture (ex : « Ses 4 cœurs permettent le \textbf{traitement parallèle} des tâches bureautiques et du navigateur sans gel de l'interface, grâce à l'ordonnancement du système d'exploitation sur l'\textbf{unité de contrôle} »).
    \item Mentionne l'impact sur le \textbf{bus système} ou la \textbf{mémoire} (ex : « La \textbf{RAM DDR4 8 Go} en \textbf{double canal} assure une bande passante suffisante pour charger l'OS et les applications sans \emph{swapping} disque »).
    \item Conclus par l'adéquation globale (ex : « Ce choix garantit la fluidité attendue pour un coût maîtrisé, adapté au budget d'une collectivité »).
\end{enumerate}
\end{methode}

\courssub{Critères techniques de choix d'un équipement}

\begin{propriete}[Critères de choix pertinents pour le BEPC]
Pour justifier un choix, tu dois mobiliser au minimum trois des critères suivants, en utilisant le vocabulaire exact de la leçon :
\begin{center}
\begin{tabularx}{\linewidth}{|l|X|X|}
\hline
\rowcolor{popblueL}
\textbf{Critère} & \textbf{Notions du cours à mobiliser} & \textbf{Exemple de formulation attendue} \\
\hline
\textbf{Usage principal} & Type d'applications (lourdes/légères), monoposte/réseau & « Pour de la bureautique (traitement de texte, tableur), la charge CPU est faible et séquentielle. » \\
\hline
\textbf{Processeur (CPU)} & Architecture (RISC/CISC), nombre de cœurs, fréquence, cache L1/L2/L3, TDP & « Un CPU \textbf{CISC x86-64} 4 cœurs / 8 threads à 3,0 GHz avec 8 Mo de cache L3 assure le multitâche bureautique. » \\
\hline
\textbf{Mémoire vive (RAM)} & Capacité (Go), type (DDR4/DDR5), fréquence, double canal, volatilité & « 8 Go \textbf{DDR4-3200} en double canal évitent le \emph{swapping} sur le disque, garantissant la réactivité. » \\
\hline
\textbf{Stockage} & SSD (NVMe/SATA) vs HDD, capacité, débit, non-volatilité & « Un \textbf{SSD NVMe 256 Go} offre des temps d'accès $\approx$ 0,1 ms contre 10 ms pour un HDD : démarrage OS < 15 s. » \\
\hline
\textbf{Autonomie (Mobile)} & Architecture RISC (ARM), gravure (nm), gestion big.LITTLE, capacité batterie (mAh/Wh) & « Le \textbf{SoC ARMv8} gravé en 4 nm avec cœurs \emph{efficiency} gère l'affichage et la 4G en faible consommation. » \\
\hline
\textbf{Coût / Durabilité} & Rapport performance/prix, garantie, réparabilité, consommation électrique (TDP) & « Un poste \textbf{fanless} (sans ventilateur) à base de CISC basse consommation (TDP 15 W) réduit la facture électrique et le bruit. » \\
\hline
\end{tabularx}
\end{center}
\end{propriete}

\courssub{Étude de cas 1 : Choisir un ordinateur pour un secrétariat de mairie (bureautique)}

\begin{exoresolu}[Choix d'un poste de travail bureautique pour un secrétariat]
\textbf{Situation :} La mairie de Dschang souhaite équiper le secrétariat (5 postes). Usage : traitement de texte (LibreOffice Writer), tableur (Calc), messagerie, navigateur web (démarches administratives en ligne), impression. Budget serré, maintenance interne faible, environnement poussiéreux. Durée de vie visée : 5 ans.

\textbf{Question :} Propose une configuration type justifiée pour un poste.

\tcblower

\textbf{Solution détaillée et justifiée :}

\textbf{1. Processeur (CPU) : Architecture CISC x86-64, 4 cœurs / 8 threads (SMT), fréquence de base 3,3 GHz (Turbo 4,3 GHz), Cache L3 12 Mo, TDP 60 W (ex: gamme Intel Core i3 / AMD Ryzen 3 récents).}
\begin{itemize}
    \item \textbf{Justification :} La bureautique moderne (onglets multiples, JavaScript lourd dans le navigateur) bénéficie du \textbf{SMT (Simultaneous Multithreading, Hyper-Threading chez Intel)} pour maintenir l'interface réactive. L'architecture \cle{CISC} x86-64 garantit la compatibilité totale avec Windows 10/11 et LibreOffice. Le TDP de 60 W permet un refroidissement passif ou semi-passif (ventilateur lent), crucial en milieu poussiéreux (moins d'aspiration de particules).
\end{itemize}

\textbf{2. Mémoire vive (RAM) : 8 Go DDR4-3200 (1 × 8 Go), extensible à 32 Go (2 slots).}
\begin{itemize}
    \item \textbf{Justification :} 8 Go est le \textbf{minimum confortable} aujourd'hui pour Windows 11 + navigateur (5-10 onglets) + suite bureautique sans déclencher le \emph{swapping} (échange RAM/Disque) qui use le SSD et gèle le système. Le choix d'une seule barrette laisse un slot libre pour une mise à niveau future (évolutivité 5 ans) sans jeter l'existant.
\end{itemize}

\textbf{3. Stockage : SSD M.2 NVMe 256 Go (bus PCIe 3.0 x4, débit $\approx$ 3500 Mo/s).}
\begin{itemize}
    \item \textbf{Justification :} Un \textbf{SSD NVMe} réduit le démarrage à $\approx$ 12 s et l'ouverture des applications à $\approx$ 1-2 s. L'absence de pièces mobiles (contrairement au HDD) assure la \textbf{robustesse} aux chocs et vibrations (transport, nettoyage local) et le silence total. 256 Go suffit pour l'OS, les logiciels et les documents (quelques Go/an), les archives lourdes pouvant aller sur un serveur partagé ou clé USB.
\end{itemize}

\textbf{4. Boîtier / Alimentation : Format SFF (Small Form Factor) ou Tour mini, Alim 80+ Bronze 300 W.}
\begin{itemize}
    \item \textbf{Justification :} Gain de place sur le bureau (mairie exiguë). Alimentation certifiée \textbf{80+ Bronze} : rendement $\ge$ 85\% $\rightarrow$ économie d'électricité sur 5 ans (5 postes $\times$ 8 h/j $\times$ 220 j/an).
\end{itemize}

\textbf{5. Écran : 24 pouces IPS Full HD (1920$\times$1080), anti-reflet, faible lumière bleue.}
\begin{itemize}
    \item \textbf{Justification :} Confort visuel pour 7 h/j de travail (santé au travail). IPS = angles de vision larges (partage écran avec usager au guichet).
\end{itemize}

\textbf{Synthèse de l'argumentation (modèle de réponse d'examen) :}
\emph{« Pour un usage bureautique intensif en milieu administratif, je recommande un poste à architecture \textbf{CISC x86-64} 4 cœurs / 8 threads (SMT). Cette architecture à jeu d'instructions complexe gère nativement les logiciels propriétaires de l'administration. Les 8 Go de \textbf{RAM DDR4} (évolutif) évitent le \emph{swapping} sur le \textbf{SSD NVMe} 256 Go, dont l'accès aléatoire rapide (IOPS élevés) garantit la fluidité du multitâche. Le format compact et l'alimentation efficace assurent la durabilité et la sobriété énergétique demandées. »}
\end{exoresolu}

\courssub{Étude de cas 2 : Choisir un smartphone pour un étudiant (autonomie, prix)}

\begin{exoresolu}[Choix d'un smartphone étudiant : autonomie et budget]
\textbf{Situation :} Étudiant en 3ème à Yaoundé. Besoins : WhatsApp, Facebook, recherches web, photos occasionnelles, GPS, révision de cours (PDF, vidéos YouTube). Budget strict : 100 000 - 150 000 FCFA. Exigence : tenir une journée complète (7 h - 22 h) sans recharge (coupures électriques fréquentes). Double SIM indispensable.

\textbf{Question :} Propose un choix justifié en t'appuyant sur l'architecture matérielle.

\tcblower

\textbf{Solution détaillée et justifiée :}

\textbf{Choix visé :} Smartphone d'entrée/milieu de gamme récent — \textbf{SoC ARM (RISC) gravé en 6 nm (ou 12 nm), 8 cœurs big.LITTLE, RAM 6/8 Go LPDDR4X (extensible virtuellement), Stockage 128/256 Go UFS 2.2, Batterie 5000-6000 mAh, Charge 18-33 W.}

\textbf{1. Architecture Processeur (SoC) : ARMv8-A (RISC) — Ex: 6 nm, 2×Cortex-A76 2,2 GHz + 6×Cortex-A55 2,0 GHz (type Helio G99 / Dimensity 6000) ou 12 nm, 2×A75 + 6×A55 (type Unisoc T606).}
\begin{itemize}
    \item \textbf{Justification :} L'architecture \cle{RISC} (ARM) est \textbf{incontournable sur mobile} pour son rapport performance/Watt. La configuration \textbf{big.LITTLE} (2 cœurs \emph{Performance} + 6 cœurs \emph{Efficiency}) est idéale : les cœurs \emph{Efficiency} gèrent l'affichage, la 4G, la réception messages (WhatsApp en fond) à très basse consommation ; les cœurs \emph{Performance} s'activent seulement pour le chargement de pages web lourdes ou le défilement fluide de l'interface. La gravure \textbf{6 nm} divise la consommation par rapport aux anciennes 12 nm pour une performance équivalente $\rightarrow$ \textbf{gain direct sur l'autonomie}.
\end{itemize}

\textbf{2. Mémoire (RAM) : 6 Go LPDDR4X (physique) + 6 Go \emph{virtuel} (swap sur stockage UFS/eMMC).}
\begin{itemize}
    \item \textbf{Justification :} 6 Go physiques permettent de garder 5-6 apps en mémoire (WhatsApp, Messenger, Chrome, PDF Reader, Galerie) sans rechargement. L'extension virtuelle (fonction \texttt{RAM Boost}) utilise le stockage rapide comme débordement, évitant les \emph{kills} d'apps par l'OS (Android LMKD) quand la RAM physique sature. \textbf{Attention :} la RAM virtuelle est plus lente (bus stockage vs bus mémoire) et use le stockage (écritures NAND).
\end{itemize}

\textbf{3. Stockage : 128 Go UFS 2.2 (ou eMMC 5.1 sur les moins chers).}
\begin{itemize}
    \item \textbf{Justification :} \textbf{UFS 2.2} (Universal Flash Storage) offre un débit séquentiel $\approx$ 800-1200 Mo/s et surtout des \textbf{IOPS aléatoires} bien supérieurs à l'eMMC. Cela se ressent sur la fluidité de l'interface (animation, ouverture apps) et la vitesse de copie des cours (PDF, vidéos). 128 Go est le minimum pour l'OS (15-20 Go) + apps + photos/vidéos + cours hors-ligne.
\end{itemize}

\textbf{4. Batterie \& Gestion énergie : 5000 mAh (typique) / 19,2 Wh + Charge 18-33 W.}
\begin{itemize}
    \item \textbf{Justification :} La capacité \textbf{5000 mAh} est le standard « 1 jour garanti » pour un écran LCD 90 Hz 6,6-6,8 pouces. L'architecture RISC + gravure fine + gestion \texttt{Doze} d'Android (mise en veille profonde CPU/radios écran éteint) permet d'atteindre 8-10 h \emph{Screen-On Time} (SOT) en usage mixte. La charge 18-33 W (PD/QC) recharge 50\% en 30 min (pratique entre deux cours).
\end{itemize}

\textbf{5. Connectivité : Double SIM 4G (VoLTE) + Wi-Fi 5 (ac) + Bluetooth 5.x.}
\begin{itemize}
    \item \textbf{Justification :} Double SIM physique (nano) indispensable pour gérer forfait « Data pas cher » (MTN/Orange/Nexttel) et forfait « Appels famille ». VoLTE = appels HD sur réseau 4G sans bascule 2G/3G (économie batterie + qualité).
\end{itemize}

\textbf{Synthèse de l'argumentation (modèle de réponse d'examen) :}
\emph{« Pour un étudiant avec budget limité et besoin d'autonomie, je préconise un smartphone à \textbf{SoC ARMv8 (RISC)} gravé en \textbf{6 nm} avec architecture \textbf{big.LITTLE} (2 cœurs Perf + 6 cœurs Eff). Cette architecture RISC à jeu d'instructions réduit minimise la consommation au repos (cœurs Efficiency) tout en offrant la réactivité nécessaire (cœurs Performance). Les \textbf{6 Go de RAM LPDDR4X} (avec extension virtuelle) assurent le multitâche éducatif (PDF, navigateur, messagerie) sans rechargement d'apps. Le \textbf{stockage UFS 2.2 128 Go} garantit des IOPS élevés pour la fluidité système. Enfin, la \textbf{batterie 5000 mAh} couplée à la gravure 6 nm et à la gestion \texttt{Doze} d'Android assure l'autonomie « journée complète » requise malgré les coupures secteur. »}
\end{exoresolu}

\courssub{Relire et vérifier sa justification}

\begin{attention}[Erreur fréquente : l'absence de justification technique]
\textbf{Le piège :} Répondre « Je choisis un Core i5 avec 16 Go de RAM et 512 Go SSD parce que c'est puissant et rapide. »
\begin{itemize}
    \item \textbf{Pourquoi c'est insuffisant :} « Puissant » et « rapide » sont des adjectifs \textbf{subjectifs}. Ils ne démontrent pas que tu comprends \emph{comment} l'architecture produit cette performance.
    \item \textbf{Correction attendue :} « Je choisis un \textbf{processeur CISC x86-64 6 cœurs / 12 threads (SMT), cache L3 18 Mo (ex: gamme Core i5 / Ryzen 5 récents)}. Ses \textbf{6 cœurs physiques} permettent l'exécution \textbf{parallèle} des threads de l'OS, du navigateur (multi-processus) et de la compilation (make -j6) sans contention sur l'\textbf{unité de contrôle}. Les \textbf{16 Go DDR4-3200 double canal} offrent une \textbf{bande passante $\approx$ 50 Go/s} au \textbf{bus mémoire}, évitant le goulot d'étranglement lors du chargement de gros jeux de données en RAM. Le \textbf{SSD NVMe PCIe 4.0} (7000 Mo/s) réduit la latence d'accès aux fichiers projet (IOPS aléatoires 4K $\approx$ 600k) vs un SATA (80k). »
\end{itemize}
\textbf{Règle d'or :} Chaque affirmation technique (ex: « 6 cœurs ») doit être suivie de sa \textbf{conséquence architecturale} (ex: « parallélisme thread / unité de contrôle ») et de son \textbf{impact utilisateur} (ex: « compilation sans gel »).
\end{attention}

\begin{savaistu}[Au BEPC, le vocabulaire technique fait la différence]
Dans la grille d'évaluation officielle MINESEC pour l'épreuve pratique ou écrite d'Informatique 3ème, la compétence « \textbf{Rédiger et justifier une démarche} » est notée sur \textbf{4 à 6 points} sur 20.
\begin{itemize}
    \item \textbf{0-1 pt :} Réponse vague, termes courants (« rapide », « grand », « fort »).
    \item \textbf{2-3 pts :} Choix correct mais justification incomplète (ex: cite la RAM mais pas le bus, cite le CPU mais pas l'architecture RISC/CISC).
    \item \textbf{4-6 pts :} \textbf{Justification structurée} : Besoin identifié $\rightarrow$ Composant nommé avec specs exactes $\rightarrow$ Lien explicite avec une notion du cours (cycle instruction, bus, pipeline, cache, volatilité, gravure, TDP, big.LITTLE, swapping, IOPS).
\end{itemize}
\textbf{Astuce :} Apprends par cœur 3 « phrases-types » de justification (CPU, RAM, Stockage, Batterie) que tu adaptes au contexte. Ex: « La capacité de \textbf{X Go} de \textbf{RAM DDR4} en \textbf{double canal} assure une bande passante mémoire suffisante pour [usage] sans recourir au \emph{swapping} sur le disque, ce qui préserverait la réactivité et la durée de vie du SSD. »
\end{savaistu}

\courssub{Synthèse de la leçon : Carte mentale de l'architecture de l'ordinateur}

\begin{popfigure}
\begin{tikzpicture}[
    mindmap,
    concept color=popblue,
    level 1/.style={level distance=4.5cm, sibling angle=90},
    level 2/.style={level distance=3.2cm, sibling angle=45},
    level 3/.style={level distance=2.8cm, sibling angle=30},
    every node/.style={font=\small},
    grow cyclic,
    text=popdark,
    edge from parent/.style={draw=popline, thick}
]
\node[concept, minimum size=2.2cm, fill=popblue!90, text=white, align=center]{Architecture\\ Ordinateur}
    child[concept color=poporange] {
        node[concept, minimum size=1.8cm] {CPU (UC + UAL)}
        child[concept color=poporange!70] { node[concept] {Unité de Contrôle (UC)} child { node[concept, scale=0.7] {Séquencement} } child { node[concept, scale=0.7] {Décodage} } }
        child[concept color=poporange!70] { node[concept] {Unité Arithmétique \& Logique (UAL)} child { node[concept, scale=0.7] {Calculs entiers/flottants} } child { node[concept, scale=0.7] {Opérations logiques} } child { node[concept, scale=0.7] {Comparaisons} } }
        child[concept color=poporange!70] { node[concept] {Registres} child { node[concept, scale=0.7] {PC, IR, ACC, SP} } child { node[concept, scale=0.7] {Registres généraux} } }
        child[concept color=poporange!70] { node[concept] {Architectures} child { node[concept, scale=0.7] {CISC (x86-64)} } child { node[concept, scale=0.7] {RISC (ARM, RISC-V)} } }
    }
    child[concept color=popgreen] {
        node[concept, minimum size=1.8cm] {Mémoires}
        child[concept color=popgreen!70] { node[concept] {Cache L1/L2/L3 (SRAM)} child { node[concept, scale=0.7] {Rapide, proche CPU} } }
        child[concept color=popgreen!70] { node[concept] {RAM (DRAM)} child { node[concept, scale=0.7] {Volatile, programmes/données} } child { node[concept, scale=0.7] {DDR4/DDR5, Double canal} } }
        child[concept color=popgreen!70] { node[concept] {ROM / Firmware} child { node[concept, scale=0.7] {Non-volatile, BIOS/UEFI} } }
        child[concept color=popgreen!70] { node[concept] {Stockage Secondaire} child { node[concept, scale=0.7] {SSD NVMe / SATA} } child { node[concept, scale=0.7] {HDD (archivage)} } }
    }
    child[concept color=poppurple] {
        node[concept, minimum size=1.8cm] {Bus Système}
        child[concept color=poppurple!70] { node[concept] {Bus de Données} child { node[concept, scale=0.7] {Bidirectionnel, largeur (64b)} } }
        child[concept color=poppurple!70] { node[concept] {Bus d'Adresses} child { node[concept, scale=0.7] {Unidirectionnel, 2$^n$ octets} } }
        child[concept color=poppurple!70] { node[concept] {Bus de Contrôle} child { node[concept, scale=0.7] {Signaux: Read, Write, IRQ} } }
    }
    child[concept color=poppink] {
        node[concept, minimum size=1.8cm] {Cycle Instruction}
        child[concept color=poppink!70] { node[concept] {1. Fetch (Chargement)} }
        child[concept color=poppink!70] { node[concept] {2. Decode (Décodage)} }
        child[concept color=poppink!70] { node[concept] {3. Execute (Exécution)} }
        child[concept color=poppink!70] { node[concept] {4. Writeback (Écriture)} }
        child[concept color=poppink!70] { node[concept] {Pipeline / Superscalaire} }
    };
\end{tikzpicture}
\legende{Carte mentale de synthèse : les 4 piliers de l'architecture de l'ordinateur vus en 3ème (CPU, Mémoires, Bus, Cycle d'exécution) avec la distinction RISC/CISC.}
\end{popfigure}

\textbf{Tu as maintenant toutes les cartes en main pour réussir l'épreuve :} connais ton vocabulaire (registres, bus, cache, pipeline, RISC, CISC, swapping, IOPS, big.LITTLE), structure ta réponse en 3 étapes (Besoin $\rightarrow$ Choix $\rightarrow$ Argument technique), et relie \emph{systématiquement} chaque composant choisi à une notion du cours. C'est cette rigueur qui fait la différence entre une note moyenne et l'excellence au BEPC. \emph{Bon courage pour la suite du programme !}

\newpage

\coursec{Activité d'intégration}

\coursec{Leçon 2 : Description de l'architecture de l'ordinateur}

\courssub{Objectifs de l'activité}

À la fin de cette activité d'intégration, tu seras capable de :

\begin{itemize}
\item \cle{identifier}, sur une fiche technique réelle, les différentes unités fonctionnelles d'un ordinateur étudiées dans la leçon (unité centrale de traitement, mémoires, bus, périphériques) ;
\item \cle{comparer} plusieurs configurations matérielles en mobilisant les notions d'architecture de l'ordinateur (rôle du processeur, fréquence, nombre de cœurs, capacité de la mémoire vive, type de stockage) ;
\item \cle{distinguer} les architectures \cle{RISC} et \cle{CISC} et expliquer leurs conséquences pratiques sur le choix d'une machine ;
\item \cle{rédiger et justifier} une réponse argumentée, c'est-à-dire un conseil technique appuyé sur des notions du cours et non sur des impressions ;
\item \cle{appliquer} les savoirs de l'unité à une situation concrète de la vie courante : l'équipement d'un cybercafé au Cameroun.
\end{itemize}

\begin{aretenir}[Compétence visée]
Cette activité mobilise le savoir-faire essentiel de la leçon : \emph{« appliquer les notions de l'unité à des situations concrètes de la vie courante »} et \emph{« rédiger et justifier une démarche, une réponse ou une conclusion »}. Un bon technicien, même débutant, ne dit pas « cette machine est meilleure » : il dit \emph{pourquoi}, en citant les composants et leur rôle.
\end{aretenir}

\courssub{Situation}

Monsieur Kamdem est le propriétaire du cybercafé « \emph{Espace Net} », situé au quartier Marché B à Bafoussam, dans la région de l'Ouest. Son cybercafé compte dix postes de travail utilisés chaque jour par des élèves, des étudiants et des commerçants, principalement pour :

\begin{itemize}
\item la bureautique (rédaction de lettres et de CV sous un traitement de texte, tableaux sous un tableur, impressions) ;
\item la navigation sur Internet (recherche d'informations, messagerie, réseaux sociaux, démarches administratives en ligne) ;
\item occasionnellement, les jeux vidéo, très demandés par les jeunes clients le soir et le week-end.
\end{itemize}

Trois de ses vieilles machines tombent régulièrement en panne. Un fournisseur de matériel informatique de Douala lui propose trois ordinateurs neufs, dont voici les fiches techniques simplifiées :

\begin{center}
\begin{tabularx}{\linewidth}{|l|X|X|X|}
\hline
\rowcolor{popblueL} \textbf{Caractéristique} & \textbf{Machine A} & \textbf{Machine B} & \textbf{Machine C} \\
\hline
Type de processeur & Architecture CISC & Architecture CISC & Architecture RISC \\
\hline
Fréquence du processeur & 2,4 GHz & 3,2 GHz & 1,8 GHz \\
\hline
Nombre de cœurs & 2 cœurs & 4 cœurs & 8 cœurs \\
\hline
Mémoire vive (RAM) & 4 Go & 8 Go & 16 Go \\
\hline
Stockage & Disque dur 500 Go & Disque dur 1 To & Disque SSD 256 Go \\
\hline
Prix & 120 000 FCFA & 185 000 FCFA & 230 000 FCFA \\
\hline
\end{tabularx}
\end{center}

Monsieur Kamdem hésite. Son neveu lui a dit : « Prends la machine C, c'est la plus chère, donc la meilleure ! » Mais il se demande si la machine la plus chère est vraiment la plus adaptée à \emph{chacun} de ses usages. Il sollicite ta classe de 3ème, qui vient d'étudier l'architecture de l'ordinateur, pour l'aider à prendre une décision raisonnée.

\textbf{Problème posé :} quelle machine conseiller à Monsieur Kamdem pour un usage bureautique + Internet ? Pour un usage jeux vidéo ? Et surtout : comment \emph{justifier} ce conseil avec les notions du cours ?

\courssub{Consignes}

Par groupes de quatre élèves, travaillez les questions suivantes dans l'ordre. Chaque réponse doit être rédigée en phrases et \emph{justifiée} par une notion du cours.

\begin{enumerate}
\item \textbf{Repérage.} Pour chacune des trois machines, identifie sur la fiche technique les éléments qui correspondent aux unités fonctionnelles étudiées dans la leçon : l'unité centrale de traitement (CPU), la mémoire centrale (mémoire vive), la mémoire de masse (stockage). Précise le rôle de chacune de ces unités en une phrase.
\item \textbf{Analyse du processeur.} Compare les trois processeurs selon trois critères : l'architecture (RISC ou CISC), la fréquence d'horloge et le nombre de cœurs. Pour chaque critère, explique ce qu'il signifie et quel avantage il procure à l'utilisateur.
\item \textbf{Analyse des mémoires.} Compare les capacités de RAM des trois machines. Rappelle pourquoi une mémoire vive plus grande permet un travail plus fluide (plusieurs applications ouvertes en même temps). Compare aussi les deux types de stockage proposés (disque dur et disque SSD).
\item \textbf{Premier avis argumenté.} Rédige un court avis (5 à 8 lignes) recommandant à Monsieur Kamdem une machine pour l'usage \emph{bureautique + Internet}. Ton avis doit contenir : la machine choisie, au moins \textbf{trois arguments techniques} tirés du cours, et une prise en compte du prix (un cybercafé est une petite entreprise qui doit maîtriser ses dépenses).
\item \textbf{Second avis argumenté.} Rédige un second avis (5 à 8 lignes) recommandant une machine pour l'usage \emph{jeux vidéo}, en justifiant pourquoi les besoins sont différents de ceux de la bureautique.
\item \textbf{Restitution.} Chaque groupe présente oralement ses deux avis en deux minutes. La classe complète ensuite au tableau la grille de critères suivante :

\begin{center}
\begin{tabularx}{\linewidth}{|l|X|X|X|}
\hline
\rowcolor{popblueL} \textbf{Critère} & \textbf{Machine A} & \textbf{Machine B} & \textbf{Machine C} \\
\hline
Rapidité du processeur & & & \\
\hline
Capacité à faire plusieurs tâches & & & \\
\hline
Confort de la mémoire vive & & & \\
\hline
Rapidité du stockage & & & \\
\hline
Rapport qualité/prix & & & \\
\hline
\end{tabularx}
\end{center}

\end{enumerate}

\courssub{Ressources à mobiliser}

Pour réussir cette activité, tu dois mobiliser les savoirs suivants, vus dans la leçon :

\begin{itemize}
\item \textbf{L'architecture d'un ordinateur} : c'est l'organisation des grandes unités fonctionnelles qui le composent : l'unité centrale de traitement (CPU), les mémoires (centrale et de masse), les bus qui les relient et les interfaces d'entrées/sorties.
\item \textbf{Le CPU} : c'est le « cerveau » de l'ordinateur. Il exécute les instructions des programmes. Sa rapidité dépend notamment de sa \cle{fréquence d'horloge} (en GHz : nombre d'opérations élémentaires par seconde) et de son \cle{nombre de cœurs} (un cœur = une unité capable de traiter des instructions ; plusieurs cœurs permettent de traiter plusieurs tâches en parallèle).
\item \textbf{La mémoire vive (RAM)} : mémoire de travail, rapide mais volatile, qui contient les programmes et les données en cours d'utilisation. Plus elle est grande, plus on peut ouvrir d'applications simultanément sans ralentissement.
\item \textbf{La mémoire de masse} : stockage permanent (disque dur ou disque SSD). Le disque \cle{SSD} est beaucoup plus rapide qu'un disque dur classique, mais coûte plus cher à capacité égale.
\item \textbf{Les architectures RISC et CISC} : dans une architecture \cle{CISC} (\emph{Complex Instruction Set Computer}), le processeur comprend un grand nombre d'instructions complexes ; dans une architecture \cle{RISC} (\emph{Reduced Instruction Set Computer}), le jeu d'instructions est réduit à des instructions simples et rapides, exécutées très efficacement, souvent avec une meilleure économie d'énergie et de très bonnes performances par cœur.
\item \textbf{Le cycle d'exécution d'une instruction} : le CPU répète sans cesse les étapes : chercher l'instruction en mémoire, la décoder, l'exécuter, puis écrire le résultat. Une fréquence plus élevée signifie davantage de cycles par seconde.
\end{itemize}

\begin{methode}[Rédiger un avis technique argumenté]
Pour convaincre un client (ou un examinateur au BEPC), suis cette démarche en 4 étapes :
\begin{enumerate}
\item \textbf{Identifier le besoin} : quel usage ? (bureautique, jeu, montage vidéo, serveur…). Les besoins ne sont pas les mêmes.
\item \textbf{Sélectionner les composants critiques} : pour cet usage, qu'est-ce qui compte le plus ? (CPU puissant ? beaucoup de RAM ? SSD rapide ?).
\item \textbf{Comparer les valeurs chiffrées} : cite les nombres (fréquence, cœurs, Go, type disque) en les reliant au besoin identifié.
\item \textbf{Conclure avec le critère économique} : le choix est-il justifié par le rapport performance/prix ? Une machine surdimensionnée est un gaspillage ; une machine sous-dimensionnée crée de l'insatisfaction.
\end{enumerate}
\emph{Astuce :} utilise toujours la structure « Je choisis X \textbf{car} (argument technique 1), \textbf{car} (argument technique 2), \textbf{car} (argument technique 3). Son prix de Y FCFA est \textbf{adapté} car… »
\end{methode}

\courssub{Démarche et corrigé commenté}

\begin{exoresolu}[Conseiller le cybercafé « Espace Net »]
\textbf{Énoncé.} À partir des trois fiches techniques de Monsieur Kamdem : (1) identifier les unités fonctionnelles sur chaque fiche ; (2) comparer les configurations ; (3) rédiger un avis argumenté pour un usage bureautique + Internet, puis pour un usage jeux vidéo.

\tcblower

\textbf{Étape 1 : repérage des unités fonctionnelles.}

Sur chaque fiche, on retrouve les mêmes grandes unités de l'architecture d'un ordinateur :
\begin{itemize}
\item \emph{Unité centrale de traitement (CPU)} : la ligne « type de processeur » + « fréquence » + « nombre de cœurs ». C'est l'unité qui exécute les instructions des programmes.
\item \emph{Mémoire centrale (RAM)} : la ligne « mémoire vive ». C'est la mémoire de travail, volatile, qui accueille les programmes en cours d'exécution.
\item \emph{Mémoire de masse} : la ligne « stockage » (disque dur ou SSD). C'est le stockage permanent des fichiers et des logiciels.
\end{itemize}
Les bus (non mentionnés sur la fiche mais présents dans toute machine) assurent le transport des données et des instructions entre ces unités.

\textbf{Étape 2 : comparaison des processeurs.}

\begin{itemize}
\item \emph{Architecture} : les machines A et B utilisent un processeur CISC (jeu d'instructions riche et complexe, classique sur les ordinateurs de bureau) ; la machine C utilise un processeur RISC (instructions simples, exécutées très efficacement).
\item \emph{Fréquence} : B est la plus rapide en cycles par seconde (3,2 GHz), puis A (2,4 GHz), puis C (1,8 GHz). Mais attention : à architecture différente, la fréquence seule ne suffit pas à conclure, car un cœur RISC accomplit son travail de manière très efficace à chaque cycle.
\item \emph{Nombre de cœurs} : C domine nettement (8 cœurs contre 4 pour B et 2 pour A) : elle peut traiter beaucoup plus de tâches en parallèle.
\end{itemize}

\textbf{Étape 3 : comparaison des mémoires.}

La RAM passe de 4 Go (A) à 8 Go (B) puis 16 Go (C) : la machine C permet d'ouvrir simultanément beaucoup plus d'applications et d'onglets Internet sans ralentir. Pour le stockage, C dispose d'un SSD : démarrage et chargement des programmes beaucoup plus rapides qu'avec les disques durs de A et B, même si sa capacité (256 Go) est plus faible — ce qui n'est pas gênant pour un poste de cybercafé où les fichiers des clients ne sont pas conservés.

\textbf{Étape 4 : avis pour l'usage bureautique + Internet (exemple de rédaction).}

\emph{« Pour la bureautique et Internet, nous conseillons la machine B. Ces usages demandent des ressources modérées : son processeur CISC à 4 cœurs cadencé à 3,2 GHz exécute rapidement les instructions d'un traitement de texte ou d'un navigateur, et ses 8 Go de RAM permettent d'ouvrir confortablement plusieurs applications et onglets en même temps. Son disque dur de 1 To offre une large capacité de stockage pour les documents des clients. Enfin, à 185 000 FCFA, elle est 45 000 FCFA moins chère que la machine C, dont la puissance serait gaspillée pour ces tâches simples. La machine A, avec seulement 2 cœurs et 4 Go de RAM, risque des ralentissements dès que plusieurs programmes sont ouverts. »}

\textbf{Étape 5 : avis pour l'usage jeux vidéo (exemple de rédaction).}

\emph{« Pour les jeux vidéo, nous conseillons la machine C. Les jeux sont des programmes gourmands qui sollicitent fortement le processeur et la mémoire : les 8 cœurs de son processeur RISC traitent efficacement les nombreuses tâches simultanées du jeu, ses 16 Go de RAM évitent les ralentissements et les temps de chargement, et son disque SSD lance les jeux beaucoup plus vite qu'un disque dur classique. Sa fréquence plus basse (1,8 GHz) est compensée par l'efficacité de l'architecture RISC et le grand nombre de cœurs. Son prix élevé (230 000 FCFA) se justifie ici par l'usage intensif, alors qu'il serait inutile pour la simple bureautique. »}

\textbf{Étape 6 : grille de critères complétée (correction collective).}

\begin{center}
\begin{tabularx}{\linewidth}{|l|X|X|X|}
\hline
\rowcolor{popblueL} \textbf{Critère} & \textbf{Machine A} & \textbf{Machine B} & \textbf{Machine C} \\
\hline
Rapidité du processeur & Moyenne & Bonne & Très bonne (8 cœurs RISC) \\
\hline
Capacité multitâche & Faible (2 cœurs) & Bonne (4 cœurs) & Excellente (8 cœurs) \\
\hline
Confort de la RAM & Limite (4 Go) & Confortable (8 Go) & Très large (16 Go) \\
\hline
Rapidité du stockage & Lente (disque dur) & Lente (disque dur) & Très rapide (SSD) \\
\hline
Rapport qualité/prix & Correct & Très bon & Bon (usage intensif) \\
\hline
\end{tabularx}
\end{center}

\textbf{Conclusion.} Il n'existe pas de machine « meilleure » dans l'absolu : le bon choix dépend de l'\emph{usage}. La machine B est le meilleur compromis pour la bureautique et Internet ; la machine C est justifiée pour les jeux vidéo.
\end{exoresolu}

\begin{attention}[Erreurs fréquentes à éviter]
\begin{itemize}
\item Conclure uniquement à partir du prix (« la plus chère est la meilleure ») ou d'un seul critère (« la fréquence la plus haute gagne ») : il faut croiser \emph{plusieurs} critères et les relier à l'usage.
\item Comparer directement les fréquences de deux architectures différentes (RISC et CISC) sans précaution : l'efficacité par cycle n'est pas la même.
\item Confondre mémoire vive (RAM, volatile, mémoire de travail) et mémoire de masse (disque dur ou SSD, stockage permanent).
\item Oublier de justifier : un avis technique sans argument tiré du cours n'est pas une réponse attendue au BEPC.
\end{itemize}
\end{attention}

\courssub{Bilan de l'activité}

Cette activité t'a permis de mettre en évidence plusieurs idées essentielles de la leçon :

\begin{itemize}
\item \textbf{L'architecture n'est pas une notion abstraite} : derrière chaque ligne d'une fiche technique (processeur, RAM, stockage) se cache une unité fonctionnelle précise, dont tu sais désormais expliquer le rôle.
\item \textbf{Le choix d'une machine se justifie par l'usage} : une configuration excellente pour les jeux vidéo peut être un gaspillage d'argent pour de la bureautique. Le bon raisonnement part du besoin, puis remonte vers les composants.
\item \textbf{Les architectures RISC et CISC ont des conséquences concrètes} : elles montrent qu'on ne peut pas juger un processeur à sa seule fréquence ; le nombre de cœurs et l'efficacité du jeu d'instructions comptent tout autant.
\item \textbf{Tu as pratiqué une démarche complète} : observer des données, les analyser avec les notions du cours, comparer selon des critères explicites, puis rédiger et défendre une conclusion argumentée — exactement la compétence « rédiger et justifier une réponse » évaluée au BEPC.
\end{itemize}

\begin{exercice}[Pour aller plus loin]
Monsieur Kamdem te consulte à nouveau : un client lui propose un ordinateur d'occasion à 90 000 FCFA, avec un processeur CISC à 2 cœurs cadencé à 3,0 GHz, 4 Go de RAM et un disque dur de 320 Go. En rédigeant un avis argumenté de 5 à 8 lignes, explique si cette offre est intéressante pour remplacer un poste de bureautique, en la comparant à la machine A (120 000 FCFA) et en tenant compte du fait qu'il s'agit d'un matériel d'occasion.
\end{exercice}

\newpage

\coursec{Méthodes et exercices résolus}

\coursec{Leçon N02 : Description de l'architecture de l'ordinateur}

\courssub{Qu'est-ce que l'architecture d'un ordinateur ?}

\begin{definition}[Architecture d'un ordinateur]
L'\cle{architecture d'un ordinateur} désigne l'organisation logique et physique de ses composants principaux, ainsi que les règles qui régissent leurs échanges d'informations. Au collège, on étudie le modèle de \cle{von Neumann} (architecture à programme enregistré) : le programme et les données sont stockés dans une même mémoire, et le processeur les traite séquentiellement.
\end{definition}

\begin{popfigure}
\begin{tikzpicture}[
    node distance=1.2cm and 1.5cm,
    block/.style={draw=popblue, thick, rounded corners, fill=popblueL, minimum height=1.8cm, minimum width=3.2cm, align=center, font=\small},
    io/.style={draw=poporange, thick, rounded corners, fill=poporangeL, minimum height=1.4cm, minimum width=2.8cm, align=center, font=\small},
    bus/.style={draw=popdark, very thick, -{Stealth[length=3mm]}, popdark},
    busbi/.style={draw=popdark, very thick, -{Stealth[length=3mm]}, popdark, double distance=1.5pt},
    labelbus/.style={font=\scriptsize, text=popmuted, midway, above, sloped}
]
% Nœuds
\node[block, align=center] (cpu){\textbf{UNITÉ CENTRALE (CPU)}\\UAL +\ UC\ +\ Registres};
\node[block, above=of cpu, align=center] (mem){\textbf{MÉMOIRE VIVE (RAM)}\\\emph{Programmes \& Données en cours}};
\node[io, left=of cpu, align=center] (entree){\textbf{PÉRIPHÉRIQUES}\\\textbf{D'ENTRÉE}\\(Clavier, Souris, Scanner\ldots)};
\node[io, right=of cpu, align=center] (sortie){\textbf{PÉRIPHÉRIQUES}\\\textbf{DE SORTIE}\\(Écran, Imprimante\ldots)};
\node[io, below=of cpu, align=center] (masse){\textbf{MÉMOIRE DE MASSE}\\(Disque dur, Clé USB, SSD\ldots)};
% Bus système (représenté comme une colonne centrale)
\draw[busbi] (mem) -- node[labelbus, left] {Bus de données / Adresses / Commande} (cpu);
\draw[bus] (entree) -- node[labelbus, above] {Bus d'entrée} (cpu);
\draw[bus] (cpu) -- node[labelbus, above] {Bus de sortie} (sortie);
\draw[bus] (masse) -- node[labelbus, left] {Bus E/S} (cpu);
% Flèches annotées "Entrée -> Traitement -> Sortie"
\draw[popgreen, very thick, -{Stealth[length=3mm]}, dashed] (entree.east) -- ++(0.5,0) |- (cpu.north west);
\draw[popgreen, very thick, -{Stealth[length=3mm]}, dashed] (cpu.north east) |- ++(0.5,0) -- (sortie.west);
\end{tikzpicture}
\legende{Architecture de von Neumann : les composants sont reliés par des bus (données, adresses, commande). L'information suit le chemin : Entrée $\rightarrow$ Mémoire $\rightarrow$ CPU $\rightarrow$ Mémoire $\rightarrow$ Sortie.}
\end{popfigure}

\begin{aretenir}[Les 4 sous-ensembles fondamentaux]
Tout ordinateur, du smartphone au serveur, repose sur quatre sous-ensembles reliés par les \cle{bus} :
\begin{enumerate}
\item les \cle{périphériques d'entrée} (clavier, souris, webcam, scanner) : introduisent l'information ;
\item l'\cle{unité centrale de traitement} (CPU) : exécute les instructions (calcul, logique, contrôle) ;
\item les \cle{mémoires} : \cle{vive} (RAM, volatile) pour le travail en cours, \cle{de masse} (disque, SSD, clé USB, permanente) pour la sauvegarde ;
\item les \cle{périphériques de sortie} (écran, imprimante, haut-parleur) : restituent le résultat.
\end{enumerate}
\end{aretenir}

\begin{methode}[Identifier et justifier le rôle des composants dans une situation de vie]
Face à une situation concrète (achat pour un cybercafé, choix d'une machine pour la saisie des notes, panne d'un poste), suis cette démarche :
\begin{enumerate}
\item \textbf{Lire la situation} et repérer ce que l'utilisateur veut faire (saisir, calculer, stocker, afficher, imprimer).
\item \textbf{Associer chaque besoin à un composant} de l'architecture : périphériques d'\cle{entrée}, périphériques de \cle{sortie}, \cle{CPU} pour les calculs, \cle{mémoires} (RAM, disque) pour le stockage.
\item \textbf{Justifier chaque choix} en expliquant le rôle du composant : ne jamais se contenter de citer un nom.
\item \textbf{Schématiser} si demandé : placer l'unité centrale (UC = UAL + UT), les mémoires et les périphériques, reliés par les \cle{bus}.
\item \textbf{Conclure} par une phrase qui répond précisément à la question posée.
\end{enumerate}
\textbf{Réflexe à retenir :} toute information suit le chemin : \emph{entrée} $\rightarrow$ \emph{mémoire} $\rightarrow$ \emph{traitement (CPU)} $\rightarrow$ \emph{mémoire} $\rightarrow$ \emph{sortie}.
\end{methode}

\begin{exoresolu}[Le cybercafé de quartier]
Monsieur Etoundi veut ouvrir un cybercafé à Bafoussam. Il achète des ordinateurs et te demande de lui expliquer, pour chaque poste, quels composants permettront aux clients de : (a) taper un texte, (b) voir le texte à l'écran, (c) enregistrer le texte, (d) imprimer le texte. Présente ta réponse sous forme de tableau en justifiant le rôle de chaque composant.
\tcblower
\textbf{Raisonnement :} on applique la démarche : chaque besoin (a, b, c, d) est associé à un composant dont on justifie le rôle.
\begin{center}
\begin{tabularx}{\linewidth}{|l|X|X|}
\hline
\rowcolor{popblueL}
\textbf{Besoin du client} & \textbf{Composant utilisé} & \textbf{Justification (rôle)} \\
\hline
(a) Taper un texte & Clavier (périphérique d'entrée) & Il permet d'introduire les caractères dans l'ordinateur. \\
\hline
(b) Voir le texte & Écran (périphérique de sortie) & Il affiche le résultat du traitement à l'utilisateur. \\
\hline
(c) Enregistrer le texte & Disque dur / SSD (mémoire de masse) & Il conserve les données de façon permanente, même ordinateur éteint. \\
\hline
(d) Imprimer le texte & Imprimante (périphérique de sortie) & Elle reproduit le document sur papier. \\
\hline
\end{tabularx}
\end{center}
On ajoute que le texte tapé est d'abord placé en \cle{mémoire vive} (RAM), traité par le \cle{CPU} (mise en forme, correction orthographique), puis enregistré sur le disque dur.
\textbf{Conclusion :} chaque poste du cybercafé doit donc comporter des périphériques d'entrée (clavier, souris), des périphériques de sortie (écran, imprimante), une mémoire de masse (disque dur/SSD) et une unité centrale de traitement, tous reliés par les bus : c'est l'architecture de base de tout ordinateur.
\end{exoresolu}

\courssub{L'unité centrale de traitement (CPU)}

\begin{definition}[CPU / Processeur]
Le \cle{CPU} (\emph{Central Processing Unit}), ou \cle{processeur}, est le « cerveau » de l'ordinateur. Il lit, décode et exécute les instructions des programmes stockés en mémoire vive. Sa vitesse est donnée par sa \cle{fréquence d'horloge} (en GHz : milliards de cycles par seconde).
\end{definition}

\begin{popfigure}
\begin{tikzpicture}[
    node distance=1cm and 1.5cm,
    box/.style={draw=poppurple, thick, rounded corners, fill=poppurpleL, minimum height=1.5cm, minimum width=2.5cm, align=center, font=\small},
    reg/.style={draw=popgreen, thick, rounded corners, fill=popgreenL, minimum height=1cm, minimum width=2cm, align=center, font=\small},
    arr/.style={-Stealth, thick, popdark}
]
\node[box, align=center] (uc){\textbf{Unité de Commande (UC)}\\\emph{Dirige, séquence, contrôle}};
\node[box, right=of uc, align=center] (ual){\textbf{Unité Arithmétique \& Logique (UAL)}\\\emph{Calcule (+, -, $\times$, $\div$) \& Compare ($<, >, =$)}};
\node[reg, above=of uc, align=center] (reg1){Registre\\(Instruction)};
\node[reg, above=of ual, align=center] (reg2){Registres\\(Données / Accumulateur)};
\node[reg, below=of uc, align=center] (reg3){Registre\\(Adresse / Pointeur)};
\node[reg, below=of ual, align=center] (reg4){Registre\\(État / Flags)};
\draw[arr] (uc) -- node[midway, above, font=\scriptsize] {Signaux de contrôle} (ual);
\draw[arr, double distance=1pt] (reg1) -- (uc);
\draw[arr, double distance=1pt] (reg2) -- (ual);
\draw[arr, double distance=1pt] (reg3) -- (uc);
\draw[arr, double distance=1pt] (reg4) -- (ual);
\draw[arr, dashed] (uc.north west) -- ++(-1,0) node[left, font=\scriptsize] {Bus de commande};
\draw[arr, dashed] (ual.north east) -- ++(1,0) node[right, font=\scriptsize] {Bus de données};
\draw[arr, dashed] (reg3.south west) -- ++(-1,0) node[left, font=\scriptsize] {Bus d'adresses};
\end{tikzpicture}
\legende{Organisation interne simplifiée du CPU : l'UC commande, l'UAL calcule, les registres stockent temporairement. Les trois bus sortent vers la carte mère.}
\end{popfigure}

\begin{propriete}[Composition du CPU]
Le CPU comprend trois éléments indissociables :
\begin{itemize}
\item l'\cle{Unité de Commande (UC)} : elle \textbf{dirige} et \textbf{coordonne} toutes les opérations (lecture instruction, envoi signaux de contrôle vers UAL, mémoires, E/S) ;
\item l'\cle{Unité Arithmétique et Logique (UAL)} : elle \textbf{effectue} les calculs arithmétiques (addition, soustraction\ldots) et les opérations logiques (comparaisons ET, OU, NON) ;
\item les \cle{registres} : petites mémoires \emph{ultra-rapides} internes au CPU (quelques octets à quelques centaines d'octets) qui conservent l'instruction en cours, les données opérandes, l'adresse de l'instruction suivante, l'état du processeur (drapeaux).
\end{itemize}
\end{propriete}

\begin{methode}[Décrire l'unité centrale de traitement]
Pour décrire le CPU dans une copie d'examen :
\begin{enumerate}
\item \textbf{Définir} le CPU : c'est le « cerveau » de l'ordinateur, il exécute les instructions des programmes.
\item \textbf{Citer ses deux parties principales} et leur rôle :
\begin{itemize}
\item l'\cle{unité arithmétique et logique} (UAL) : effectue les calculs et les comparaisons ;
\item l'\cle{unité de commande} (UC) : dirige et coordonne le travail de tous les autres éléments.
\end{itemize}
\item \textbf{Citer les registres} : petites mémoires très rapides à l'intérieur du CPU qui stockent temporairement les données en cours de traitement.
\item \textbf{Préciser la fréquence d'horloge} (en GHz) : elle indique la rapidité du processeur ; plus elle est élevée, plus le CPU exécute d'instructions par seconde (à architecture comparable).
\end{enumerate}
\textbf{Formule de rédaction type :} « Le CPU est composé de l'UAL qui calcule, de l'UC qui commande, et de registres qui stockent temporairement les données. »
\end{methode}

\begin{exoresolu}[Deux ordinateurs pour le secrétariat]
Le principal d'un collège hésite entre deux ordinateurs : le poste A possède un processeur à \SI{2,4}{GHz}, le poste B un processeur à \SI{3,2}{GHz}. (1) Définis le CPU et cite ses deux parties principales. (2) Lequel des deux postes exécutera plus vite les calculs des bulletins de notes ? Justifie.
\tcblower
\textbf{(1)} Le \cle{CPU} (unité centrale de traitement, ou processeur) est le composant qui exécute les instructions des programmes : c'est le « cerveau » de l'ordinateur. Ses deux parties principales sont :
\begin{itemize}
\item l'\textbf{unité arithmétique et logique (UAL)}, qui effectue les opérations de calcul et les comparaisons ;
\item l'\textbf{unité de commande (UC)}, qui coordonne et dirige le fonctionnement de tous les éléments de l'ordinateur.
\end{itemize}
\textbf{(2)} La fréquence d'horloge indique le nombre de cycles d'horloge par seconde. À architecture équivalente, une fréquence plus élevée permet d'enchaîner plus de cycles d'instruction par seconde. Or :
\[ 3,2\ \text{GHz} > 2,4\ \text{GHz} \]
Le poste B réalise donc plus de cycles par seconde que le poste A.
\textbf{Conclusion :} le poste B (\SI{3,2}{GHz}) traitera plus rapidement les calculs des bulletins de notes, car sa fréquence d'horloge est plus élevée ; à architecture comparable, un processeur plus rapide exécute davantage d'instructions par seconde.
\end{exoresolu}

\courssub{Les mémoires et les bus}

\begin{definition}[Mémoire vive (RAM) et Mémoire de masse]
\begin{itemize}
\item La \cle{mémoire vive} (RAM, \emph{Random Access Memory}) est une mémoire \emph{volatile} : son contenu disparaît dès que l'ordinateur est éteint. Elle est \emph{rapide} et stocke les programmes et données \emph{en cours d'utilisation} par le CPU.
\item La \cle{mémoire de masse} (disque dur HDD, SSD, clé USB, carte SD) est une mémoire \emph{permanente} (non volatile) : elle conserve les fichiers sans alimentation électrique. Elle est \emph{plus lente} que la RAM et sert à la \emph{sauvegarde} et au \emph{transport} des données.
\end{itemize}
\end{definition}

\begin{propriete}[Unités de capacité et conversions]
Les capacités s'expriment en octets (o) et leurs multiples binaires (standard JEDEC / éducation nationale) :
\[ 1\ \text{Kio} = 1024\ \text{o} \quad;\quad 1\ \text{Mio} = 1024\ \text{Kio} \quad;\quad 1\ \text{Gio} = 1024\ \text{Mio} \quad;\quad 1\ \text{Tio} = 1024\ \text{Gio} \]
\textbf{Repère :} en 3ème, on utilise souvent les notations Ko, Mo, Go, To avec la convention $1\ \text{Go} = 1024\ \text{Mo}$.
\begin{itemize}
\item Pour passer d'une unité supérieure à une inférieure : \textbf{multiplier} par 1024 (ex : $8\ \text{Go} \rightarrow 8 \times 1024 = 8192\ \text{Mo}$).
\item Pour passer d'une unité inférieure à une supérieure : \textbf{diviser} par 1024 (ex : $2048\ \text{Mo} \rightarrow 2048 / 1024 = 2\ \text{Go}$).
\end{itemize}
\end{propriete}

\begin{definition}[Les trois bus du système]
Les composants communiquent via trois ensembles de fils électriques (bus) distincts :
\begin{itemize}
\item le \cle{bus de données} (bidirectionnel) : transporte les \textbf{valeurs} (instructions, opérandes, résultats) entre CPU, mémoire et périphériques ;
\item le \cle{bus d'adresses} (unidirectionnel CPU $\rightarrow$ autres) : transporte le \textbf{numéro de case mémoire} (adresse) où lire ou écrire ;
\item le \cle{bus de commande} (unidirectionnel CPU $\rightarrow$ autres) : transporte les \textbf{ordres} de l'UC : \texttt{LECTURE}, \texttt{ÉCRITURE}, \texttt{RÉINITIALISATION}, \texttt{INTERRUPTION}.
\end{itemize}
La \cle{largeur du bus d'adresses} (ex: 32 bits, 64 bits) détermine la mémoire maximale adressable ($2^{32}$ octets = 4 Gio pour 32 bits).
\end{definition}

\begin{methode}[Classer les mémoires et identifier les bus]
\begin{enumerate}
\item \textbf{Distinguer les deux grandes familles} : RAM (volatile, rapide, travail en cours) vs Mémoire de masse (permanente, plus lente, sauvegarde).
\item \textbf{Comparer les capacités} en utilisant les conversions $1\ \text{Go} = 1024\ \text{Mo}$.
\item \textbf{Identifier les trois bus} et leur rôle précis : Données (valeurs), Adresses (localisation), Commande (ordres).
\end{enumerate}
\textbf{Réflexe :} pour savoir si une mémoire est volatile, pose la question : « Que reste-t-il quand j'éteins ? »
\end{methode}

\begin{exoresolu}[La clé USB et la mémoire vive]
Amina enregistre son exposé sur une clé USB de \SI{8}{Go}. Son ordinateur possède \SI{4}{Go} de mémoire vive (RAM). (1) Quelle est la différence entre la RAM et la clé USB ? (2) Un film occupe \SI{2048}{Mo}. Combien de tels films Amina peut-elle enregistrer sur sa clé ? (3) Cite les trois bus d'un ordinateur et leur rôle.
\tcblower
\textbf{(1)} La \textbf{RAM} (mémoire vive) est une mémoire \emph{volatile} : son contenu disparaît dès que l'ordinateur s'éteint ; elle sert à stocker temporairement les programmes et données en cours d'utilisation. La \textbf{clé USB} est une mémoire de masse \emph{permanente} : elle conserve les fichiers même hors tension ; elle sert à sauvegarder et transporter des données.
\textbf{(2)} Conversion de la capacité de la clé en Mo :
\[ 8\ \text{Go} = 8 \times 1024\ \text{Mo} = 8192\ \text{Mo} \]
Nombre de films :
\[ \frac{8192\ \text{Mo}}{2048\ \text{Mo}} = 4 \]
Amina peut donc enregistrer \textbf{4 films} de \SI{2048}{Mo} sur sa clé.
\textbf{(3)} Les trois bus sont :
\begin{itemize}
\item le \textbf{bus de données}, qui transporte les données entre les composants ;
\item le \textbf{bus d'adresses}, qui indique l'emplacement des données en mémoire ;
\item le \textbf{bus de commande}, qui transporte les ordres de l'unité de commande (lecture, écriture).
\end{itemize}
\textbf{Conclusion :} la RAM et la clé USB se distinguent par leur volatilité et leur rôle ; la clé de \SI{8}{Go} contient exactement 4 films de \SI{2048}{Mo} ; les échanges entre composants passent par les bus de données, d'adresses et de commande.
\end{exoresolu}

\courssub{Le cycle d'exécution d'une instruction}

\begin{definition}[Cycle machine (Fetch–Decode–Execute)]
Le CPU exécute les instructions les unes après les autres en répétant un cycle élémentaire, cadencé par l'\cle{horloge} (fréquence en GHz). Chaque cycle d'instruction comprend 5 étapes :
\end{definition}

\begin{popfigure}
\begin{tikzpicture}[
    node distance=1.3cm,
    stepbox/.style={draw=popblue, thick, rounded corners, fill=popblueL, minimum height=1.2cm, minimum width=3.5cm, align=center, font=\small},
    arrow/.style={-Stealth, thick, popdark},
    decision/.style={draw=poporange, thick, diamond, aspect=2, fill=poporangeL, minimum width=2.5cm, align=center, font=\small}
]
\node[stepbox, align=center] (e1){\textbf{1. FETCH (Recherche)}\\\emph{UC lit instruction @ PC en RAM via Bus Adresses/Données}};
\node[stepbox, below=of e1, align=center] (e2){\textbf{2. DECODE (Décodage)}\\\emph{UC analyse l'instruction : opcode + opérandes}};
\node[stepbox, below=of e2, align=center] (e3){\textbf{3. FETCH OPERANDS}\\\emph{UC charge les données dans Registres (Bus Données)}};
\node[stepbox, below=of e3, align=center] (e4){\textbf{4. EXECUTE (Exécution)}\\\emph{UAL effectue calcul / comparaison / saut}};
\node[stepbox, below=of e4, align=center] (e5){\textbf{5. WRITEBACK (Rangement)}\\\emph{Résultat $\rightarrow$ Registre ou RAM (Bus Données)}};
\node[decision, right=of e3, xshift=2cm, align=center] (clock){\textbf{HORLOGE}\\\emph{Tick / Top}};
\draw[arrow] (e1) -- (e2);
\draw[arrow] (e2) -- (e3);
\draw[arrow] (e3) -- (e4);
\draw[arrow] (e4) -- (e5);
\draw[arrow, dashed, popmuted] (e5) -- ++(0,-0.5) -| (e1.north) node[midway, above, font=\scriptsize, text=popmuted] {Instruction suivante (PC+1)};
\draw[arrow, poporange] (clock) -- node[midway, above, font=\scriptsize] {Cadence} (e3);
\end{tikzpicture}
\legende{Le cycle d'exécution (Fetch–Decode–Execute–Writeback) répété à chaque top d'horloge. L'UC pilote les étapes 1, 2, 3, 5 ; l'UAL réalise l'étape 4.}
\end{popfigure}

\begin{methode}[Raconter le cycle d'exécution en étapes]
Pour décrire comment le CPU exécute une instruction (ex: addition) :
\begin{enumerate}
\item \textbf{Recherche de l'instruction (Fetch)} : l'unité de commande va chercher l'instruction en mémoire vive à l'adresse indiquée par le pointeur de programme (PC), via les bus.
\item \textbf{Décodage (Decode)} : l'unité de commande analyse l'instruction (code opération + opérandes) pour comprendre l'action à mener.
\item \textbf{Recherche des données (Fetch Operands)} : les données nécessaires sont amenées de la mémoire vers les registres du CPU via le bus de données.
\item \textbf{Exécution (Execute)} : l'UAL effectue l'opération demandée (calcul arithmétique ou test logique).
\item \textbf{Rangement du résultat (Writeback)} : le résultat est écrit dans un registre ou en mémoire vive.
\end{enumerate}
Ce cycle se répète des milliards de fois par seconde, au rythme de l'\cle{horloge}.
\textbf{Réflexe de rédaction :} présente toujours les étapes \emph{dans l'ordre} et précise quel élément du CPU (UC ou UAL) intervient à chaque étape.
\end{methode}

\begin{exoresolu}[Calcul de la moyenne d'un élève]
Pour calculer la moyenne de Tchoupo, le secrétaire tape ses notes dans un tableur, qui exécute l'instruction « additionner les notes ». Décris, en étapes ordonnées, comment le CPU exécute cette instruction, en précisant le rôle de l'unité de commande et de l'UAL.
\tcblower
\textbf{Étape 1 — Recherche de l'instruction (Fetch) :} l'\textbf{unité de commande} va chercher en mémoire vive l'instruction « additionner les notes » grâce aux bus (bus d'adresses pour l'adresse, bus de données pour l'instruction).
\textbf{Étape 2 — Décodage (Decode) :} l'unité de commande analyse l'instruction et comprend qu'il s'agit d'une addition : elle prépare l'UAL en sélectionnant l'opération « + ».
\textbf{Étape 3 — Recherche des données (Fetch Operands) :} les notes de Tchoupo, stockées en mémoire vive, sont amenées dans les \textbf{registres} du CPU via le bus de données.
\textbf{Étape 4 — Exécution (Execute) :} l'\textbf{UAL} effectue l'addition des notes (opération arithmétique).
\textbf{Étape 5 — Rangement (Writeback) :} le résultat (la somme des notes) est rangé en mémoire vive (ou dans un registre), puis pourra être affiché à l'écran via le bus de données.
\textbf{Conclusion :} l'exécution d'une instruction suit toujours le même cycle : recherche, décodage, recherche des données, exécution, rangement ; l'unité de commande dirige les étapes 1, 2, 3 et 5, tandis que l'UAL réalise le calcul à l'étape 4. Ce cycle, répété au rythme de l'horloge, permet au tableur d'obtenir la moyenne presque instantanément.
\end{exoresolu}

\courssub{Les architectures RISC et CISC}

\begin{definition}[Architectures CISC et RISC]
Deux philosophies de conception du jeu d'instructions (ISA) du processeur coexistent :
\begin{itemize}
\item \cle{CISC} (\emph{Complex Instruction Set Computer}) : le processeur comprend un \emph{grand nombre} d'instructions (souvent plusieurs centaines), de \emph{longueurs variables}, certaines très \emph{complexes} (réalisent plusieurs opérations de bas niveau en une seule instruction, ex: \texttt{STRING\_COPY}, calcul polynôme). Exemple historique : Intel x86 (PC de bureau, portables classiques).
\item \cle{RISC} (\emph{Reduced Instruction Set Computer}) : le processeur comprend un \emph{petit nombre} d'instructions (souvent < 100), de \emph{longueur fixe}, \emph{simples} et \emph{rapides} (une instruction = une opération élémentaire : charger, stocker, additionner, brancher). Les opérations complexes sont décomposées en plusieurs instructions simples par le compilateur. Exemple : ARM (smartphones, tablettes, Raspberry Pi, serveurs modernes Apple M-series, AWS Graviton).
\end{itemize}
\end{definition}

\begin{propriete}[Comparatif RISC vs CISC]
\begin{center}
\begin{tabularx}{\linewidth}{|l|X|X|}
\hline
\rowcolor{popblueL}
\textbf{Critère} & \textbf{RISC (ex: ARM)} & \textbf{CISC (ex: x86)} \\
\hline
Nombre d'instructions & Réduit (quelques dizaines) & Grand (plusieurs centaines) \\
\hline
Complexité des instructions & Simples, 1 cycle d'horloge idéal & Complexes, plusieurs cycles \\
\hline
Longueur des instructions & Fixe (ex: 32 bits) & Variable (1 à 15 octets) \\
\hline
Modes d'adressage & Peu (Load/Store uniquement) & Nombreux (mémoire-mémoire) \\
\hline
Registres & Nombreux (32 ou plus) & Historiquement peu (8-16) \\
\hline
Consommation d'énergie & \textbf{Faible} (idéal mobile/embarqué) & Plus élevée (refroidissement actif) \\
\hline
Décodage matériel & Simple, rapide, peu de transistors & Complexe, microcode, plus de transistors \\
\hline
Usage typique & Smartphones, tablettes, IoT, serveurs basse conso & PC de bureau, stations de travail, serveurs legacy \\
\hline
\end{tabularx}
\end{center}
\end{propriete}

\begin{methode}[Rédiger une comparaison RISC / CISC]
\begin{enumerate}
\item \textbf{Définir chaque architecture} en donnant le sens de l'acronyme et la caractéristique principale (nombre/complexité instructions).
\item \textbf{Comparer selon des critères précis} : nombre d'instructions, complexité, longueur, consommation, usage.
\item \textbf{Donner un exemple d'usage concret} : CISC pour les ordinateurs de bureau classiques (compatibilité logicielle historique) ; RISC pour les téléphones portables, tablettes, objets connectés (faible consommation, rapidité, dissipation thermique faible).
\item \textbf{Conclure} en rappelant l'idée directrice : RISC = peu d'instructions simples exécutées vite ; CISC = beaucoup d'instructions, parfois complexes.
\end{enumerate}
\textbf{Réflexe :} une comparaison se présente idéalement en \emph{tableau} avec les mêmes critères en lignes.
\end{methode}

\begin{exoresolu}[Ordinateur de bureau ou smartphone ?]
Le téléphone Android de Nadège et l'ordinateur de bureau de son frère n'utilisent pas le même type d'architecture de processeur. (1) Définis les architectures RISC et CISC. (2) Présente dans un tableau trois différences entre ces deux architectures. (3) Indique, en justifiant, quelle architecture convient le mieux au smartphone de Nadège.
\tcblower
\textbf{(1)} L'architecture \textbf{CISC} (\emph{Complex Instruction Set Computer}) est une architecture où le processeur dispose d'un grand nombre d'instructions, dont certaines complexes réalisant plusieurs opérations à la fois. L'architecture \textbf{RISC} (\emph{Reduced Instruction Set Computer}) est une architecture où le processeur dispose d'un nombre réduit d'instructions simples, chacune exécutée très rapidement (souvent en un cycle d'horloge).
\textbf{(2)} Tableau comparatif (trois différences) :
\begin{center}
\begin{tabularx}{\linewidth}{|l|X|X|}
\hline
\rowcolor{popblueL}
\textbf{Critère} & \textbf{RISC} & \textbf{CISC} \\
\hline
Nombre d'instructions & Réduit (peu d'instructions) & Grand (beaucoup d'instructions) \\
\hline
Complexité des instructions & Simples, une opération élémentaire chacune & Parfois complexes, plusieurs opérations combinées \\
\hline
Consommation d'énergie & Faible (peu de transistors, décodage simple) & Plus élevée (décodage complexe, microcode) \\
\hline
\end{tabularx}
\end{center}
\textbf{(3)} Le smartphone de Nadège fonctionne sur batterie et doit être réactif tout en chauffant peu. L'architecture RISC, avec ses instructions simples exécutées rapidement, son décodage matériel allégé et sa faible consommation d'énergie, lui convient donc le mieux.
\textbf{Conclusion :} RISC et CISC sont deux manières de concevoir le jeu d'instructions d'un processeur : RISC privilégie la simplicité, la rapidité et l'efficacité énergétique (appareils portables, embarqués), CISC la richesse des instructions et la compatibilité ascendante (ordinateurs de bureau classiques) ; c'est pourquoi le smartphone de Nadège utilise une architecture de type RISC (processeur ARM).
\end{exoresolu}

\courssub{Appliquer et justifier : démarche argumentée}

\begin{methode}[Construire une réponse argumentée au BEPC]
Pour rédiger une réponse complète, rigoureuse et bien justifiée :
\begin{enumerate}
\item \textbf{Reformuler la question} en une phrase pour montrer qu'elle est comprise.
\item \textbf{Annoncer la notion du cours} utilisée (définition exacte, avec le vocabulaire technique : CPU, UAL, UC, bus, mémoire vive, RISC, CISC\ldots).
\item \textbf{Développer le raisonnement} étape par étape : chaque affirmation est suivie d'une justification (\og car\ldots \fg{}, \og or\ldots \fg{}, \og donc\ldots \fg{}).
\item \textbf{Appuyer sur un schéma ou un tableau} si la question s'y prête (schéma de l'architecture, tableau comparatif, organigramme du cycle).
\item \textbf{Conclure} par une phrase qui répond directement à la question posée, sans introduire d'idée nouvelle.
\end{enumerate}
\textbf{Barème implicite :} la définition exacte rapporte des points, la justification rapporte des points, la conclusion rapporte des points ; une réponse juste mais non justifiée n'obtient jamais la totalité des points.
\end{methode}

\begin{exoresolu}[Pourquoi l'ordinateur ralentit-il ?]
Dans la salle informatique du collège, un ordinateur possède \SI{2}{Go} de RAM. Quand la professeure ouvre à la fois le traitement de texte, le navigateur Internet et le logiciel de notes, l'ordinateur devient très lent. Rédige une réponse argumentée expliquant ce phénomène et proposant une solution matérielle.
\tcblower
\textbf{Reformulation :} il s'agit d'expliquer pourquoi l'ouverture simultanée de plusieurs programmes ralentit l'ordinateur, et de proposer une amélioration matérielle.
\textbf{Notion utilisée :} la \cle{mémoire vive} (RAM) est la mémoire qui contient les programmes et les données \emph{en cours d'utilisation} ; chaque programme ouvert y occupe une place. Le CPU ne travaille qu'avec les données présentes en RAM.
\textbf{Raisonnement :}
\begin{itemize}
\item Chaque logiciel ouvert (traitement de texte, navigateur, logiciel de notes) est chargé dans la RAM pour être exécuté par le CPU.
\item Or la RAM de cet ordinateur est limitée à \SI{2}{Go}. Quand la somme des besoins mémoire des trois programmes approche ou dépasse cette capacité, l'ordinateur manque d'espace de travail.
\item Le système d'exploitation utilise alors la \cle{mémoire virtuelle} : il déplace les parties moins utilisées de la RAM vers le \textbf{disque dur} (fichier d'échange / swap), qui est \emph{beaucoup plus lent} (accès mécanique ou latence SSD supérieure à la RAM).
\item Les allers-retours constants entre la RAM rapide et le disque lent (thrashing) saturant le bus de données, ralentissent fortement tous les traitements : d'où la lenteur observée.
\end{itemize}
\textbf{Solution matérielle :} augmenter la capacité de la mémoire vive, par exemple passer de \SI{2}{Go} à \SI{4}{Go} ou \SI{8}{Go} de RAM, afin que les trois programmes puissent y résider simultanément sans recourir au disque dur.
\textbf{Conclusion :} l'ordinateur ralentit parce que la RAM de \SI{2}{Go} est insuffisante pour contenir les trois programmes ouverts en même temps, ce qui oblige la machine à utiliser le disque dur, bien plus lent, comme mémoire d'appoint ; la solution consiste à augmenter la quantité de mémoire vive.
\end{exoresolu}

\begin{attention}[Les 5 erreurs qui coûtent des points au BEPC]
\begin{enumerate}
\item \textbf{Confondre mémoire vive et mémoire de masse.} La RAM est volatile (elle s'efface à l'extinction), le disque dur/SSD est permanent. Écrire « la RAM conserve les fichiers » est une erreur qui fait perdre tous les points de la question.
\item \textbf{Confondre UAL et unité de commande.} L'UAL \emph{calcule} (arithmétique, logique), l'unité de commande \emph{dirige} (séquencement, signaux de contrôle). Attribuer les calculs à l'unité de commande est une inversion classique et sanctionnée.
\item \textbf{Se tromper dans les conversions d'unités.} Retenir $1\ \text{Go} = 1024\ \text{Mo}$ (et non $1000$) dans les exercices du cours ; poser la multiplication ou la division dans le bon sens (Go $\rightarrow$ Mo : on $\times 1024$ ; Mo $\rightarrow$ Go : on $/ 1024$) et vérifier l'ordre de grandeur du résultat.
\item \textbf{Mélanger RISC et CISC.} RISC = jeu d'instructions \emph{réduit} et simple (téléphones, tablettes, embarqué) ; CISC = jeu d'instructions \emph{complexe} et étendu (PC bureau x86). L'acronyme lui-même donne la réponse : \textbf{R}educed vs \textbf{C}omplex.
\item \textbf{Répondre sans justifier ni conclure.} Citer un composant sans expliquer son rôle, ou donner un résultat de calcul sans phrase de conclusion, fait perdre la moitié des points : toute réponse doit comporter la définition exacte, la justification (\og car\ldots \fg{}) et une conclusion qui répond à la question posée.
\end{enumerate}
\end{attention}

\begin{savaistu}[Le Cameroun et l'architecture matérielle]
Au Cameroun, l'équipement informatique des établissements (programme \og Un élève, un ordinateur \fg{}, cybercafés, centres de formation) repose massivement sur l'architecture \textbf{x86 (CISC)} pour les postes fixes (processeurs Intel Core / AMD Ryzen) et sur \textbf{ARM (RISC)} pour les tablettes éducatives, les smartphones et les mini-PC type Raspberry Pi utilisés dans les clubs robotique. La fibre optique CAMTEL (backbone national) relie les serveurs (souvent architecture x86 serveur ou ARM serveur) aux établissements. Comprendre l'architecture permet au technicien de maintenance de choisir la bonne barrette de RAM (DDR4/DDR5, format DIMM/SODIMM), le bon disque (SATA/NVMe) et de diagnostiquer une panne de bus (bétonnage, nappe défectueuse) ou une surchauffe CPU (pâte thermique, ventilateur).
\end{savaistu}

\begin{exercice}[Entraînement BEPC — Architecture complète]
\begin{enumerate}
\item \textbf{Schématise} l'architecture de von Neumann en plaçant : CPU (UC, UAL, Registres), Mémoire vive, Disque dur, Clavier, Écran, et les trois bus (légende chaque bus).
\item Un ordinateur affiche \texttt{Processeur : 2,8 GHz / 4 cœurs / Cache L3 8 Mo}. Explique le sens de \og 2,8 GHz \fg{} et de \og 4 cœurs \fg{}.
\item Convertis : (a) \SI{16}{Go} en Mo ; (b) \SI{32768}{Mo} en Go.
\item Classifie : Processeur Intel Core i5 (CISC ou RISC ?) ; Processeur Apple M2 (CISC ou RISC ?) ; Processeur ARM Cortex-A78 (CISC ou RISC ?). Justifie chaque choix par une caractéristique.
\item Un élève affirme : « La clé USB est plus rapide que la RAM car elle garde les fichiers. » Corrige cette affirmation en utilisant les termes \emph{volatile}, \emph{permanent}, \emph{débit}.
\end{enumerate}
\end{exercice}

\begin{corrige}[Exercice n°1 -- Architecture complète]
\begin{enumerate}
\item \textbf{Schéma :} Voir figure de la section 1 (Architecture de von Neumann). Vérifier : CPU au centre (UC+UAL+Registres), RAM au-dessus, Disque dur en dessous, Clavier à gauche (Entrée), Écran à droite (Sortie), Bus Données/Adresses/Commande entre CPU et RAM, Bus E/S vers périphériques.
\item \textbf{2,8 GHz} = fréquence d'horloge : \num{2,8} milliards de cycles par seconde. \textbf{4 cœurs} = 4 unités de traitement (CPU logiques) indépendants dans le même boîtier, capables d'exécuter 4 threads simultanément (parallélisme).
\item \textbf{Conversions :} (a) $16 \times 1024 = \SI{16384}{Mo}$. (b) $32768 / 1024 = \SI{32}{Go}$.
\item \textbf{Classification :}
\begin{itemize}
\item Intel Core i5 : \textbf{CISC} (architecture x86-64, jeu d'instructions complexe, longueur variable, microcode, forte consommation, ventilateur obligatoire).
\item Apple M2 : \textbf{RISC} (architecture ARMv8/ARM64, jeu d'instructions réduit, longueur fixe 32 bits, load/store, faible consommation, passive cooling possible).
\item ARM Cortex-A78 : \textbf{RISC} (par définition, architecture ARM = RISC).
\end{itemize}
\item \textbf{Correction :} La clé USB est une mémoire de masse \emph{permanente} (non volatile), mais son \emph{débit} (vitesse de transfert) est \emph{bien inférieur} à celui de la RAM. La RAM est \emph{volatile} (perte à l'arrêt) mais \emph{ultra-rapide} (dizaines de Go/s), c'est pourquoi le CPU ne travaille qu'avec elle. La clé USB sert au stockage/transport, pas au travail en cours.
\end{enumerate}
\end{corrige}

\newpage

\coursec{Exercices}

\coursec{Description de l'architecture de l'ordinateur}

\courssub{Je vérifie mes connaissances}

\begin{exercice}[QCM : Composants et fonctions]
\textbf{Associez} chaque composant de l'architecture de Von Neumann (colonne de gauche) à sa fonction principale (colonne de droite). \textbf{Une seule réponse par ligne.}

\begin{center}
\begin{tabularx}{\linewidth}{|l|X|}
\hline
\rowcolor{popblueL}
\textbf{Composant} & \textbf{Fonction} \\
\hline
1. Unité Arithmétique et Logique (UAL) & A. Stocke les programmes et données en cours d'exécution (volatil). \\
\hline
2. Unité de Contrôle (UC) & B. Effectue les calculs arithmétiques et les opérations logiques. \\
\hline
3. Mémoire vive (RAM) & C. Décode les instructions et commande la séquence des opérations. \\
\hline
4. Mémoire morte (ROM) & D. Contient le programme de démarrage (BIOS/UEFI), non volatile. \\
\hline
5. Bus de données & E. Transporte les informations (opérandes, résultats) entre CPU, mémoire et périphériques. \\
\hline
6. Bus d'adresses & F. Transporte l'adresse de la case mémoire ou du périphérique visé (unidirectionnel). \\
\hline
\end{tabularx}
\end{center}
\end{exercice}

\begin{exercice}[Vrai ou Faux ? Justifiez]
Pour chaque affirmation, cochez \textbf{Vrai} ou \textbf{Faux} et \textbf{justifiez} votre réponse en une phrase précise.

\begin{enumerate}
    \item \textbf{La mémoire ROM s'efface complètement lorsque l'on éteint l'ordinateur.}
    \item \textbf{L'architecture RISC utilise un jeu d'instructions plus complexe et plus fourni que l'architecture CISC.}
    \item \textbf{Le bus de données est unidirectionnel : il ne transporte les informations que du CPU vers la mémoire.}
    \item \textbf{L'Unité de Contrôle (UC) et l'Unité Arithmétique et Logique (UAL) forment ensemble le Processeur (CPU).}
    \item \textbf{Un périphérique de stockage de masse (disque dur, SSD) est directement adressable par le CPU comme la RAM.}
\end{enumerate}
\end{exercice}

\begin{exercice}[Définitions et vocabulaire]
\textbf{Donnez une définition claire et concise} des termes suivants en utilisant le vocabulaire technique de la leçon :
\begin{enumerate}
    \item Architecture de Von Neumann (principe de base).
    \item Cycle d'horloge (et lien avec la fréquence).
    \item Bus système (citer les trois bus et leur sens de circulation).
    \item Architecture CISC vs RISC (différence fondamentale sur le jeu d'instructions).
    \item Registre (rôle dans le CPU, exemples : PC, IR, ACC).
\end{enumerate}
\end{exercice}

\begin{exercice}[Classification des périphériques et mémoires]
\textbf{Classez} les éléments suivants dans le tableau ci-dessous en cochant la (ou les) bonne(s) case(s) : \textbf{Entrée}, \textbf{Sortie}, \textbf{Stockage de masse}, \textbf{Mémoire centrale (RAM/ROM)}.

\begin{center}
\begin{tabularx}{\linewidth}{|l|X|X|X|X|}
\hline
\rowcolor{popblueL}
\textbf{Élément} & \textbf{Entrée} & \textbf{Sortie} & \textbf{Stockage de masse} & \textbf{Mémoire centrale} \\
\hline
Clavier & & & & \\
\hline
Écran tactile & & & & \\
\hline
RAM (DDR4) & & & & \\
\hline
Disque Dur (HDD 1 To) & & & & \\
\hline
ROM (BIOS) & & & & \\
\hline
SSD NVMe (512 Go) & & & & \\
\hline
Imprimante laser & & & & \\
\hline
Souris optique & & & & \\
\hline
Clé USB (32 Go) & & & & \\
\hline
\end{tabularx}
\end{center}
\end{exercice}

\courssub{Je m'entraîne}

\begin{exercice}[Schéma de l'architecture de Von Neumann]
\textbf{Complétez} le schéma ci-dessous en nommant les \textbf{3 unités internes du CPU (1, 2, 3)} et en indiquant la nature et le \textbf{sens des flèches (4, 5, 6)} pour les trois bus (Bus de données, Bus d'adresses, Bus de contrôle) entre le CPU et la Mémoire, ainsi qu'entre le CPU et les Périphériques E/S.

\begin{center}
\begin{tikzpicture}[node distance=1.2cm and 1.5cm, every node/.style={font=\small}, >=Stealth]
% Nodes
\node[draw, rectangle, rounded corners, minimum width=3.5cm, minimum height=2cm, fill=popblue!10, align=center] (cpu){\textbf{CPU}\\(Unité Centrale)};
\node[draw, rectangle, rounded corners, minimum width=2.5cm, minimum height=1.5cm, fill=popgreen!10, right=of cpu, align=center] (mem){\textbf{Mémoire Principale}\\(RAM / ROM)};
\node[draw, rectangle, rounded corners, minimum width=2.5cm, minimum height=1.5cm, fill=poporange!10, below=of cpu] (io) {\textbf{Périphériques E/S}};

% Internal CPU labels (positioned relative to cpu center)
\node[draw, rectangle, minimum width=1.3cm, minimum height=0.8cm, fill=poppink!20] (ual) at ($(cpu.center)+(-0.9,0.5)$) {\textbf{1. ?}};
\node[draw, rectangle, minimum width=1.3cm, minimum height=0.8cm, fill=poppink!20] (uc) at ($(cpu.center)+(-0.9,-0.5)$) {\textbf{2. ?}};
\node[draw, rectangle, minimum width=1.3cm, minimum height=0.8cm, fill=popgold!20] (reg) at ($(cpu.center)+(0.9,0)$) {\textbf{3. ?}};

% Arrows CPU <-> Memory
\draw[<->, thick, popblue] (cpu.east) -- node[above, midway, font=\small] {\textbf{4. ?}} (mem.west);
\draw[<->, thick, popblue] ($(cpu.east)+(0,-0.3)$) -- node[below, midway, font=\small] {\textbf{5. ?}} ($(mem.west)+(0,-0.3)$);

% Arrows CPU <-> IO
\draw[<->, thick, poporange] (cpu.south) -- node[right, midway, font=\small] {\textbf{6. ?}} (io.north);

% Legend hints
\node[align=center, text width=3cm, font=\tiny\itshape, text=popmuted] at ($(cpu.east)!0.5!(mem.west) + (0,1.0)$) {Bus de données (bidir.) / Bus d'adresses (CPU$\to$Mem) / Bus de contrôle ?};
\node[align=center, text width=2.5cm, font=\tiny\itshape, text=popmuted] at ($(cpu.south)!0.5!(io.north) + (1.0,0)$) {Bus système (E/S) ?};
\end{tikzpicture}
\end{center}

\textbf{Légendez} ensuite le rôle de chaque unité numérotée 1, 2, 3 dans le CPU.
\end{exercice}

\begin{exercice}[Cycle d'exécution : « Multiplier 6 par 7 »]
\textbf{Décrivez pas à pas} le déroulement du cycle \textbf{Fetch – Decode – Execute – Writeback} pour l'instruction machine équivalente à \texttt{MUL R1, R2, R3} (contenu de R2 $\times$ contenu de R3 $\to$ R1). Précisez pour chaque étape :
\begin{enumerate}
    \item Le rôle du \textbf{PC (Program Counter)} et de l'\textbf{IR (Instruction Register)}.
    \item Le chemin emprunté par les données sur les \textbf{bus} (adresses, données).
    \item L'action de l'\textbf{UAL} et de l'\textbf{UC}.
    \item La mise à jour du \textbf{PC} pour l'instruction suivante.
\end{enumerate}
\end{exercice}

\begin{exercice}[Comparatif RISC vs CISC : Choix pour un smartphone]
Voici un tableau comparatif simplifié de deux processeurs hypothétiques : \textbf{Proc-A} (architecture CISC) et \textbf{Proc-B} (architecture RISC).

\begin{center}
\begin{tabularx}{\linewidth}{|l|X|X|}
\hline
\rowcolor{popblueL}
\textbf{Caractéristique} & \textbf{Proc-A (CISC)} & \textbf{Proc-B (RISC)} \\
\hline
Nombre d'instructions du jeu & $\approx$ 300 (complexes, longueur variable) & $\approx$ 50 (simples, longueur fixe 32 bits) \\
\hline
Modes d'adressage & Nombreux (mémoire-mémoire, registre-mémoire) & Peu (Load/Store uniquement registre-registre) \\
\hline
Exécution d'une instruction & Plusieurs cycles d'horloge (microcode) & 1 cycle d'horloge (pipeline efficace) \\
\hline
Densité de code (taille programme) & Élevée (code compact) & Faible (plus d'instructions nécessaires) \\
\hline
Complexité du matériel (décodage) & Élevée (décodeur complexe, microcode) & Faible (décodeur simple, câblé) \\
\hline
Consommation énergétique (à perf. égales) & Plus élevée & Plus faible \\
\hline
Nombre de registres généraux & Faible (ex: 8--16) & Élevé (ex: 32) \\
\hline
\end{tabularx}
\end{center}

\textbf{Questions :}
\begin{enumerate}
    \item Quel processeur (\textbf{Proc-A} ou \textbf{Proc-B}) est le plus adapté pour un \textbf{smartphone} moderne ? Justifiez en citant \textbf{au moins trois lignes du tableau}.
    \item Pourquoi la longueur fixe des instructions (RISC) facilite-t-elle le \textbf{pipeline} (exécution en pipeline) ?
    \item Expliquez pourquoi le compilateur a un rôle plus critique dans l'architecture RISC que dans l'architecture CISC.
\end{enumerate}
\end{exercice}

\begin{exercice}[Les bus système : rôles et sens]
Un élève affirme : \og \textit{Le bus de données transporte l'adresse de la case mémoire à lire, et le bus d'adresses transporte la valeur lue.\fg}
\begin{enumerate}
    \item Cette affirmation est-elle correcte ? \textbf{Corrigez-la} précisément.
    \item Quel bus est \textbf{unidirectionnel} (CPU $\to$ Mémoire/Périphérique) ? Quel bus est \textbf{bidirectionnel} ? Justifiez physiquement pourquoi.
    \item À quel moment du cycle \textbf{Fetch} le \textbf{Bus d'adresses} est-il utilisé ? Et le \textbf{Bus de données} ?
    \item Citez deux signaux de contrôle transportés par le \textbf{Bus de contrôle} lors d'une lecture mémoire.
\end{enumerate}
\end{exercice}

\courssub{Je me prépare au Probatoire}

\begin{exercice}[Cas pratique : Achat d'un ordinateur portable pour le BEPC]
\textbf{Contexte :} Armand, élève en classe de 3ème, dispose d'un budget de \textbf{250\,000 FCFA} pour s'acheter un ordinateur portable d'occasion. Il doit rédiger ses exposés (traitement de texte), consulter des cours en ligne (navigateur web avec plusieurs onglets), et regarder des vidéos tutorielles. Il hésite entre deux offres :

\begin{center}
\begin{tabularx}{\linewidth}{|l|X|X|}
\hline
\rowcolor{popblueL}
\textbf{Caractéristique} & \textbf{Offre Alpha} & \textbf{Offre Beta} \\
\hline
Processeur & Intel Core i3-5005U (2 cœurs, \textbf{4 threads}, 2.0 GHz, \textbf{CISC}, 15W TDP) & MediaTek MT8183 (8 cœurs : 4$\times$Cortex-A73 + 4$\times$Cortex-A53, \textbf{RISC}, $\sim$5W TDP) \\
\hline
RAM & 4 Go DDR3L (soudée, non extensible) & 4 Go LPDDR4X (soudée, non extensible) \\
\hline
Stockage & HDD 500 Go (5400 tr/min) & eMMC 64 Go \\
\hline
Écran & 14'' HD (1366$\times$768) TN & 11.6'' Full HD (1920$\times$1080) IPS Tactile \\
\hline
OS installé & Windows 10 Pro 64 bits & Chrome OS (Linux based) \\
\hline
Batterie (état) & 45 Wh (usure 40\%) & 38 Wh (usure 10\%) \\
\hline
Prix & 240\,000 FCFA & 230\,000 FCFA \\
\hline
\end{tabularx}
\end{center}

\textbf{Travail demandé :}
\begin{enumerate}
    \item \textbf{Analysez l'architecture mémoire :} Pour chaque offre, expliquez l'impact du type de stockage (HDD vs eMMC) sur le \textbf{temps de démarrage de l'OS} et l'\textbf{ouverture des applications}. Utilisez les notions de \textit{débit}, \textit{latence} et \textit{mécanique vs électronique}.
    \item \textbf{Analysez l'architecture processeur :} Comparez l'approche \textbf{CISC (Offre Alpha)} et \textbf{RISC (Offre Beta)} pour la gestion du \textbf{multitâche} (plusieurs onglets + vidéo + Word). Quel processeur gère le mieux le parallélisme ici ? Justifiez par le nombre de cœurs/threads et l'architecture \textit{big.LITTLE} de l'Offre Beta.
    \item \textbf{Identifiez le goulot d'étranglement (bottleneck) principal} de chaque configuration pour l'usage décrit (bureautique + navigation web).
    \item \textbf{Rédigez un conseil argumenté} (10--15 lignes) pour aider Armand à choisir, en mobilisant \textbf{au moins quatre notions distinctes} de la leçon (ex: cycle d'horloge, bus, mémoire volatile/non volatile, architecture CPU, périphériques E/S, consommation énergétique).
\end{enumerate}
\end{exercice}

\begin{exercice}[Analyse de fiches techniques : Serveur de fichiers vs Embarqué]
Deux fiches techniques de processeurs sont fournies :

\begin{center}
\begin{tabularx}{\linewidth}{|l|X|X|}
\hline
\rowcolor{popblueL}
\textbf{Paramètre} & \textbf{Processeur X (Xeon Silver 4214)} & \textbf{Processeur Y (ARM Cortex-A53)} \\
\hline
Architecture & x86-64 (\textbf{CISC}) & ARMv8-A (\textbf{RISC}) \\
\hline
Cœurs / Threads & 12 cœurs / 24 threads & 4 cœurs / 4 threads \\
\hline
Fréquence de base / Turbo & 2.2 GHz / 3.2 GHz & 1.5 GHz (fixe) \\
\hline
Cache L3 & 16.5 Mo (partagé) & 512 Ko -- 1 Mo (partagé) \\
\hline
TDP (Enveloppe thermique) & 85 W & $\sim$ 2--3 W \\
\hline
Jeu d'instructions & Complexe (CISC), micro-opérations & RISC, longueur fixe 32 bits (AArch64) \\
\hline
Support mémoire max & 1 To DDR4 ECC (6 canaux) & 4 Go -- 8 Go LPDDR3/4 (2 canaux) \\
\hline
Virtualisation & Matérielle avancée (VT-x, VT-d) & Matérielle (EL2) \\
\hline
\end{tabularx}
\end{center}

\textbf{Questions :}
\begin{enumerate}
    \item \textbf{Identifiez} l'architecture (CISC ou RISC) de chaque processeur et \textbf{expliquez} comment la différence de \textbf{complexité du décodage} impacte la surface de puce et la consommation (TDP).
    \item Le \textbf{Processeur X} est destiné à un \textbf{serveur de base de données} (nombreuses requêtes simultanées, gros volumes de données). Justifiez ce choix par \textbf{trois caractéristiques} du tableau (cœurs, cache, mémoire, virtualisation).
    \item Le \textbf{Processeur Y} équipe une \textbf{box internet (routeur domestique)} ou une \textbf{tablette d'entrée de gamme}. Justifiez par \textbf{trois caractéristiques} (TDP, fréquence, architecture mémoire).
    \item \textbf{Calcul de performance théorique (PIPS) :} En supposant un CPI (Cycles Par Instruction) moyen de 1.5 pour X (grâce au superscalaire/out-of-order) et 1.0 pour Y (pipeline RISC simple), calculez le nombre d'instructions par seconde (IPS) max théorique pour chacun à leur fréquence max. Lequel est le plus puissant en brute force ? Pourquoi ce chiffre ne suffit-il pas à comparer les performances réelles sur une application donnée ?
\end{enumerate}
\end{exercice}

\begin{exercice}[Scénario : Le cybercafé de l'école]
L'administration du Lycée Bilingue de Ngaoundéré veut installer un mini cybercafé de 10 postes pour les élèves. Le budget est serré. Le technicien propose de recycler de vieilles tours (CPU Core 2 Duo E8400 @ 3.0 GHz, 4 Go DDR2, HDD 160 Go) en installant \textbf{Linux Mint XFCE} (léger) au lieu de Windows 10.

\textbf{Questions :}
\begin{enumerate}
    \item \textbf{Architecture CPU :} Le Core 2 Duo est une architecture \textbf{CISC} (x86). Expliquez pourquoi, malgré son âge (gravure 45 nm, 2008), il reste capable de faire tourner un OS moderne léger. Citez le rôle du \textbf{cache L2 partagé (6 Mo)} et de la \textbf{fréquence élevée} pour le mono-thread.
    \item \textbf{Mémoire et Bus :} La RAM est en \textbf{DDR2} (bus de données 64 bits, fréquence 800 MHz). Calculez le \textbf{débit théorique maximal} de ce bus en Go/s. Ce débit est-il suffisant pour afficher une vidéo HD (1080p, $\approx$ 0,2--0,4 Go/s débit brut non compressé 24 bits 30 fps) ? Expliquez le rôle de la \textbf{carte graphique intégrée (IGP)} et de sa mémoire dédiée (ou partagée) dans ce cas.
    \item \textbf{Stockage :} Le HDD 160 Go (7200 tr/min) est le maillon faible. Proposez une \textbf{upgrade à faible coût} (type, interface, capacité) qui transformerait la réactivité du poste. Justifiez par les temps d'accès (latence) et débits.
    \item \textbf{Cycle d'instruction concret :} Un élève clique sur l'icône \og Firefox \fg. Décrivez la chaîne complète : \textbf{Périphérique d'entrée} $\to$ \textbf{Contrôleur d'interruption} $\to$ \textbf{CPU (Fetch/Decode/Execute)} $\to$ \textbf{Disque (lecture binaire)} $\to$ \textbf{RAM} $\to$ \textbf{CPU (exécution code)} $\to$ \textbf{Carte vidéo} $\to$ \textbf{Écran}. Identifiez à chaque flèche le \textbf{bus concerné} (PCIe, SATA, DDR, PCIe/AGP, HDMI/VGA).
    \item \textbf{Éthique et maintenance :} Citez \textbf{deux risques} liés à l'utilisation de matériel ancien (sécurité, fiabilité) et \textbf{deux avantages} (écologique, pédagogique) de cette démarche de recyclage.
\end{enumerate}
\end{exercice}

\newpage

\coursec{Corrigés des exercices}

\begin{corrige}[Exercice 1 — QCM : Composants et fonctions]
\textbf{Correspondances :}
\begin{center}
\begin{tabularx}{\linewidth}{|l|X|}
\hline
\rowcolor{popblueL}
\textbf{Composant} & \textbf{Réponse justifiée} \\
\hline
1. \texttt{UAL} & \textbf{B} — L'\texttt{Unité Arithmétique et Logique} est l'atelier de calcul : additions, soustractions, comparaisons, ET/OU logiques. \\
\hline
2. \texttt{UC} & \textbf{C} — L'\texttt{Unité de Contrôle} lit le programme instruction par instruction, les décode et commande les autres unités. \\
\hline
3. \texttt{RAM} & \textbf{A} — Mémoire vive, \cle{volatile} : elle contient les programmes et données en cours d'utilisation et s'efface à l'extinction. \\
\hline
4. \texttt{ROM} & \textbf{D} — Mémoire morte, \cle{non volatile} : elle conserve le programme de démarrage (BIOS/UEFI) même hors tension. \\
\hline
5. Bus de données & \textbf{E} — Il transporte les valeurs (données, résultats) dans les deux sens entre CPU, mémoire et périphériques (\cle{bidirectionnel}). \\
\hline
6. Bus d'adresses & \textbf{F} — Il transporte l'adresse de la case mémoire visée, uniquement du CPU vers la mémoire (\cle{unidirectionnel}). \\
\hline
\end{tabularx}
\end{center}
\textbf{Points de vigilance :} ne pas confondre bus de données (bidirectionnel) et bus d'adresses (unidirectionnel) ; la ROM n'est pas une mémoire « de travail » mais une mémoire de démarrage.
\end{corrige}

\begin{corrige}[Exercice 2 — Vrai ou Faux ? Justifiez]
\begin{enumerate}
    \item \textbf{Faux.} La ROM est une mémoire \cle{non volatile} : son contenu (BIOS/UEFI) est conservé même ordinateur éteint ; c'est la RAM qui s'efface.
    \item \textbf{Faux.} C'est l'inverse : RISC (\emph{Reduced Instruction Set Computer}) utilise un jeu d'instructions \cle{réduit et simple}, tandis que CISC (\emph{Complex Instruction Set Computer}) en propose un grand nombre, complexes.
    \item \textbf{Faux.} Le bus de données est \cle{bidirectionnel} : le CPU envoie des données à écrire vers la mémoire et reçoit les données lues depuis la mémoire. C'est le bus d'adresses qui est unidirectionnel.
    \item \textbf{Vrai.} Le processeur (CPU) est constitué de l'UC (commande), de l'UAL (calcul) et des registres (stockage interne très rapide).
    \item \textbf{Faux.} Le disque dur ou le SSD n'est pas adressable directement par le CPU : les données doivent d'abord être \cle{recopiées en RAM} (via le contrôleur de stockage) avant d'être traitées.
\end{enumerate}
\textbf{Points de vigilance :} une justification est exigée pour chaque item ; « vrai » ou « faux » seul ne rapporte pas tous les points.
\end{corrige}

\begin{corrige}[Exercice 3 — Définitions et vocabulaire]
\begin{enumerate}
    \item \textbf{Architecture de Von Neumann :} modèle d'ordinateur dans lequel un \cle{programme et ses données} sont stockés dans la \emph{même mémoire} ; le CPU exécute les instructions \emph{séquentiellement}, une par une, en les allant chercher en mémoire (programme enregistré).
    \item \textbf{Cycle d'horloge :} plus petit intervalle de temps rythmant le fonctionnement du processeur, imposé par son horloge interne. La \cle{fréquence} est le nombre de cycles par seconde : $f = 1/T$ (ex. \SI{2}{GHz} $=$ 2 milliards de cycles par seconde).
    \item \textbf{Bus système :} ensemble de liaisons reliant CPU, mémoire et périphériques. Il comprend : le \textbf{bus de données} (bidirectionnel), le \textbf{bus d'adresses} (unidirectionnel, CPU $\to$ mémoire) et le \textbf{bus de contrôle} (signaux de commande : lecture, écriture, interruptions).
    \item \textbf{CISC vs RISC :} le CISC possède un jeu d'instructions \emph{nombreux et complexe} (instructions puissantes, longueur variable, exécution en plusieurs cycles) ; le RISC possède un jeu \emph{réduit et simple} (longueur fixe, exécution en un cycle environ), ce qui simplifie le matériel et réduit la consommation.
    \item \textbf{Registre :} petite mémoire interne au CPU, très rapide, stockant temporairement une donnée ou une adresse. Exemples : \textbf{\texttt{PC}} (compteur de programme : adresse de la prochaine instruction), \textbf{\texttt{IR}} (registre d'instruction : instruction en cours), \textbf{\texttt{ACC}} (accumulateur : résultat des calculs de l'UAL).
\end{enumerate}
\end{corrige}

\begin{corrige}[Exercice 4 — Classification des périphériques et mémoires]
\begin{center}
\begin{tabular}{|l|c|c|c|c|}
\hline
\rowcolor{popblueL}
\textbf{Élément} & \textbf{Entrée} & \textbf{Sortie} & \textbf{Stockage de masse} & \textbf{Mémoire centrale} \\
\hline
Clavier & X & & & \\
\hline
Écran tactile & X & X & & \\
\hline
RAM (DDR4) & & & & X \\
\hline
Disque dur (HDD \SI{1}{To}) & & & X & \\
\hline
ROM (BIOS) & & & & X \\
\hline
SSD NVMe (\SI{512}{Go}) & & & X & \\
\hline
Imprimante laser & & X & & \\
\hline
Souris optique & X & & & \\
\hline
Clé USB (\SI{32}{Go}) & & & X & \\
\hline
\end{tabular}
\end{center}
\textbf{Justifications :} l'écran tactile est à la fois entrée (le toucher) et sortie (l'affichage) ; la RAM et la ROM sont des mémoires \emph{centrales} car directement accessibles au CPU ; HDD, SSD et clé USB sont des mémoires de \emph{masse} : non volatiles, grandes capacités, mais non adressables directement par le CPU.
\textbf{Point de vigilance :} une clé USB n'est pas de la RAM malgré le mot « mémoire » : c'est un stockage de masse amovible.
\end{corrige}

\begin{corrige}[Exercice 5 — Schéma de l'architecture de Von Neumann]
\textbf{Unités internes du CPU :}
\begin{itemize}
    \item \textbf{1. \texttt{UAL}} (Unité Arithmétique et Logique) : effectue les calculs et les opérations logiques.
    \item \textbf{2. \texttt{UC}} (Unité de Contrôle) : décode les instructions et commande la séquence des opérations.
    \item \textbf{3. Registres} (\texttt{PC}, \texttt{IR}, \texttt{ACC}...) : mémoires internes très rapides stockant données, adresses et instruction en cours.
\end{itemize}
\textbf{Bus et sens des flèches :}
\begin{itemize}
    \item \textbf{4. Bus de données} : bidirectionnel (CPU $\leftrightarrow$ Mémoire) — transporte les valeurs lues ou à écrire.
    \item \textbf{5. Bus d'adresses} : unidirectionnel (CPU $\to$ Mémoire) — transporte l'adresse de la case visée.
    \item \textbf{6. Bus système vers les E/S} : bus de données (bidirectionnel) + bus de contrôle (signaux lecture/écriture, interruptions) reliant le CPU aux périphériques d'entrée/sortie.
\end{itemize}
\textbf{Points de vigilance :} le bus d'adresses part \emph{toujours} du CPU ; le bus de contrôle accompagne les deux autres bus dans les échanges avec la mémoire comme avec les périphériques.
\end{corrige}

\begin{corrige}[Exercice 6 — Cycle d'exécution : « Multiplier 6 par 7 »]
Supposons \texttt{R2 = 6} et \texttt{R3 = 7} ; l'instruction \texttt{MUL R1, R2, R3} se trouve à l'adresse mémoire $A$.
\begin{enumerate}
    \item \textbf{Fetch (recherche) :} le \texttt{PC} contient l'adresse $A$ ; l'UC place $A$ sur le \textbf{bus d'adresses} et active le signal de lecture sur le bus de contrôle. L'instruction codée remonte par le \textbf{bus de données} et est rangée dans l'\texttt{IR}.
    \item \textbf{Decode (décodage) :} l'\textbf{UC} analyse le contenu de l'\texttt{IR} : elle reconnaît une multiplication, identifie les opérandes (registres \texttt{R2}, \texttt{R3}) et la destination (\texttt{R1}). Le \texttt{PC} est incrémenté : $\text{PC} \leftarrow A + 1$ (adresse de l'instruction suivante).
    \item \textbf{Execute (exécution) :} l'UC commande l'\textbf{UAL} qui reçoit les contenus de \texttt{R2} (6) et \texttt{R3} (7) et calcule $6 \times 7 = 42$.
    \item \textbf{Writeback (rangement) :} le résultat 42 est écrit dans le registre \textbf{\texttt{R1}}. Le cycle recommence avec la nouvelle valeur du \texttt{PC}.
\end{enumerate}
\textbf{Bilan :} bus d'adresses utilisé au Fetch (envoi de $A$), bus de données utilisé au Fetch (retour de l'instruction) ; ici les opérandes étant déjà dans des registres, aucun accès mémoire supplémentaire n'est nécessaire.
\textbf{Point de vigilance :} le \texttt{PC} est mis à jour \emph{avant ou pendant} l'exécution, jamais après le Writeback : il doit toujours pointer sur la prochaine instruction.
\end{corrige}

\begin{corrige}[Exercice 7 — Comparatif RISC vs CISC : choix pour un smartphone]
\begin{enumerate}
    \item \textbf{Choix : Proc-B (RISC).} Justifications tirées du tableau :
    \begin{itemize}
        \item \emph{Consommation énergétique plus faible} : critère décisif pour un appareil sur batterie ;
        \item \emph{1 cycle par instruction / pipeline efficace} : exécution rapide et régulière malgré un matériel simple ;
        \item \emph{Complexité du matériel faible} (décodeur simple, câblé) : puce plus petite, moins chère, qui chauffe peu — pas de ventilateur dans un smartphone.
    \end{itemize}
    \item \textbf{Longueur fixe et pipeline :} avec des instructions de taille fixe (\SI{32}{bits}), le processeur sait \emph{exactement} où commence et où finit chaque instruction ; le décodage est uniforme et chaque étage du pipeline (Fetch, Decode, Execute...) met le même temps. On peut donc chevaucher plusieurs instructions comme sur une chaîne de montage. En CISC, la longueur variable impose un décodage complexe qui déséquilibre les étages.
    \item \textbf{Rôle du compilateur :} en RISC, le matériel est simple et n'offre que des instructions élémentaires : c'est le \cle{compilateur} qui doit décomposer les opérations complexes en suites d'instructions simples, choisir les registres (nombreux) et ordonner le code pour remplir le pipeline. En CISC, une partie de cette « intelligence » est câblée dans le processeur (microcode).
\end{enumerate}
\textbf{Point de vigilance :} citer explicitement les lignes du tableau dans la justification, comme demandé.
\end{corrige}

\begin{corrige}[Exercice 8 — Les bus système : rôles et sens]
\begin{enumerate}
    \item \textbf{Affirmation incorrecte.} Correction : c'est le \textbf{bus d'adresses} qui transporte l'adresse de la case mémoire à lire, et le \textbf{bus de données} qui transporte la valeur lue (ou à écrire). L'élève a inversé les deux bus.
    \item Le \textbf{bus d'adresses} est \cle{unidirectionnel} (CPU $\to$ mémoire/périphérique) : seul le CPU décide \emph{où} lire ou écrire, la mémoire ne choisit jamais d'adresse. Le \textbf{bus de données} est \cle{bidirectionnel} : les données circulent du CPU vers la mémoire (écriture) et de la mémoire vers le CPU (lecture).
    \item Pendant le \textbf{Fetch} : d'abord le \textbf{bus d'adresses} transporte l'adresse contenue dans le \texttt{PC} vers la mémoire ; ensuite le \textbf{bus de données} ramène l'instruction lue vers le registre \texttt{IR} du CPU.
    \item Deux signaux du \textbf{bus de contrôle} lors d'une lecture : le signal \textbf{MEMR} (lecture mémoire, \emph{Memory Read}) et le signal de validation d'adresse ; on peut aussi citer \textbf{MEMW} (écriture) ou les signaux d'\textbf{interruption} (IRQ).
\end{enumerate}
\end{corrige}

\begin{corrige}[Exercice 9 — Cas pratique : achat d'un ordinateur portable pour le BEPC]
\begin{enumerate}
    \item \textbf{Architecture mémoire :} le HDD de l'offre Alpha est un support \emph{mécanique} (plateaux à \SI{5400}{tr/min}, tête de lecture mobile) : latence élevée (plusieurs millisecondes) et débit modeste ($\approx$ \SI{80--100}{Mo/s}) ; le démarrage de Windows et l'ouverture des applications seront lents. L'eMMC de l'offre Beta est une mémoire \emph{électronique} (flash, sans pièce mobile) : latence bien plus faible, accès quasi instantané aux petits fichiers, d'où un démarrage de Chrome OS très rapide. Contrepartie : \SI{64}{Go} seulement, contre \SI{500}{Go}.
    \item \textbf{Architecture processeur :} l'offre Alpha (CISC, 2 cœurs / 4 threads, \SI{2}{GHz}) offre de bonnes performances par cœur mais peu de parallélisme. L'offre Beta (RISC, 8 cœurs \emph{big.LITTLE} : 4 cœurs puissants A73 + 4 cœurs économes A53) répartit mieux les tâches simultanées : les onglets et la vidéo sur les gros cœurs, les tâches de fond sur les petits. Pour le multitâche léger décrit, \textbf{Proc-B gère mieux le parallélisme}, avec en prime une consommation $\sim$3 fois plus faible (\SI{5}{W} contre \SI{15}{W}).
    \item \textbf{Goulots d'étranglement :} offre Alpha : le \textbf{HDD mécanique} (et \SI{4}{Go} de RAM juste pour Windows 10 + onglets) ; offre Beta : le \textbf{stockage de \SI{64}{Go}} vite saturé et l'absence de logiciels bureautiques « lourds » natifs sous Chrome OS.
    \item \textbf{Conseil argumenté (modèle de rédaction) :} « Armand, je te conseille l'\textbf{offre Beta}. Son processeur RISC exécute chaque instruction en environ un cycle d'horloge et ses 8 cœurs répartissent tes onglets, ta vidéo et ton traitement de texte sans ralentir, tout en consommant trois fois moins : ta batterie (usée à \SI{10}{\%} seulement) tiendra la journée de cours. Son stockage eMMC, mémoire flash non volatile sans mécanique, démarre le système en quelques secondes, contrairement au disque mécanique de l'Alpha dont la latence pénalise chaque ouverture de fichier. Les échanges entre CPU, mémoire et périphériques passent par les bus système : avec \SI{4}{Go} de RAM volatile et un système léger, ces bus ne seront pas saturés. Son écran tactile Full HD est un périphérique d'entrée/sortie plus confortable pour lire tes cours. Réserve toutefois une clé USB pour tes documents, car \SI{64}{Go} se remplissent vite. »
\end{enumerate}
\textbf{Barème indicatif (sur 20) :} Q1 : 4 pts ; Q2 : 5 pts ; Q3 : 3 pts ; Q4 : 8 pts (dont 4 pts pour la mobilisation correcte d'au moins quatre notions de la leçon, 2 pts pour la cohérence de l'argumentation, 2 pts pour la clarté).
\textbf{Points de vigilance :} exiger les notions de latence, débit, volatile/non volatile, cycle d'horloge, bus, architecture RISC/CISC ; le choix final importe moins que la qualité de la justification.
\end{corrige}

\begin{corrige}[Exercice 10 — Analyse de fiches techniques : serveur vs embarqué]
\begin{enumerate}
    \item \textbf{X : CISC} (x86-64) ; \textbf{Y : RISC} (ARM). En CISC, le décodeur doit gérer des instructions complexes de longueur variable, souvent via du microcode : cela occupe une grande \emph{surface de puce} et consomme beaucoup (TDP \SI{85}{W}). En RISC, le décodage câblé d'instructions simples de longueur fixe est compact et économe (TDP \SI{2--3}{W}).
    \item \textbf{Processeur X pour un serveur de base de données :} (a) \textbf{12 cœurs / 24 threads} pour traiter de nombreuses requêtes simultanées ; (b) \textbf{\SI{16,5}{Mo} de cache L3} pour garder les données fréquentes au plus près du CPU ; (c) \textbf{\SI{1}{To} de RAM ECC sur 6 canaux} pour de gros volumes avec correction d'erreurs ; (d) la \textbf{virtualisation matérielle} pour héberger plusieurs serveurs virtuels.
    \item \textbf{Processeur Y pour une box ou une tablette :} (a) \textbf{TDP de \SI{2--3}{W}} : pas de ventilateur, fonctionnement 24 h/24 sans surchauffe ; (b) \textbf{architecture RISC} simple et économe ; (c) mémoire \textbf{LPDDR} basse consommation (\SI{4--8}{Go} suffisants pour router ou naviguer).
    \item \textbf{Calcul des IPS :} $\text{IPS} = \dfrac{f}{\text{CPI}}$.
    \begin{align*}
    \text{X : } & \frac{\num{3,2} \times 10^{9}}{1,5} \approx \num{2,13} \times 10^{9} \text{ instructions/s par c\oe ur} \\
    \text{Y : } & \frac{\num{1,5} \times 10^{9}}{1,0} = \num{1,5} \times 10^{9} \text{ instructions/s par c\oe ur}
    \end{align*}
    Par cœur, X est plus puissant ; avec ses 12 cœurs, il écrase Y en « brute force ». Mais ce chiffre ne suffit pas : les performances réelles dépendent aussi du \textbf{cache}, de la \textbf{mémoire} (débit, latence), du \textbf{stockage}, du \textbf{compilateur} et de la nature de l'application (une instruction CISC peut faire le travail de plusieurs instructions RISC).
\end{enumerate}
\textbf{Barème indicatif (sur 20) :} Q1 : 4 pts ; Q2 : 5 pts ; Q3 : 4 pts ; Q4 : 7 pts (calculs 4 pts, discussion 3 pts).
\textbf{Points de vigilance :} vérifier la formule IPS $= f/\text{CPI}$ ; exiger la mention « par cœur » ; la comparaison réaliste doit citer au moins deux facteurs externes au CPU.
\end{corrige}

\begin{corrige}[Exercice 11 — Scénario : le cybercafé de l'école]
\begin{enumerate}
    \item \textbf{CPU :} bien que gravé en \SI{45}{nm} (2008), le Core 2 Duo reste un CISC x86 compatible avec tous les logiciels actuels ; sa \textbf{fréquence élevée (\SI{3}{GHz})} donne de bonnes performances en mono-tâche, et son \textbf{cache L2 de \SI{6}{Mo}} évite de nombreux accès à la lente RAM DDR2. Un OS léger comme Linux Mint XFCE demandant peu de ressources, le poste reste utilisable pour la bureautique et le web.
    \item \textbf{Débit du bus mémoire DDR2-800 :}
    \[ \text{Débit} = \frac{64 \text{ bits}}{8} \times 2 \times 800 \times 10^{6} = \num{12,8} \times 10^{9} \text{ octets/s} \approx \SI{12,8}{Go/s} \]
    (le facteur 2 vient du « double débit » DDR). Ce débit est \textbf{largement suffisant} pour une vidéo HD ($\approx$ \SI{0,2--0,4}{Go/s}). La \textbf{carte graphique intégrée (IGP)} décode et met en forme l'image ; elle utilise une partie de la RAM (mémoire partagée) ou sa mémoire dédiée comme mémoire vidéo, déchargeant le CPU.
    \item \textbf{Upgrade du stockage :} remplacer le HDD par un \textbf{SSD SATA de \SI{240--256}{Go}} (environ \SI{10000--15000}{FCFA}). Latence : de $\sim$\SI{10}{ms} (mécanique) à $\sim$\SI{0,1}{ms} (électronique) ; débit : de $\sim$\SI{100}{Mo/s} à $\sim$\SI{500}{Mo/s}. Démarrage et lancement des applications deviennent quasi instantanés.
    \item \textbf{Chaîne complète du clic sur « Firefox » :}
    \begin{itemize}
        \item Souris (périphérique d'entrée) $\to$ contrôleur d'interruption via le port \textbf{USB} ;
        \item Interruption $\to$ CPU : Fetch/Decode/Execute de la routine de traitement ;
        \item Le CPU ordonne la lecture du programme sur le disque via le bus \textbf{SATA} ;
        \item Le binaire de Firefox est copié du disque vers la \textbf{RAM} (bus mémoire \textbf{DDR2}) ;
        \item Le CPU exécute le code depuis la RAM (cycle Fetch--Decode--Execute) ;
        \item Le CPU envoie l'image à la carte vidéo (bus \textbf{PCIe}) puis à l'écran (liaison \textbf{VGA/HDMI}).
    \end{itemize}
    \item \textbf{Éthique et maintenance :} \emph{Risques} : failles de sécurité non corrigées sur du matériel ancien ; pannes matérielles (disques usés, condensateurs fatigués). \emph{Avantages} : geste écologique (lutte contre les déchets électroniques) et pédagogique (les élèves apprennent la maintenance et Linux à moindre coût).
\end{enumerate}
\textbf{Barème indicatif (sur 20) :} Q1 : 4 pts ; Q2 : 4 pts (calcul 2 pts) ; Q3 : 3 pts ; Q4 : 6 pts ; Q5 : 3 pts.
\textbf{Points de vigilance :} dans le calcul de débit, ne pas oublier la conversion bits $\to$ octets (division par 8) ni le double débit DDR ; dans la chaîne du clic, chaque flèche doit être associée à un bus nommé.
\end{corrige}

\newpage

\coursec{Fiche bilan}

\begin{popfigure}
\centering
\begin{tikzpicture}[
  mindmap,
  grow cyclic,
  every node/.style={concept, execute at begin node=\hskip0pt},
  concept color=popblue,
  level 1/.append style={level distance=4.5cm, sibling angle=60},
  level 2/.append style={level distance=3cm, sibling angle=45},
  text=white,
  font=\small\bfseries
]
\node[concept, scale=1.2, align=center]{Architecture\\ de l'ordinateur}
  child[concept color=poporange] { node[concept, align=center]{1. Définition\\ \& Rôle} }
  child[concept color=popgreen] { node[concept, align=center]{2. CPU\\ (UC, ALU, Registres)} }
  child[concept color=poppurple] { node[concept, align=center]{3. Mémoires\\ \& Bus} }
  child[concept color=poppink] { node[concept, align=center]{4. Cycle\\ d'exécution} }
  child[concept color=popgold] { node[concept, align=center]{5. Architecture\\ CISC} }
  child[concept color=popblue] { node[concept, align=center]{6. Architecture\\ RISC} };
\end{tikzpicture}
\legende{Carte mentale : Architecture de l'ordinateur (Classe de 3ème)}
\end{popfigure}

\begin{bilan}
\courssub{Définitions clés}
\begin{itemize}
  \item \textbf{Architecture d'un ordinateur} : Organisation logique et physique des composants (CPU, mémoire, entrées/sorties) et de leurs interconnexions (bus) qui détermine le fonctionnement de la machine.
  \item \textbf{CPU (Unité Centrale de Traitement)} : « Cerveau » de l'ordinateur. Composé de l'\textbf{Unité de Contrôle (UC)} (séquencement), de l'\textbf{ALU} (calculs arithmétiques/logiques) et des \textbf{registres} (stockage ultra-rapide).
  \item \textbf{Mémoire centrale (RAM)} : Volatile, rapide, stocke programmes et données en cours d'exécution. Adressable par octet.
  \item \textbf{Bus} : Ensemble de fils conducteurs assurant le transport de l'information. \textbf{Bus d'adresses} (unidir., sélection), \textbf{Bus de données} (bidir., valeurs), \textbf{Bus de contrôle} (ordres : lecture/écriture).
  \item \textbf{Cycle machine} : Suite d'étapes répétées pour chaque instruction : \textbf{Chargement} (Fetch) $\to$ \textbf{Décodage} (Decode) $\to$ \textbf{Exécution} (Execute) $\to$ \textbf{Écriture résultat} (Writeback).
  \item \textbf{CISC (Complex Instruction Set Computer)} : Jeu d'instructions riche, complexe, longueur variable, microcodé (ex. x86/Intel/AMD). Privilégie la densité de code.
  \item \textbf{RISC (Reduced Instruction Set Computer)} : Jeu d'instructions réduit, simple, longueur fixe, exécution directe (pipeline), architecture \emph{load/store} (ex. ARM, MIPS, RISC-V). Privilégie la vitesse d'horloge et le parallélisme.
\end{itemize}

\courssub{Comparatif RISC vs CISC (à connaître par cœur)}
\begin{center}
\begin{tabularx}{\linewidth}{|>{\bfseries}l|X|X|}
\hline
\rowcolor{popblueL}
\textbf{Critère} & \textbf{CISC} & \textbf{RISC} \\ \hline
Nombre d'instructions & Élevé ($>$ 200) & Faible ($<$ 100) \\ \hline
Longueur instructions & Variable & Fixe (ex. 32 bits) \\ \hline
Modes d'adressage & Nombreux & Réduits (souvent 3 max) \\ \hline
Accès mémoire & Direct dans instructions & Uniquement \texttt{LOAD}/\texttt{STORE} \\ \hline
Exécution & Microcode, cycles variables & Pipeline, 1 cycle/ins. (idéal) \\ \hline
Registres & Peu (8--16) & Nombreux (32+) \\ \hline
Exemples & x86, x86-64 (PC, serveurs) & ARM (smartphones, embarqué), RISC-V \\ \hline
\end{tabularx}
\end{center}

\courssub{Méthodes express}
\begin{methode}[Décrire le cycle d'exécution d'une instruction]
  \begin{enumerate}
    \item Identifier l'instruction pointée par le registre \texttt{PC} (Program Counter).
    \item \textbf{Chargement} : Copier l'instruction de la RAM (adresse \texttt{PC}) vers le registre \texttt{RI} (Registre d'Instruction) ; incrémenter \texttt{PC}.
    \item \textbf{Décodage} : L'UC analyse le code opération (opcode) dans \texttt{RI} et détermine l'opération et les opérandes (registres, adresse mémoire, valeur immédiate).
    \item \textbf{Exécution} : L'ALU réalise l'opération (calcul, test, saut, accès mémoire via \texttt{MAR}/\texttt{MDR}).
    \item \textbf{Écriture} : Stocker le résultat dans un registre ou en mémoire (via bus de données).
  \end{enumerate}
\end{methode}

\begin{methode}[Justifier le choix RISC ou CISC pour un usage donné]
  \begin{enumerate}
    \item Identifier la contrainte principale : \textbf{performance brute / faible conso} (embarqué, mobile) $\to$ RISC ; \textbf{compatibilité logicielle legacy / densité code} (PC, serveur classique) $\to$ CISC.
    \item Citer l'avantage clé : RISC = pipeline efficace, fréquence haute, faible chaleur ; CISC = instructions puissantes, moins d'instructions par programme, compatibilité historique.
    \item Conclure en reliant au contexte camerounais : ex. smartphone (ARM/RISC) vs poste cybercafé (x86/CISC).
  \end{enumerate}
\end{methode}
\end{bilan}

\begin{aretenir}[Checklist avant le BEPC]
  \begin{itemize}
    \item $\square$ Je sais définir l'architecture d'un ordinateur et nommer ses 3 grands blocs (CPU, Mémoire, E/S) reliés par les bus.
    \item $\square$ Je distingue l'Unité de Contrôle (UC) de l'ALU et je cite le rôle des registres (PC, RI, MAR, MDR, Accumulateur).
    \item $\square$ Je caractérise la RAM (volatile, rapide, adressable) vs mémoire de masse (non volatile, lente, grande capacité).
    \item $\square$ Je nomme les 3 bus (adresses, données, contrôle) et leur sens (uni/bidirectionnel).
    \item $\square$ Je récite les 4 étapes du cycle machine dans l'ordre : Chargement $\to$ Décodage $\to$ Exécution $\to$ Écriture.
    \item $\square$ Je sais expliquer le rôle du registre \texttt{PC} (compteur ordinal) et de l'incrémentation.
    \item $\square$ Je cite 3 différences majeures entre RISC et CISC (nb instructions, longueur, modes adressage, mémoire, registres).
    \item $\square$ Je donne un exemple de processeur CISC (x86) et un exemple RISC (ARM, RISC-V).
    \item $\square$ Je sais argumenter pourquoi les smartphones utilisent de l'ARM (RISC) : faible consommation, dissipation thermique, pipeline.
    \item $\square$ Je sais lire un organigramme simple du cycle d'exécution (losange décision, rectangles actions).
    \item $\square$ Je sais rédiger une réponse structurée : Définition $\to$ Schéma/Description $\to$ Comparatif $\to$ Exemple concret $\to$ Conclusion.
    \item $\square$ Je connais le vocabulaire anglais clé : \texttt{Fetch, Decode, Execute, Writeback, Pipeline, Load/Store}.
  \end{itemize}
\end{aretenir}

\findecours

\end{document}
